\documentclass[11pt]{amsart}
\usepackage[localtheorems]{raeez-math-template}

% ============================================================================
% Vol III --- The Gluing Chapter: unified integration of Parts I through X.
%
% Source: ten attack-heal-convergent Parts under notes/gluing_ch_part_*.tex,
% produced by a fifteen-agent elite-voice adversarial swarm (Wave-5 Phase-B,
% 2026-04-23). Each Part is a first-principles elementary derivation of one
% gluing mechanism; together they span the four equivariance strata and the
% sixteen-minus-one non-empty cells of the stratum combinatorics.
%
% Integration approach: this file is the amsart wrapper. The ten Part files
% are \input-ed as fragments (with their original standalone preambles
% stripped at \iffalse or commented out; package loads and theorem envs
% supplied here once). Every Part opens with \part{...}; its sections and
% subsections nest under that Part.
%
% Primary audit: notes/antipatterns_catalogue.md AP-CY293 through AP-CY344
% records every error-and-fix cycle from Waves 1-5 that fed into the current
% Part drafts.
% ============================================================================

\usepackage{stmaryrd}
\usepackage[capitalise,noabbrev]{cleveref}

% Theorem environments (supplied once at wrapper level).
\newtheorem{theorem}{Theorem}[section]
\newtheorem{proposition}[theorem]{Proposition}
\newtheorem{lemma}[theorem]{Lemma}
\newtheorem{corollary}[theorem]{Corollary}
\newtheorem{conjecture}[theorem]{Conjecture}
\theoremstyle{definition}
\newtheorem{definition}[theorem]{Definition}
\newtheorem{remark}[theorem]{Remark}
\newtheorem{example}[theorem]{Example}

% Claim-status discipline (Vol III manifesto).
\providecommand{\ClaimStatusTheorem}{\textbf{\textsf{[Theorem]}}}
\providecommand{\ClaimStatusProposition}{\textbf{\textsf{[Proposition]}}}
\providecommand{\ClaimStatusDefinition}{\textbf{\textsf{[Definition]}}}
\providecommand{\ClaimStatusProvedElsewhere}{\textbf{\textsf{[Proved elsewhere]}}}
\providecommand{\ClaimStatusProvedHere}{\textbf{\textsf{[Proved here]}}}
\providecommand{\ClaimStatusDerivation}{\textbf{\textsf{[Derivation]}}}
\providecommand{\ClaimStatusConjectured}{\textbf{\textsf{[Conjectural]}}}
\providecommand{\ClaimStatusOpen}{\textbf{\textsf{[Open]}}}
\providecommand{\ClaimStatusConventionDependent}{\textbf{\textsf{[Convention-dependent]}}}
\providecommand{\ClaimStatusPartialPrimaryLit}{\textbf{\textsf{[Primary-lit partial]}}}
\providecommand{\ClaimStatusPrincipleLevel}{\textbf{\textsf{[Principle-level]}}}

% Macro corpus (union of all Part providecommands).
\providecommand{\CoHA}{\mathrm{CoHA}}
\providecommand{\bY}{\mathbf{Y}}
\providecommand{\conBy}{Y_{\mathrm{con}}}
\providecommand{\Hom}{\mathrm{Hom}}
\providecommand{\Ext}{\mathrm{Ext}}
\providecommand{\End}{\mathrm{End}}
\providecommand{\Coh}{\mathrm{Coh}}
\providecommand{\Perf}{\mathrm{Perf}}
\providecommand{\Aut}{\mathrm{Aut}}
\providecommand{\Hilb}{\mathrm{Hilb}}
\providecommand{\Pic}{\mathrm{Pic}}
\providecommand{\Crit}{\mathrm{Crit}}
\providecommand{\HH}{\mathrm{HH}}
\providecommand{\Sp}{\mathrm{Sp}}
\providecommand{\Mp}{\mathrm{Mp}}
\providecommand{\RHom}{\mathrm{R}\mathrm{Hom}}
\providecommand{\Jac}{\mathrm{Jac}}
\providecommand{\Stab}{\mathrm{Stab}}
\providecommand{\Str}{\mathrm{Str}}
\providecommand{\Ran}{\mathrm{Ran}}
\providecommand{\Obs}{\mathrm{Obs}}
\providecommand{\Tot}{\mathrm{Tot}}
\providecommand{\Sym}{\mathrm{Sym}}
\providecommand{\Gm}{\mathbb{G}_m}
\providecommand{\cO}{\mathcal{O}}
\providecommand{\cH}{\mathcal{H}}
\providecommand{\cM}{\mathcal{M}}
\providecommand{\cN}{\mathcal{N}}
\providecommand{\cA}{\mathcal{A}}
\providecommand{\cC}{\mathcal{C}}
\providecommand{\cD}{\mathcal{D}}
\providecommand{\cE}{\mathcal{E}}
\providecommand{\cF}{\mathcal{F}}
\providecommand{\cG}{\mathcal{G}}
\providecommand{\cL}{\mathcal{L}}
\providecommand{\cS}{\mathcal{S}}
\providecommand{\cT}{\mathcal{T}}
\providecommand{\cU}{\mathcal{U}}
\providecommand{\cV}{\mathcal{V}}
\providecommand{\cX}{\mathcal{X}}
\providecommand{\cY}{\mathcal{Y}}
\providecommand{\cZ}{\mathcal{Z}}
\providecommand{\Db}{D^{b}}
\providecommand{\Z}{\mathbb{Z}}
\providecommand{\N}{\mathbb{N}}
\providecommand{\Q}{\mathbb{Q}}
\providecommand{\R}{\mathbb{R}}
\providecommand{\C}{\mathbb{C}}
\providecommand{\bZ}{\mathbb{Z}}
\providecommand{\bN}{\mathbb{N}}
\providecommand{\bQ}{\mathbb{Q}}
\providecommand{\bR}{\mathbb{R}}
\providecommand{\bC}{\mathbb{C}}
\providecommand{\bP}{\mathbb{P}}
\providecommand{\bL}{\mathbb{L}}
\providecommand{\fgl}{\mathfrak{gl}}
\providecommand{\fsl}{\mathfrak{sl}}
\providecommand{\fosp}{\mathfrak{osp}}
\providecommand{\fn}{\mathfrak{n}}
\providecommand{\fg}{\mathfrak{g}}
\providecommand{\fh}{\mathfrak{h}}
\providecommand{\fD}{\mathfrak{D}}
\providecommand{\fU}{\mathfrak{U}}
\providecommand{\fu}{\mathfrak{u}}
\providecommand{\sVect}{\mathrm{sVect}}
\providecommand{\kch}{\kappa_{\mathrm{ch}}}
\providecommand{\kcat}{\kappa_{\mathrm{cat}}}
\providecommand{\kBKM}{\kappa_{\mathrm{BKM}}}
\providecommand{\kfib}{\kappa_{\mathrm{fiber}}}
\providecommand{\kchBV}{\kappa_{\mathrm{ch,BV}}}
\providecommand{\SpCh}{\mathrm{Sp}^{\mathrm{ch}}}
\providecommand{\PhiFA}{\Phi^{\mathrm{FA}}}
\providecommand{\Eone}{E_{1}}
\providecommand{\Etwo}{E_{2}}
\providecommand{\Ethree}{E_{3}}
\providecommand{\Yplus}{Y^{+}}
\providecommand{\CY}{\mathrm{CY}}
\providecommand{\Cat}{\mathrm{Cat}}
\providecommand{\GlCoHA}{\mathrm{GlCoHA}}
\providecommand{\Mu}{\mathrm{M}}
\providecommand{\defeq}{\coloneqq}
\providecommand{\LL}{\mathbb{L}}
\providecommand{\B}{\mathrm{B}}
\providecommand{\bF}{\mathbb{F}}
\providecommand{\Ch}{\mathrm{Ch}}
\providecommand{\coCom}{\mathrm{coCom}}
\providecommand{\Conf}{\mathrm{Conf}}
\providecommand{\FA}{\mathrm{FA}}
\providecommand{\fl}{\mathfrak{l}}
\providecommand{\G}{\mathrm{G}}
\providecommand{\GRT}{\mathrm{GRT}}
\providecommand{\HolFA}{\mathrm{HolFA}}
\providecommand{\Lie}{\mathrm{Lie}}
% amsmath defines \Bar = \bar (math accent); Part V uses \Bar^{ch} as a
% bar-construction functor. Rebind locally for this chapter.
\let\BarAccent\Bar
\renewcommand{\Bar}{\mathrm{Bar}}

\title{CoHA gluing: ten routes, four strata, fifteen cells}
\author{The Gluing Chapter}
\date{2026-04-23}

\begin{document}

\maketitle

\begin{abstract}
\noindent
The critical cohomological Hall algebra $\cH(X) = \CoHA(\Db\Coh(X), \Omega_X)$
of a Calabi--Yau threefold is assembled from local data. This chapter classifies
how the assembly happens: the four equivariance strata ($T^d$-toric, reduced
$\Gm + \Aut$, orbifold inertia, lattice-polarised period), the sixteen-minus-one
non-empty cells of the stratum combinatorics, and the explicit local-to-global
cocycles that glue chart-wise constructions to a globally-defined Hall bialgebra.
Each Part resolves one construction: abstract 2-categorical framework (Part~I);
\v{C}ech gluing on toric atlases (Part~II); McKay and orbifold inertia
(Part~III); Van den Bergh derived Morita via tilting and NCCR (Part~IV);
factorisation algebras on Ran spaces and the two-stage $\Phi_3 =
\SpCh_{\Sigma_2, C} \circ \PhiFA_3$ (Part~V); Kontsevich--Soibelman wall-crossing
(Part~VI); lattice-polarised Borcherds automorphic gluing (Part~VII); the
$K3\times E$ master example (Part~VIII); obstructions on compact CY$_3$
(Part~IX); and the unifying sheaf of CoHAs on the $(2,1)$-stack of
equivariance strata (Part~X).
\end{abstract}

% ============================================================================
% Chapter-opening framing
% ============================================================================

A CY$_3$ category $\cC$ carries a Hall bialgebra $\cH(\cC)$. Locally -- on a
smooth formal chart, on a toric atlas, on an orbifold point, on a derived Morita
patch -- the Hall bialgebra has a tractable presentation: a shuffle algebra, a
Jacobi-algebra module category, a tilted dg-algebra, a vertex operator algebra.
The question addressed by this chapter is how the local presentations glue into
the global $\cH(\cC)$.

The answer is stratified. Four equivariance layers on $(X, \Omega_X)$ source
four distinct cocycle types:
\begin{itemize}[leftmargin=1.5em, itemsep=2pt]
\item \textbf{(i) Toric} -- chart transitions on a rational polyhedral fan;
cocycle in $H^1(\Sigma, \underline{\Pic}_{T^d})$.
\item \textbf{(ii) Reduced $\Gm$ and $\Aut$} -- connected automorphism action;
cocycle in $\Aut_s$-character data.
\item \textbf{(iii) Orbifold inertia} -- McKay / Galois data on $[X/G]$;
cocycle in $H^1_{\mathrm{Gal}}(G, \ldots)$.
\item \textbf{(iv) Lattice-polarised period} -- Borcherds automorphic data on
$(\Gamma, Z)$; cocycle in the Fourier coefficients of Gritsenko--Cl\'ery
weight-$k_\Psi = c_\phi(0,0)/2$ forms.
\end{itemize}
Every CY$_3$ carries a non-empty stratum-set $\Str(X) \subseteq \{\mathrm{i},
\mathrm{ii}, \mathrm{iii}, \mathrm{iv}\}$; the sixteen power-set cells partition
into one empty cell (the stratum-free compact generic case, Part~IX) and fifteen
non-empty cells whose cocycle theories compose into the global gluing datum.

The single object that coherently binds the four strata is the $(2,1)$-stack
$\mathrm{Stk}_{\Str}$ of equivariance strata. Part~X constructs a sheaf of
stable $\infty$-categorical CoHAs on $\mathrm{Stk}_{\Str}$, identifies its
stalks on each of the fifteen cells, and verifies the descent axiom for the
master example $K3 \times E$ (cell fifteen; all four strata meet).

% ============================================================================
% Ten Parts, \input as fragments.
% ============================================================================

\input{gluing_ch_part_1_abstract_framework}

% !TEX root = ../main.tex
% ============================================================================
% Gluing Chapter, Part II:
%   Chart-wise gluing on toric CY$_3$ (Stratum i)
%
% First-principles elementary-derivation draft.  Every step is pinned to an
% explicit chart computation: Jacobian of the transition, action on the
% $(3,0)$-form, induced action on polyvectors, Fourier--Mukai multiplier,
% Koszul sign.  This draft is standalone (CG voice, no "Wave N" language,
% no process narration, no retraction bookkeeping); it subsumes the
% conifold two-chart construction of
% \texttt{chapters/examples/toric\_cy3\_coha.tex} \S conifold-two-chart-cech,
% feeds the McKay/orbifold Part III, and states the local-to-global
% theorem that terminates Stratum~i.
% ============================================================================

\providecommand{\bY}{\mathbb{Y}}

\part{Chart-wise gluing on toric CY$_3$ (Stratum i)}

\medskip

Stratum~i is the toric regime: every top-dimensional cone of the fan
contributes a $\C^3$-chart carrying the Schiffmann--Vasserot positive
half $Y^+(\widehat{\fgl}_1)$, and the global CoHA is the
\v{C}ech--$E_3$-factorisation-algebra homotopy colimit of the chart
data along the face lattice of the fan.  The simplest non-trivial
instance is the resolved conifold
\(
  \bY = \mathrm{Tot}(\cO_{\bP^1}(-1) \oplus \cO_{\bP^1}(-1))
\)
with its two-chart atlas $\bY = U_+ \cup U_-$.  The multi-chart
generalisation indexes over a toric fan; the compact-$4$-cycle
obstruction separates zero-interior-point diagrams from diagrams
requiring McKay resolution; the local-to-global theorem closes the
toric programme at the zero-interior-point boundary.

% ============================================================================
\section{\v{C}ech gluing of two charts: the conifold paradigm}
\label{sec:gluing-p2-conifold}
% ============================================================================

The obstruction resolved by a two-chart atlas is the absence of a global
$\C^3$-parametrisation: the resolved conifold $\bY$ is a rank-$2$ vector
bundle over $\bP^1$, not an affine variety.  Any chart-wise
Schiffmann--Vasserot construction must name the transition datum that turns
two copies of $\C^3$ into $\bY$, then transport that datum through every
intermediate step --- the $(3,0)$-form, the polyvector fields, the Hall
product, the super structure.

\subsection{The two-chart atlas}

Let $z \in \bP^1$ be the standard affine coordinate on the base, with
$\tilde z = z^{-1}$ on the complementary chart.  On the total space
\(
  \bY = \mathrm{Tot}(\cO_{\bP^1}(-1) \oplus \cO_{\bP^1}(-1))
\)
introduce fibre coordinates $u_+, v_+$ in the trivialisation over the
$z$-chart and $u_-, v_-$ in the trivialisation over the $\tilde z$-chart.
The cover
\[
  U_+ \;=\; \{z \neq \infty\} \times \C^2 \;\cong\; \C^3,
  \qquad
  U_- \;=\; \{\tilde z \neq \infty\} \times \C^2 \;\cong\; \C^3,
\]
has overlap $U_+ \cap U_- = (\bP^1 \setminus \{0, \infty\}) \times \C^2
\cong \bC^\times \times \C^2$.  A section of $\cO_{\bP^1}(-1)$ is a function
$f(z)$ on $U_+$ that transforms to $\tilde z f(\tilde z^{-1}) = z^{-1} f(z)$
on $U_-$.  Writing $u_\pm, v_\pm$ as the fibre coordinates of the two $\cO(-1)$
summands one finds the transition
\begin{equation}\label{eq:conifold-transition-chart}
  \tilde z \;=\; z^{-1},
  \qquad
  u_- \;=\; z\, u_+,
  \qquad
  v_- \;=\; z\, v_+,
\end{equation}
as the natural datum forced by the $\cO(-1) \oplus \cO(-1)$ twist.  The
standard atlas of $\bP^1$ carries an orientation flip across
$z \leftrightarrow \tilde z$: the holomorphic $1$-form $dz = -\tilde z^{-2}\,
d\tilde z$ reverses sign through the change of chart, and this sign
propagates to the $(3,0)$-form computation of the next subsection.

\subsection{Jacobian of the transition on the $(3,0)$-form}

The Calabi--Yau $(3,0)$-form $\Omega_{\bY}$ is specified on $U_+$ by the
requirement that it trivialise the canonical bundle of $U_+ \cong \C^3$:
\[
  \Omega_{\bY}\big|_{U_+} \;=\; dz \wedge du_+ \wedge dv_+.
\]
To transport $\Omega_{\bY}$ to $U_-$ one computes the Jacobian of the
transition~\eqref{eq:conifold-transition-chart} step by step.  The exterior
derivatives are
\[
  dz \;=\; -\tilde z^{-2}\, d\tilde z,
  \qquad
  du_+ \;=\; d(z^{-1} u_-)
          \;=\; -z^{-2}\, u_-\, dz + z^{-1}\, du_-,
  \qquad
  dv_+ \;=\; -z^{-2}\, v_-\, dz + z^{-1}\, dv_-,
\]
so that
\[
  dz \wedge du_+ \;=\; dz \wedge z^{-1}\, du_-,
  \qquad
  dz \wedge du_+ \wedge dv_+ \;=\; dz \wedge z^{-1}\, du_- \wedge z^{-1}\, dv_-
  \;=\; z^{-2}\, dz \wedge du_- \wedge dv_-.
\]
Substituting $dz = -\tilde z^{-2}\, d\tilde z$ and $z^{-2} = \tilde z^{\,2}$:
\begin{equation}\label{eq:jacobian-calculation}
  \Omega_{\bY}\big|_{U_+}
  \;=\; dz \wedge du_+ \wedge dv_+
  \;=\; z^{-2}\, dz \wedge du_- \wedge dv_-
  \;=\; (\tilde z^{\,2}) \cdot (-\tilde z^{-2})\, d\tilde z \wedge du_- \wedge dv_-
  \;=\; -d\tilde z \wedge du_- \wedge dv_-.
\end{equation}
The fibre Jacobian $z^{-2} = \tilde z^{\,2}$ cancels the cotangent factor
$-\tilde z^{-2}$ from $dz$ exactly; the surviving sign is the
base-orientation flip of $\bP^1$.  This is the statement that $K_{\bY}$
is trivial as a line bundle on $\bY$ --- its transition cocycle on the
overlap $U_+ \cap U_-$ collapses to $\pm 1$ --- with the base-orientation
sign $-1$ read off the global trivialisation~\eqref{eq:jacobian-calculation}.
The fibre twist $z^{-2}$ (two factors of $\cO(-1)$, one from each
$\cO_{\bP^1}(-1)$ summand) compensates the cotangent factor $z^{-2}$ on
the base $1$-form $dz$, leaving the constant $-1$ as the whole transition
multiplier on $\Omega_\bY$.

\subsection{The chart-wise Schiffmann--Vasserot CoHA}

Each chart $U_\pm \cong \C^3$ carries the Jordan-triple quiver with
potential
\[
  Q_{\mathrm{Jor}} : \;\; \underbrace{\bullet}_{0}
  \;\;\text{with three loops}\;\; x_\pm, y_\pm, z_\pm,
  \qquad
  W_\pm \;=\; \mathrm{tr}\bigl(x_\pm[y_\pm, z_\pm]\bigr).
\]
The critical CoHA of $(Q_{\mathrm{Jor}}, W_\pm)$ is the positive half of the
affine Yangian of $\widehat{\fgl}_1$ (Schiffmann--Vasserot 2012
\texttt{arXiv:1202.2756}; Tsymbaliuk 2014 \texttt{arXiv:1404.5240}),
\[
  \CoHA\bigl(U_\pm, W_\pm\bigr) \;\cong\; Y^+\bigl(\widehat{\fgl}_1\bigr),
\]
with generators $e_r^{(\pm)}, f_r^{(\pm)}, h_r^{(\pm)}$ indexed by $r \in
\Z_{\geq 0}$ and structure functions given by the RSYZ shuffle kernel
(Rapcak--Soibelman--Yang--Zhao 2020 \texttt{arXiv:1810.10402} \S5,
equation (5.3), specialised to the single-vertex Jordan-triple case).
At weight $(1,1)$ the shuffle kernel evaluates to
\begin{equation}\label{eq:structure-constant-1-1}
  \varphi_{11}(u_1; u_2) \;=\; \frac{1}{(u_1 - u_2 - \varepsilon_1)(u_1 - u_2 - \varepsilon_2)}
  \;\;\xrightarrow{\;u_2 \to u_1\;}\;\;
  \frac{1}{\varepsilon_1 \varepsilon_2},
\end{equation}
the structure constant that every downstream computation must reproduce.

\subsection{Fourier--Mukai transition kernel}

On the overlap $U_+ \cap U_- = \bC^\times \times \C^2$ the derived category
$D^b(\mathrm{Coh}(U_+ \cap U_-))$ carries the autoequivalence implementing
the transition~\eqref{eq:conifold-transition-chart} at the level of sheaves.
Its integral kernel is
\begin{equation}\label{eq:FM-kernel-conifold}
  K_{+-} \;=\; \cO_\Delta \otimes \pi_1^*\,\cO_{\bP^1}(-1)
  \;\in\; D^b\bigl(\mathrm{Coh}\,(U_+ \cap U_-) \times (U_+ \cap U_-)\bigr),
\end{equation}
with $\Delta$ the diagonal and $\pi_1$ the projection onto the first factor.
The transform $T_{+-}(\cF) := \pi_{2,*}\bigl(\pi_1^*\cF \otimes K_{+-}\bigr)$
acts on polyvector fields
$\mathrm{PV}^\bullet(U_\pm) = \bigoplus_p \bigwedge^p T_{U_\pm}$ through the
chain rule applied to~\eqref{eq:conifold-transition-chart} combined with the
$\cO_{\bP^1}(-1)$-twist.

\paragraph{Chain rule on the cotangent basis.}
Differentiate~\eqref{eq:conifold-transition-chart} to obtain
\[
  du_+ \;=\; d(z^{-1} u_-) \;=\; -z^{-2}\, u_-\, dz + z^{-1}\, du_-,
  \qquad
  dv_+ \;=\; -z^{-2}\, v_-\, dz + z^{-1}\, dv_-.
\]
On the vertical (fibre-degree) sector --- the sector entering the
Schiffmann--Vasserot Hall product --- only the $z^{-1}\,du_-$ and
$z^{-1}\,dv_-$ terms survive (the $dz$-terms drop under pairing with the
Jordan-triple fibre polyvectors by~\eqref{eq:structure-constant-1-1}).
The cotangent multiplier is therefore $z^{-1}$ per fibre slot, giving
$z^{-(a+b)}$ on a fibre-form $du_+^a \wedge dv_+^b$.

\paragraph{Polyvector transport via $\Omega_\bY$-duality.}
The critical CoHA identifies polyvectors with differential forms via
Serre duality against the $(3,0)$-form $\Omega_\bY$:
\[
  \Lambda^p\, T_{\bY} \;\cong\; \Omega^{3-p}_\bY \otimes K_\bY^{-1}
  \;\cong\; \Omega^{3-p}_\bY
\]
using triviality of $K_\bY$ (equation~\eqref{eq:jacobian-calculation}).
Under this identification the fibre-polyvector $\partial_{u_+}^a
\partial_{v_+}^b$ corresponds (up to sign) to the fibre-form
$du_+^a \wedge dv_+^b$ contracted against the base $dz$-direction of
$\Omega_\bY$.  Its transport under $T_{+-}$ is the cotangent transformation
$du_+^a \wedge dv_+^b = z^{-(a+b)}\, du_-^a \wedge dv_-^b$.

\paragraph{$\cO(-1)$-kernel twist.}
The factor $\pi_1^*\cO_{\bP^1}(-1)$ in $K_{+-}$ contributes a further
multiplier $z^{-1}$ on the scalar sector: a section of
$\cO_{\bP^1}(-1)$ on $U_+$ written as $g(z)$ is expressed on $U_-$ as
$\tilde z\, g(\tilde z^{-1}) = z^{-1}\, g(z)$.

\paragraph{Net multiplier.}
Combining the cotangent re-section and the kernel twist,
\[
  \bigl[z^{-(a+b)}\bigr]_{\text{cotangent chain rule}}
    \;\cdot\;
  \bigl[z^{-1}\bigr]_{\cO(-1)\text{-kernel}}
  \;=\;
  z^{-(1 + a + b)}.
\]
The Fourier--Mukai transport of a Jordan-triple fibre polyvector therefore
reads
\begin{equation}\label{eq:FM-multiplier}
  T_{+-}^* :\;
  f(z, u_+, v_+)\, \partial_{u_+}^a\, \partial_{v_+}^b
  \;\longmapsto\;
  f(\tilde z^{-1}, \tilde z\, u_-, \tilde z\, v_-)\;
  z^{-(1 + a + b)}\;
  \partial_{u_-}^a\, \partial_{v_-}^b.
\end{equation}
The base-polyvector sector (powers of $\partial_z$) picks up the
base-orientation multiplier $(-\tilde z^{\,2})$ per $\partial_z$-slot; the
Jordan-triple selection rule~\eqref{eq:structure-constant-1-1} restricts
the Hall product to the fibre sector $\alpha = 0$ on which the base
orientation enters only through the global sign
of~\eqref{eq:jacobian-calculation}.

\subsection{Homotopy colimit and super-Yangian output}

The pair $(U_+, U_-)$ with overlap $U_+ \cap U_-$ forms a two-object diagram
in the category of $\C^3$-charts.  The chart-wise factorisation algebras
$\cF_\pm = \Phi^{FA}_3(\Gamma_\pm)$ (Ginzburg dgas of the Jordan-triple
quiver with potential, \S\ref{sec:gluing-p2-conifold}; Kontsevich--Tamarkin
$E_3$-formality + Costello--Li holomorphic locality) glue by the homotopy
colimit
\begin{equation}\label{eq:homotopy-colimit}
  \cF_{\bY} \;=\; \mathrm{hocolim}\Bigl(
    \cF_+\big|_{U_+ \cap U_-}
    \;\xrightarrow{\;T_{+-}\;}\;
    \cF_-\big|_{U_+ \cap U_-}
  \Bigr),
\end{equation}
an $E_3$-factorisation algebra on $\bY$ whose cohomology is
$H^\bullet(\cF_{\bY}) = \mathrm{PV}^\bullet(\bY)$.  Under the CoHA assignment
the homotopy colimit produces a new associative algebra whose generators
split into two families indexed by the two charts and whose structure
constants are determined by the multiplier~\eqref{eq:FM-multiplier}.

\paragraph{The super structure as a Koszul sign.}
The base-orientation flip in $\Omega_{\bY}|_{U_-} = -d\tilde z \wedge du_-
\wedge dv_-$ (equation~\eqref{eq:jacobian-calculation}) contributes a
Koszul sign $\sigma = -1$ on cross-chart Hall products: the orientation of
the measure against which the critical CoHA integrates polyvectors
reverses between charts, so bilinear pairings $\langle \cdot, \cdot
\rangle_\pm$ that glue across the overlap pick up a factor of $-1$.  In
the chart-wise shuffle presentation this is the parity assignment of
the Schiffmann--Vasserot generators across the two charts, producing
super-Yangian-type anticommutators on the combinations
\[
  E^{(0)} \;=\; e_0^{(+)} + e_0^{(-)},
  \qquad
  E^{(1)} \;=\; e_0^{(+)} - e_0^{(-)},
\]
where the even combination $E^{(0)}$ generates a bosonic subalgebra and the
odd combination $E^{(1)}$ generates a fermionic sector with
\[
  \bigl\{E^{(1)}, E^{(0)}\bigr\}_+ \;=\; K^{(0)},
\]
the central element extracted from the residue at $z_1 = z_2$ of the
chart-wise Jordan shuffle difference~\eqref{eq:structure-constant-1-1}.
The two-chart assembly therefore yields
\begin{equation}\label{eq:conifold-master}
  \CoHA(\bY) \;\cong\; Y^+\bigl(\widehat{\fgl}(1|1)\bigr)^{\mathrm{con}},
\end{equation}
the positive half of the conifold affine super-Yangian.  The super
structure on $\widehat{\fgl}(1|1)$ arises from the Koszul sign
$\sigma = -1$ of the base-orientation flip~\eqref{eq:jacobian-calculation}
combined with the isotropy $(\alpha_0, \alpha_0) = (\alpha_1, \alpha_1) = 0$
(Kac 1977 Thm.~2.4 on odd simple roots).

\paragraph{Three-route verification of the $(1,1)$ structure constant.}
The $(1,1)$ structure constant of $\CoHA(\bY)$ equals $(\varepsilon_1
\varepsilon_2)^{-1}$ through three independent routes:
\begin{enumerate}[label=\emph{(\roman*)}, leftmargin=2em]
\item \emph{\v{C}ech two-chart gluing}: the limit $u_2 \to u_1$ of the
  chart-wise Jordan shuffle~\eqref{eq:structure-constant-1-1}.
\item \emph{Van den Bergh tilting}: the integral of the $\varphi_W$-twisted
  Euler class over $\mathrm{Ext}^1(S_0, S_1) = \C^2$, the rank-$2$
  Ext-space between the simple modules of the Klebanov--Witten quiver.
\item \emph{Klebanov--Witten super-shuffle}: the Jordan kernel residue at
  $z_1 = z_2$ of the direct super-shuffle algebra.
\end{enumerate}
Agreement across the three routes pins $(\varepsilon_1 \varepsilon_2)^{-1}$
as a lane-independent invariant, confirming~\eqref{eq:conifold-master} at
the coefficient level.

% ============================================================================
\section{Multi-chart toric fan generalisation}
\label{sec:gluing-p2-multi-chart}
% ============================================================================

The two-chart assembly generalises to any local toric CY$_3$.  The obstruction
resolved is the absence of a global affine chart: a toric variety is covered
by affine opens $U_\sigma$ indexed by the maximal cones $\sigma$ of the fan,
and the global CoHA is assembled by running the \v{C}ech machinery along the
face lattice of the fan.

\subsection{Toric fan and chart decomposition}

Let $X = X_\Sigma$ be a smooth toric CY$_3$ with fan $\Sigma$ in the lattice
$N \cong \Z^3$.  The Calabi--Yau condition is that the primitive generators
of all rays of $\Sigma$ lie in a common affine hyperplane, equivalently
that there exists a lattice vector $m \in M = N^\vee$ such that
$\langle m, v_\rho \rangle = 1$ for every ray generator $v_\rho$.  A
natural choice places all ray generators at height $1$ along the
last coordinate; the toric diagram is then the $2$-dimensional polytope
obtained by projecting the rays to the plane $\{x_3 = 1\}$.

\paragraph{Chart structure at top cones.}
Each maximal (top-dimensional) cone $\sigma \in \Sigma(3)$ of a smooth
toric CY$_3$ is generated by exactly three rays
$\rho_1, \rho_2, \rho_3$ forming a $\Z$-basis of $N$, and determines an
affine chart
\[
  U_\sigma \;=\; \mathrm{Spec}\,\C[\sigma^\vee \cap M] \;\cong\; \C^3.
\]
Coordinates $(x_1^\sigma, x_2^\sigma, x_3^\sigma)$ on $U_\sigma$ are the
characters of the dual basis of $\sigma^\vee$.  The Jordan-triple
potential on each chart is
\[
  W_\sigma \;=\; \mathrm{tr}\bigl(x_1^\sigma [x_2^\sigma, x_3^\sigma]\bigr),
\]
and the chart-wise CoHA is $\CoHA(U_\sigma, W_\sigma) \cong
Y^+(\widehat{\fgl}_1)$ by the Schiffmann--Vasserot theorem applied to each
$\C^3$ with the standard Jordan-triple.

\paragraph{Equivariant weights and the CY condition.}
The three-torus $T^3 = (\bC^\times)^3$ acts on each $U_\sigma \cong \C^3$
with equivariant weights $(\varepsilon_1^\sigma, \varepsilon_2^\sigma,
\varepsilon_3^\sigma)$.  The Calabi--Yau constraint on $W_\sigma$ to be a
tri-vector of vanishing weight is
\[
  \varepsilon_1^\sigma + \varepsilon_2^\sigma + \varepsilon_3^\sigma \;=\; 0,
\]
which specialises the chart-wise weights to those of a local CY$_3$.  Along
a compact direction (a ray $\rho$ in $\Sigma(1)$ whose dual edge in the
toric diagram is an internal edge) the equivariant parameter vanishes:
the two fixed points of the compact $\bP^1$-base on that edge would carry
opposite weights, obstructing global regularity of sections.  This is the
mechanism that set $\varepsilon_3 = 0$ for the conifold base in
\S\ref{sec:gluing-p2-conifold}.

\subsection{Gluing cocycle on the face lattice}

Two top cones $\sigma, \sigma'$ of $\Sigma(3)$ sharing a common face
$\tau = \sigma \cap \sigma' \in \Sigma(2)$ of codimension one define an
overlap $U_\sigma \cap U_{\sigma'} = U_\tau$.  The face $\tau$ is a
two-dimensional cone, generated by two rays; the corresponding overlap is
\[
  U_\tau \;\cong\; \bC^\times \times \C^2,
\]
with the $\bC^\times$-factor encoding the direction transverse to $\tau$
within the full three-dimensional lattice.  The transition between the chart coordinates
$(x_1^\sigma, x_2^\sigma, x_3^\sigma)$ and $(x_1^{\sigma'}, x_2^{\sigma'},
x_3^{\sigma'})$ follows the lattice change of basis between $\sigma^\vee$
and $(\sigma')^\vee$: when $\sigma'$ differs from $\sigma$ by a toric flop
sending ray $\rho_1$ to $\rho_1'$ across the shared face $\tau$ (generated
by $\rho_2, \rho_3$), the conifold-type transition is
\begin{equation}\label{eq:multi-chart-transition}
  x_1^{\sigma'} \;=\; (x_1^\sigma)^{-1},
  \qquad
  x_i^{\sigma'} \;=\; x_1^\sigma \cdot x_i^\sigma
  \;\; \text{for } i = 2, 3,
\end{equation}
in direct analogy with the conifold
transition~\eqref{eq:conifold-transition-chart}.  More generally, if the
toric flop twists the fibre $\rho_1$-direction with multiplicity $(k_2, k_3)$
along $\rho_2, \rho_3$ (equivalently, the fibre bundle over the shared
$\bP^1_\tau$ is $\cO(-k_2) \oplus \cO(-k_3)$), the transition reads
$x_i^{\sigma'} = (x_1^\sigma)^{k_i}\, x_i^\sigma$ on the $i$-th fibre
coordinate.  The Fourier--Mukai kernel on the overlap is
\[
  K_{\sigma \sigma'} \;=\; \cO_\Delta \otimes \pi_1^* \cO_{\bP^1_{\tau}}(-1),
\]
the $\cO(-1)$-twist sourced by the compact $\bP^1_\tau$-base transverse
to $\tau$; the polyvector multiplier
$(x_1^\sigma)^{-(1 + k_2\, a + k_3\, b)}$ generalises~\eqref{eq:FM-multiplier}
with $a, b$ the fibre polyvector indices.  The Jacobian
argument of \S\ref{sec:gluing-p2-conifold} carries over with $z$ replaced
by $x_1^\sigma$ on each codimension-one face, and with the fibre twist
exponents $(k_2, k_3)$ replacing the conifold's $(1, 1)$; the
Calabi--Yau condition $k_2 + k_3 = 2$ (sum of $\cO(-k_i)$-degrees equals
$-\deg K_{\bP^1_\tau} = -2$) enforces triviality of
$K_{X}|_{\bP^1_\tau}$.

\paragraph{The face lattice cocycle.}
The collection $\{K_{\sigma \sigma'}\}_{\sigma, \sigma' \text{ adjacent}}$
forms a cocycle on the face lattice of $\Sigma$: at any triple overlap
$U_\sigma \cap U_{\sigma'} \cap U_{\sigma''}$ (corresponding to a
codimension-two face $\tau = \sigma \cap \sigma' \cap \sigma''$), the
compatibility
\[
  K_{\sigma \sigma''} \;=\; K_{\sigma' \sigma''} \circ K_{\sigma \sigma'}
\]
holds up to canonical isomorphism, because each kernel is the pullback of
$\cO_{\bP^1}(-1)$ along a $\bP^1$-fibre direction and the fibre directions
compose additively in the toric fan.  The cocycle class is the derivative
of the data on the fan's $1$-skeleton and is classified by the sheaf
cohomology group $H^1(\Sigma, \underline{\mathrm{Pic}}(\bP^1))$ on the fan
viewed as a CW complex.

\subsection{Galakhov--Li--Yamazaki bond factors}

Galakhov--Li--Yamazaki (2021 \texttt{arXiv:2106.01230}; Li--Yamazaki 2020
\texttt{arXiv:2003.08909} eq~(8.125)) present the multi-chart gluing data
via \emph{bond factors} $\varphi^{a \Rightarrow b}(u)$ indexed by ordered
pairs $(a, b)$ of vertices in the toric quiver.  The convention is:
\begin{equation}\label{eq:gly-bond-factor}
  \varphi^{a \Rightarrow b}(u)
  \;=\;
  \frac{\prod_{I \in \mathrm{Arr}(a \to b)} (u + h^{a \to b}_I)}{%
        \prod_{J \in \mathrm{Rel}(a \to b)} (u + h^{a \to b}_J)},
\end{equation}
with $h^{a\to b}_I$ the equivariant weight of the $I$-th arrow from
vertex $a$ to vertex $b$, and $h^{a\to b}_J$ the equivariant weight of
the $J$-th Jacobi relation ($\partial W / \partial (\text{arrow})$) tying
arrows from $a$ into and out of a trip through $b$.  The \emph{numerator
is arrow weights; the denominator is Jacobi-relation weights}; the
$(u + h)$-convention signs track the Drinfeld shifted-$R$-matrix
normalisation (Neguț 2015 \texttt{arXiv:1505.01528} eq~(1.6);
Tsymbaliuk 2014 \texttt{arXiv:1404.5240}).

\paragraph{Conifold bond factor from first principles.}
The conifold quiver has two vertices $\{0, 1\}$, two arrows
$a_1, a_2 : 0 \to 1$ with equivariant weights $h_1, h_2$, two arrows
$b_1, b_2 : 1 \to 0$ with weights $-h_1, -h_2$, and the Klebanov--Witten
superpotential $W_{\mathrm{con}} = a_1 b_1 a_2 b_2 - a_1 b_2 a_2 b_1$.
The Jacobi relations $\partial W / \partial b_i = 0$ and $\partial W /
\partial a_j = 0$ pair arrows from $0 \to 1$ into and out of cycles
through vertex $1$; for the bond $0 \Rightarrow 1$, the two Jacobi
relations carry weights $0$ and $h_1 + h_2$ (derived by summing the
weights of the arrows meeting each relation-index).  Substituting into
the GLY formula~\eqref{eq:gly-bond-factor}:
\begin{equation}\label{eq:conifold-bond-factor}
  \varphi^{0 \Rightarrow 1}_{\mathrm{con}}(u)
  \;=\;
  \frac{(u + h_1)(u + h_2)}{u \cdot (u + h_1 + h_2)},
\end{equation}
matching Neguț's conifold shuffle kernel (Neguț 2015
\texttt{arXiv:1505.01528} eq~(1.6)) and Li--Yamazaki's
BPS-quiver-algebra master formula (\texttt{arXiv:2003.08909} eq~(8.125)).
At the Calabi--Yau torus-equivariant slice $\varepsilon_1 + \varepsilon_2
+ \varepsilon_3 = 0$, the conifold weights are $h_1 = \varepsilon_1$,
$h_2 = \varepsilon_2$, and $h_1 + h_2 = \varepsilon_1 + \varepsilon_2 =
-\varepsilon_3$; substituting into~\eqref{eq:conifold-bond-factor}:
\[
  \varphi^{0 \Rightarrow 1}_{\mathrm{con}}(u)
  \;=\;
  \frac{(u + \varepsilon_1)(u + \varepsilon_2)}{u\,(u + \varepsilon_1 + \varepsilon_2)}
  \;=\;
  \frac{(u + \varepsilon_1)(u + \varepsilon_2)}{u\,(u - \varepsilon_3)}.
\]
Both presentations are equivalent through the Calabi--Yau constraint
$\varepsilon_3 = -\varepsilon_1 - \varepsilon_2$.

\paragraph{Arrow-weight / relation-weight trace on the toric diagram.}
The arrow count $|\mathrm{Arr}(a \to b)|$ matches the number of toric
diagram edges joining the dual vertices of $a$ and $b$; the
Jacobi-relation count $|\mathrm{Rel}(a \to b)|$ matches the number of
toric triangles sharing the edge.  For the conifold ($2$ arrows, $2$
relations per direction) the $2 + 2$ tally reflects the two-edge,
two-triangle diagram.  For local $\bP^2$ the bond factors are labelled
by the three $\Z_3$-characters of the McKay quiver (Beilinson nine-arrow
structure, \S\ref{sec:gluing-p2-compact-cycle}); for the banana
threefold, by the three edges of the banana diagram; for the suspended
pinch point, by the pinch-point bond structure.

\subsection{Examples}

\paragraph{$\C^3/\Z_n$ (linear toric).}
The fan has $n$ maximal cones arrayed along the toric diagram edge, each
$\C^3$; the face lattice is a linear chain with $n - 1$ codimension-one
faces; the gluing cocycle is a chain of $n - 1$ FM kernels, each a
$\cO(-1)$-twist.  The multi-chart \v{C}ech assembly produces the
$A_{n-1}$-type McKay CoHA, matching the Beasley--Plesser antisymmetric
quiver on $\C^3/\Z_n$.

\paragraph{Banana threefold.}
Three ray clusters meeting at a common vertex produce a fan with three
maximal cones and three codimension-one faces.  The \v{C}ech assembly
requires three FM kernels on a triangular face lattice; the resulting
CoHA is a triangular gluing of three Schiffmann--Vasserot
$Y^+(\widehat{\fgl}_1)$ copies, governed by the three-arrow banana quiver
with cubic potential.

\paragraph{Suspended pinch point.}
Four maximal cones sharing a common pinch ray produce a fan with four
charts and six codimension-one faces.  The \v{C}ech cocycle is defined on
a square face lattice with one internal pinch; the CoHA assembly
reproduces the SPP quiver of Aganagic--Karch--L\"ust--Miemiec (1999).

% ============================================================================
\section{The compact $4$-cycle obstruction}
\label{sec:gluing-p2-compact-cycle}
% ============================================================================

The obstruction resolved here is subtle: a local toric CY$_3$ with a
compact divisor --- a maximal cone whose dual vertex lies in the
\emph{interior} of the toric diagram --- admits a chart-wise $\C^3$ cover,
but the direct \v{C}ech--$E_3$-factorisation assembly of
\S\ref{sec:gluing-p2-multi-chart} misses the cohomology class of the
compact divisor.

\subsection{Interior lattice points and compact divisors}

Let $X_\Sigma$ be a local toric CY$_3$ with toric diagram $\Delta_\Sigma$ a
lattice polygon in the plane $\{x_3 = 1\}$.  The duality between toric
fans and toric diagrams identifies:
\[
  \text{rays of } \Sigma \;\longleftrightarrow\; \text{vertices of } \Delta_\Sigma,
  \qquad
  \text{maximal cones} \;\longleftrightarrow\; \text{internal triangulation triangles},
\]
and one distinguishes:
\begin{itemize}
\item \emph{boundary rays}: rays whose dual vertex lies on the boundary of
  $\Delta_\Sigma$; these correspond to non-compact divisors in $X$;
\item \emph{interior rays}: rays whose dual vertex lies in the interior
  of $\Delta_\Sigma$; these correspond to compact divisors, i.e.\ compact
  $4$-cycles in $X$.
\end{itemize}
The resolved conifold has no interior lattice point: its toric diagram is a
unit square with vertices $(0,0), (1,0), (0,1), (1,1)$ and no interior
lattice point, consistent with $\chi(\bY) = 2$ and no compact $4$-cycle.
Local $\bP^2 = \mathrm{Tot}(K_{\bP^2})$ has exactly one interior lattice
point: its toric diagram is the triangle with vertices $(1,0), (0,1),
(-1,-1)$ (the three non-compact divisor rays at height $1$) and the
interior lattice point $(0,0)$ at the origin; the ray $(0,0,1)$ attached
to this interior point generates the zero-section $\bP^2 \subset K_{\bP^2}$,
a compact $4$-cycle.

\subsection{Why chart-wise gluing misses the compact divisor}

The chart-wise \v{C}ech assembly of \S\ref{sec:gluing-p2-multi-chart} has
three ingredients on each chart: the Schiffmann--Vasserot
$Y^+(\widehat{\fgl}_1)$, the equivariant torus action, and the FM kernel
on codimension-one overlaps.  None of these detect the cohomology of a
compact divisor:

\begin{itemize}
\item The chart-wise polyvector complex $\mathrm{PV}^\bullet(U_\sigma)$ is
  concentrated in $(0, *)$-degree (holomorphic polyvector fields on
  $\C^3$), and has no contribution from the class of a compact surface.
\item The FM kernels $K_{\sigma \sigma'}$ transport polyvectors across
  codimension-one faces but do not generate new cohomology classes:
  they are sheaves of derived Morita data, not cycle classes.
\item The $T^3$-equivariant localisation concentrates all integrals at
  the fixed points of the torus action, i.e.\ the maximal cones; the
  contribution of a compact divisor is the \emph{sum over the fixed
  points contained in the divisor}, which must be computed by a separate
  procedure.
\end{itemize}
The Neguț wheel-condition corrections --- the shuffle-algebra
relations forced by the presence of an interior lattice point --- are
missed by the direct \v{C}ech assembly.  For local $\bP^2$ with single
interior point at the origin of the triangulation, the correct shuffle
algebra is the $\Z_3$-McKay variant of $Y^+(\widehat{\fgl}_1)$ whose
kernel carries an additional wheel-condition multiplier tracking the
$\Z_3$-orbit of the interior lattice point (Neguț 2015
\texttt{arXiv:1512.06473}; Rapcak--Soibelman--Yang--Zhao 2020
\texttt{arXiv:2001.10549}).  The direct multi-chart \v{C}ech assembly of
\S\ref{sec:gluing-p2-multi-chart} picks up only the boundary contributions
to $\varphi_{a \Rightarrow b}$ and misses this wheel-condition
multiplier.

\subsection{Resolution via McKay/orbifold route}

The standard resolution of the obstruction is to pass to the orbifold
quotient: local $\bP^2 = [\C^3/\Z_3]^{\mathrm{res}}$ (the crepant
resolution of the McKay orbifold), and apply the Bridgeland--King--Reid
equivalence
\[
  D^b\bigl(\mathrm{Coh}(K_{\bP^2})\bigr) \;\simeq\; D^b\bigl(\mathrm{Coh}^{\Z_3}(\C^3)\bigr)
\]
(Bridgeland--King--Reid 2001 \texttt{arXiv:math/9908027} Thm.~1.2).  The
right-hand category is a single $\C^3$-chart with $\Z_3$-equivariance;
the CoHA of this orbifold chart carries the Neguț wheel-condition
multipliers automatically, since the $\Z_3$-action on $\C^3$ produces
exactly the interior-point shifts at each orbit element.  The compact
$4$-cycle cohomology is recovered as the $\Z_3$-fixed part of the
$\C^3/\Z_3$ inertia stack.

Alternatively, one may keep the multi-chart cover and enrich each
chart-wise shuffle kernel with a higher-rank Neguț correction: replace
$\varphi_{a \Rightarrow b}(u)$ by the shuffle kernel
$\varphi_{a \Rightarrow b}^{\mathrm{rank-}r}(u)$ of the rank-$r$ shuffle
algebra, with $r$ the number of interior lattice points of the toric
diagram.  The two resolutions --- McKay/orbifold route (Part III) and
higher-rank shuffle correction --- are Morita-equivalent under the
Rapcak--Soibelman--Yang--Zhao toric theorem (2020
\texttt{arXiv:2001.10549} Thm.~B), but differ in their chart-wise
presentations.

\begin{proposition}[Compact-$4$-cycle obstruction]
\label{prop:compact-cycle-obstruction}
Let $X_\Sigma$ be a local toric CY$_3$ with toric diagram $\Delta_\Sigma$
carrying $r \geq 1$ interior lattice points.  The direct multi-chart
\v{C}ech assembly of \S\ref{sec:gluing-p2-multi-chart} reproduces
$\CoHA_{\mathrm{DM}}(X_\Sigma)$ only after replacing the chart-wise
single-vertex Jordan-triple data by either:
\begin{enumerate}[label=\emph{(\alph*)}, leftmargin=2em]
\item the $G$-equivariant Schiffmann--Vasserot data on $[\C^3 / G]$,
  where $G \subset \mathrm{SU}(3)$ is the McKay group whose crepant
  resolution $\widetilde{\C^3/G}$ recovers $X_\Sigma$ (McKay/orbifold
  route; applies when $\Delta_\Sigma$ is a McKay triangle, including
  $X_\Sigma = K_{\bP^2}$ with $G = \Z_3$);
\item the multi-vertex Neguț shuffle algebra $Y^+(Q_X, W_X)$ whose
  vertex set matches the full toric quiver of $X_\Sigma$ (interior and
  boundary lattice points together) and whose bond factors carry the
  wheel-condition multipliers tracking the interior-point positions
  (higher-rank shuffle route).
\end{enumerate}
Direct multi-chart assembly with only the chart-wise $Y^+(\widehat{\fgl}_1)$
data misses the interior-point contributions to the Hall product.
\end{proposition}

\begin{proof}[Sketch]
The chart-wise single-vertex polyvectors generate only the
boundary-vertex portion of the full toric quiver shuffle algebra; the
interior-vertex portion is generated by compositions extending into the
$G$-McKay inertia (Bridgeland--King--Reid 2001 \texttt{arXiv:math/9908027}
Thm.~1.2).  The Rapcak--Soibelman--Yang--Zhao toric theorem (2020
\texttt{arXiv:2001.10549} Thm.~B) identifies the full quiver shuffle
data with $\CoHA_{\mathrm{DM}}(X_\Sigma)$ when the interior vertices are
included; the McKay/orbifold route recovers the interior through the
$[\C^3/G]$-orbifold chart and the BKR equivalence, and the higher-rank
shuffle route recovers it by enlarging the single-vertex chart shuffle
to the full $(Q_X, W_X)$-shuffle with interior-vertex weights inserted.
\end{proof}

% ============================================================================
\section{Local-to-global theorem for zero-interior-lattice-point diagrams}
\label{sec:gluing-p2-local-to-global}
% ============================================================================

The combination of \S\S\ref{sec:gluing-p2-conifold}--\ref{sec:gluing-p2-compact-cycle}
yields the closure theorem of Stratum~i: for a local toric CY$_3$ whose
toric diagram has no interior lattice points, the chart-wise \v{C}ech
assembly is exact and reproduces the Davison--Meinhardt critical CoHA.

\begin{theorem}[Local-to-global for zero-interior-point toric CY$_3$]
\label{thm:local-to-global-zero-interior}
Let $X_\Sigma$ be a smooth local toric Calabi--Yau threefold whose toric
diagram $\Delta_\Sigma$ has zero interior lattice points.  Then:
\begin{enumerate}[label=\emph{(\arabic*)}, leftmargin=2em]
\item The multi-chart \v{C}ech factorisation-algebra assembly
\[
  \cF_X \;=\; \mathrm{hocolim}_{\sigma \in \Sigma(3)}\,
    \cF_\sigma
    \;\big|\;
    \bigl\{T_{\sigma \sigma'} : \cF_\sigma \leftrightarrow \cF_{\sigma'}\bigr\},
\]
with $\cF_\sigma = \Phi^{FA}_3(\Gamma_\sigma)$ the Ginzburg dga of the
chart Jordan-triple $(Q_{\mathrm{Jor}}, W_\sigma)$ and $T_{\sigma \sigma'}$
the Fourier--Mukai transition with kernel $K_{\sigma \sigma'} = \cO_\Delta
\otimes \pi_1^*\cO_{\bP^1_\tau}(-1)$, is a well-defined $E_3$-factorisation
algebra on $X$.
\item The associated critical CoHA
\[
  \CoHA_{\mathrm{\check Cech}}(X)
  \;=\;
  \mathrm{hocolim}_{\sigma \in \Sigma(3)}\,
    Y^+(\widehat{\fgl}_1)_\sigma
    \;\big|\;
    \bigl\{K_{\sigma \sigma'}\bigr\}
\]
is isomorphic to the Davison--Meinhardt critical cohomological Hall
algebra
\(
  \CoHA_{\mathrm{DM}}(X) = \CoHA(Q_X, W_X)
\)
of the toric quiver with potential $(Q_X, W_X)$ attached to $\Sigma$.
\item The resulting algebra is the Drinfeld positive half of an affine
super-Yangian of the root system determined by $(Q_X, W_X)$, with rank
equal to the number of vertices of $Q_X$ and super structure controlled
by the parity of the base-orientation sign on each codimension-one face
of the fan.
\end{enumerate}
\end{theorem}

\begin{proof}[Sketch]
The \v{C}ech homotopy colimit is computed face by face: each
codimension-one face contributes an FM kernel~\eqref{eq:FM-kernel-conifold}
scaled by the multiplier of equation~\eqref{eq:FM-multiplier}; each
codimension-two face contributes a cocycle compatibility checked by
the Jacobian calculation of equation~\eqref{eq:jacobian-calculation}.
Zero interior lattice points means no codimension-three (interior)
corrections from the Neguț wheel condition, so each face contribution is
the direct Schiffmann--Vasserot shuffle kernel without higher-rank
adjustment.  The identification with $\CoHA_{\mathrm{DM}}(X)$ is the
Rapcak--Soibelman--Yang--Zhao toric theorem (2020
\texttt{arXiv:2001.10549} Thm.~A), which matches the chart-wise
Schiffmann--Vasserot data to the toric quiver BPS algebra.  The
super-structure on the output is the base-orientation cocycle of the
fan as computed in equation~\eqref{eq:jacobian-calculation}: for fans
with no internal edges (no compact $\bP^1$-bases), the cocycle is
trivial and the output is an even affine Yangian; for fans with one
internal edge (e.g.\ the conifold), the cocycle produces the super
structure $\widehat{\fgl}(1|1)$ of equation~\eqref{eq:conifold-master};
for fans with several internal edges, the cocycle produces the product
super structure of the corresponding super-Yangian.
\end{proof}

The class of local toric CY$_3$ covered by
Theorem~\ref{thm:local-to-global-zero-interior} includes:
\begin{itemize}
\item $\C^3$ (one chart, no overlaps, $\CoHA = Y^+(\widehat{\fgl}_1)$);
\item the resolved conifold (two charts, one overlap,
  $\CoHA = Y^+(\widehat{\fgl}(1|1))^{\mathrm{con}}$);
\item the banana threefold (three charts, three overlaps, three-loop fan);
\item the suspended pinch point (four charts, six overlaps);
\item generalised conifolds $X_{n, m}$ (two charts, one overlap with
  $\cO(-n) \oplus \cO(-m)$-twist).
\end{itemize}

The class \emph{not} covered by the theorem comprises exactly those local
toric CY$_3$ with one or more interior lattice points --- local $\bP^2$,
local $\bP^1 \times \bP^1$, local Hirzebruch $\bF_n$ for $n \geq 1$, local
del Pezzo $dP_k$ for $k \geq 1$.  These all carry compact $4$-cycles and
must pass through Proposition~\ref{prop:compact-cycle-obstruction}'s
McKay/orbifold or higher-rank shuffle resolutions.  Part III takes up
this resolution, replacing the direct chart-wise \v{C}ech assembly by the
BKR equivalence and the inertia-stack decomposition of the orbifold
route; the resulting picture extends to the full toric Stratum~i up to
$\Z_n$-orbifolds and recovers, in particular, the nine-arrow Beilinson
CoHA of local $\bP^2$ as the \v{C}ech assembly of the $[\C^3/\Z_3]$-chart
across its $\Z_3$-skew-group algebra.

\bigskip

The chart-wise assembly of Stratum~i is now closed.  Three invariants have
been extracted: the face-lattice cocycle with values in
$H^1(\Sigma, \underline{\mathrm{Pic}})$; the bond-factor presentation of
Galakhov--Li--Yamazaki on the toric diagram edges; the zero-interior-point
criterion that separates the direct-\v{C}ech-assembly diagrams from those
requiring McKay resolution.  Part III addresses the McKay route for
diagrams with one or more interior points, and Part IV takes up derived
Morita gluing for the Van den Bergh tilting presentation common to
both.  The master test case $K3 \times E$, which combines Stratum i with
Strata ii--iv, is resolved in Part VIII via the full four-stratum
decomposition.


% !TEX root = ../main.tex
%% Part III of the Gluing Chapter: McKay / orbifold gluing (Stratum iii)
%% Elementary derivation, Chriss-Ginzburg voice, 3000-5000 words
%% Date: 2026-04-23

\part{McKay and orbifold gluing}
\label{part:gluing-mckay-orbifold}

An orbifold $[X/G]$ with $G \subset \mathrm{SU}(d)$ finite is a CY$_d$ category whose atlas is not a union of \v{C}ech charts but a single chart bearing the $G$-action. The gluing cocycle is not valued in $H^1$ of the overlap nerve; it is valued in the representation category $\mathrm{Rep}(G)$ attached to each conjugacy class of $G$, assembled over the inertia stack. The categorical content is the Bridgeland--King--Reid equivalence
\[
D^b\bigl([X/G]\bigr) \;\simeq\; D^b_G(X),
\]
and its CoHA-level incarnation: every local CoHA on $[X/G]$ pulls back to a $G$-equivariant CoHA on $X$, and the global Hall product is the $G$-invariant part after the inertia decomposition is unpacked.

Two examples organise the material. Local $\mathbb{P}^2 = [\mathbb{C}^3/\mathbb{Z}_3]^{\mathrm{res}}$ admits two compatible presentations: a single-chart skew-group presentation via the nine-arrow Beilinson quiver, and a three-chart \v{C}ech presentation on its crepant resolution $\mathrm{Tot}(K_{\mathbb{P}^2})$. The orbifold CoHA matches the chart-wise CoHA exactly when a compact-$4$-cycle Negu\c{t} wheel-correction is installed. Mathieu $M_{24}$ on $K3$ occupies the intersection of Stratum (iii) with Stratum (iv): the orbifold gluing against a sporadic simple group is realised by twined elliptic genera whose weight-$1/2$ mock modular coefficients are the CoHA characters in each twisted sector.

\section{Orbifold inertia and the Bridgeland--King--Reid equivalence}
\label{sec:orbifold-inertia-bkr}

\subsection{The inertia stack, elementarily}
\label{subsec:inertia-stack}

Let $G \subset \mathrm{SU}(3)$ be finite, acting on $\mathbb{C}^3 = \mathrm{Spec}\,\mathbb{C}[x,y,z]$ preserving the holomorphic volume form $dx \wedge dy \wedge dz$. The quotient stack $[\mathbb{C}^3/G]$ has coherent sheaves presented as $G$-equivariant sheaves on $\mathbb{C}^3$:
\[
\mathrm{Coh}\bigl([\mathbb{C}^3/G]\bigr) \;\simeq\; \mathrm{Coh}_G(\mathbb{C}^3).
\]
This identification is elementary: a coherent sheaf on the stack $[\mathbb{C}^3/G]$ pulls back to a coherent sheaf on the atlas $\mathbb{C}^3$, equipped with descent data for the groupoid $G \times \mathbb{C}^3 \rightrightarrows \mathbb{C}^3$, and descent data for a groupoid whose unit arrows are group elements is precisely $G$-equivariant structure.

The \emph{inertia stack} $I([\mathbb{C}^3/G])$ is the stack of pairs $(p, g)$ with $p \in [\mathbb{C}^3/G]$ a point and $g \in \mathrm{Aut}(p)$ an automorphism. For an orbifold point $p$ represented by $x \in \mathbb{C}^3$, the automorphism group $\mathrm{Aut}(p)$ is the stabiliser $G_x \subset G$, and an inertia datum is an element $g \in G_x$ --- equivalently, a pair $(x, g)$ with $g \cdot x = x$. The inertia stack is the fibre product $[\mathbb{C}^3/G] \times_{[\mathbb{C}^3/G] \times [\mathbb{C}^3/G]} [\mathbb{C}^3/G]$ along the diagonal, which rewrites as the quotient $[\{(x, g) : g \cdot x = x\} / G]$ with $G$ acting by simultaneous conjugation $h \cdot (x, g) = (h \cdot x, h g h^{-1})$.

Slicing this groupoid by $G$-conjugacy on the second factor: pick a representative $g_0$ of each conjugacy class $[g] \in \mathrm{Conj}(G)$, and restrict to the slice $\{(x, g_0) : g_0 \cdot x = x\} = \mathbb{C}^{3, g_0}$ (the fixed locus of $g_0$). The residual $G$-action on this slice is by the centraliser $C_G(g_0) \subset G$ (elements commuting with $g_0$). Summing over conjugacy classes,
\[
I\bigl([\mathbb{C}^3/G]\bigr) \;=\; \bigsqcup_{[g] \in \mathrm{Conj}(G)} \bigl[\mathbb{C}^{3, g}/C_G(g)\bigr].
\]
The identity class $[e]$ yields $[\mathbb{C}^{3, e}/G] = [\mathbb{C}^3/G]$, the \emph{untwisted sector}. Each non-identity class $[g] \neq [e]$ yields a \emph{twisted sector} indexed by $g$. For $G$ abelian, every element is its own conjugacy class, $C_G(g) = G$, and the decomposition becomes $\bigsqcup_{g \in G} [\mathbb{C}^{3, g}/G]$.

This decomposition is the arithmetic of Dixon--Harvey--Vafa--Witten orbifold cohomology as refined by Chen--Ruan: the orbifold cohomology of $[\mathbb{C}^3/G]$ is
\[
H^\bullet_{\mathrm{orb}}\bigl([\mathbb{C}^3/G]\bigr) \;=\; \bigoplus_{[g]} H^{\bullet - 2\,\mathrm{age}(g)}\bigl(\mathbb{C}^{3,g}/C_G(g)\bigr),
\]
with age-grading shift $\mathrm{age}(g) = \sum_i a_i$ where $g$ acts on a tangent frame by $\mathrm{diag}(e^{2\pi i a_1}, e^{2\pi i a_2}, e^{2\pi i a_3})$ with $a_i \in [0, 1)$. The Calabi--Yau condition $\sum_i a_i \in \mathbb{Z}$ keeps age integral.

\subsection{Skew-group algebra and Morita equivalence}
\label{subsec:skew-group-morita}

The skew-group algebra $R \# G$ with $R = \mathbb{C}[x,y,z]$ is
\[
R \# G \;=\; R \otimes \mathbb{C}[G] \quad\text{with multiplication}\quad (r \otimes g)(r' \otimes g') = r \cdot g(r') \otimes gg'.
\]
Equivalently: $R \# G$ is the free $R$-module with basis $G$, and multiplication records the twist of the action. This is not an arbitrary algebra; it is the arrow algebra of the action groupoid $G \times \mathbb{C}^3 \rightrightarrows \mathbb{C}^3$, and its representation category is exactly $G$-equivariant sheaves:
\[
(R \# G)\text{-mod} \;\simeq\; \mathrm{Coh}_G(\mathbb{C}^3) \;\simeq\; \mathrm{Coh}\bigl([\mathbb{C}^3/G]\bigr).
\]
The first equivalence is the passage from arrow-algebra modules to equivariant modules; the second is the atlas identification from Section~\ref{subsec:inertia-stack}. Morita equivalence of $R \# G$ with its basic algebra (the endomorphism ring of a minimal projective generator) is the passage from infinite-dimensional to quiver-presented form.

\subsection{Bridgeland--King--Reid, first-principles sketch}
\label{subsec:bkr-first-principles}

The Bridgeland--King--Reid theorem (math/9908027) asserts: for $G \subset \mathrm{SL}(n, \mathbb{C})$ finite, $n \leq 3$, with $Y := G\text{-}\mathrm{Hilb}(\mathbb{C}^n)$ the equivariant Hilbert scheme of $G$-clusters, the Fourier--Mukai transform with kernel the universal $G$-cluster $\mathcal{Z} \subset Y \times \mathbb{C}^n$ is an equivalence
\[
\Phi_{\mathcal{Z}} \colon D^b(Y) \;\xrightarrow{\sim}\; D^b_G(\mathbb{C}^n) \;\simeq\; D^b\bigl([\mathbb{C}^n/G]\bigr).
\]
The first-principles content:
\begin{itemize}
  \item \emph{Crepancy:} For $G \subset \mathrm{SL}(n)$ with $n \leq 3$, $Y$ is smooth and the canonical class $K_Y = 0$; equivalently, the resolution $Y \to \mathbb{C}^n/G$ is crepant. The CY structure on $[\mathbb{C}^n/G]$ transports to a CY structure on $Y$.
  \item \emph{Equivariant Hilbert scheme:} A \emph{$G$-cluster} is a $G$-invariant $0$-dimensional subscheme $Z \subset \mathbb{C}^n$ with $H^0(\mathcal{O}_Z) \cong \mathbb{C}[G]$ as $G$-representation. The moduli $G\text{-}\mathrm{Hilb}(\mathbb{C}^n)$ parametrising these is the natural candidate resolution.
  \item \emph{Universal kernel:} The universal family $\mathcal{Z} \subset Y \times \mathbb{C}^n$ is the incidence locus. Its structure sheaf, equipped with its $G$-equivariant structure via the second projection, is an object of $D^b(Y \times_G \mathbb{C}^n) \simeq D^b_G(Y \times \mathbb{C}^n)$.
  \item \emph{Bridgeland--King--Reid criterion:} The Fourier--Mukai transform with kernel $\mathcal{O}_{\mathcal{Z}}$ is an equivalence if and only if the fibre product $\mathcal{Z} \times_Y \mathcal{Z}$ has dimension at most $n - 1$. This is the Bondal--Orlov criterion for FM-equivalence. The authors verify it for $n \leq 3$ by explicit fibrewise analysis.
  \item \emph{CoHA transport:} The Hall product on $D^b(Y)$ pulls back along $\Phi_{\mathcal{Z}}$ to a Hall product on $D^b_G(\mathbb{C}^n)$, which in turn decomposes via the inertia splitting into a conjugacy-class-indexed direct sum of twisted-sector CoHAs. This is the CoHA shadow of Chen--Ruan orbifold cohomology.
\end{itemize}

The BKR equivalence is thus the gluing cocycle for Stratum (iii): it assembles the $G$-equivariant critical CoHA on $\mathbb{C}^n$ with the $\mathbb{C}^3$-Jordan-triple potential $W = \mathrm{tr}(x[y,z])$ into the CoHA of the crepant resolution $Y$, with the inertia decomposition recording the conjugacy-class-graded structure. On the resolution side, $Y$ is a smooth toric CY$_3$ when $G$ is abelian, and the chart-wise Stratum-(i) machinery applies; on the orbifold side, the skew-group presentation $R \# G$ provides a single global datum. The BKR kernel is the gluing cocycle reconciling the two.

\paragraph{Class-group obstruction to McKay presentation.}
The applicability of BKR is governed by an arithmetic invariant of the singular base $X_0 = \mathbb{C}^n/G$: its divisor class group $\mathrm{Cl}(X_0)$. For $X_0 = \mathbb{C}^n/G$ with $G \subset \mathrm{SL}(n, \mathbb{C})$ finite, $\mathrm{Cl}(X_0) = G^{\mathrm{ab}}$, finite of order dividing $|G|$. Any Gorenstein singularity whose class group is \emph{infinite} is not of the form $\mathbb{C}^n/G$ and lies outside BKR's scope. The conifold $X_0^{\mathrm{con}} = \{xy - zw = 0\} \subset \mathbb{C}^4$ is the canonical outlier: its class group $\mathrm{Cl}(X_0^{\mathrm{con}}) = \Z$ is infinite, generated by the small-resolution Weil divisor $\{x = z = 0\}$, so the conifold admits no McKay orbifold presentation; its NCCR is the Klebanov--Witten two-node Kronecker quiver with quartic potential $W_{\mathrm{con}} = \mathrm{tr}(a_1 b_1 a_2 b_2 - a_1 b_2 a_2 b_1)$ (Szendr\H{o}i \texttt{arXiv:0705.3419}), not the Beilinson nine-arrow quiver with cubic potential. The degree of the NCCR potential is the brane-tiling invariant distinguishing the two regimes (Hanany--Kennaway \texttt{hep-th/0503149}; Franco--Hanany--Kennaway--Vegh--Wecht \texttt{hep-th/0511063}): McKay-orbifold quotients carry cubic Jordan-triple potentials on the $\C^3$-atlas; complete-intersection singularities such as the conifold carry quartic Kronecker-type potentials globally.

\section{Local \texorpdfstring{$\mathbb{P}^2$}{P2} as \texorpdfstring{$[\mathbb{C}^3/\mathbb{Z}_3]$}{[C3/Z3]} and the nine-arrow quiver}
\label{sec:local-P2-orbifold}

Local $\mathbb{P}^2$ is $\mathrm{Tot}(\mathcal{O}_{\mathbb{P}^2}(-3))$, the total space of the canonical bundle on $\mathbb{P}^2$. As a toric CY$_3$, it has a compact $4$-cycle (the zero-section $\mathbb{P}^2$ itself) and is not a union of $\mathbb{C}^3$-charts without corrections. The orbifold route produces it as a crepant resolution $\mathrm{Tot}(K_{\mathbb{P}^2}) = [\mathbb{C}^3/\mathbb{Z}_3]^{\mathrm{res}}$, and the CoHA story becomes transparent through the Beilinson nine-arrow quiver.

\subsection{The \texorpdfstring{$\mathbb{Z}_3$}{Z3} action and its age data}
\label{subsec:Z3-action}

Let $\zeta = e^{2\pi i/3}$. The group $\mathbb{Z}_3 = \langle \sigma \rangle$ acts on $\mathbb{C}^3$ by
\[
\sigma \cdot (x, y, z) \;=\; (\zeta x, \zeta y, \zeta z).
\]
The weights of $\sigma$ on the tangent frame are $(1/3, 1/3, 1/3)$, with age
\[
\mathrm{age}(\sigma) = 1/3 + 1/3 + 1/3 = 1 \;\in\; \mathbb{Z},
\]
confirming $\sigma \in \mathrm{SU}(3)$ and the CY orbifold condition. The three conjugacy classes $\{e, \sigma, \sigma^2\}$ (abelian, every element its own class) give ages $0, 1, 2$ respectively. The fixed loci are $\mathbb{C}^{3,e} = \mathbb{C}^3$, $\mathbb{C}^{3,\sigma} = \{0\}$, $\mathbb{C}^{3,\sigma^2} = \{0\}$, and the inertia decomposition is
\[
I\bigl([\mathbb{C}^3/\mathbb{Z}_3]\bigr) \;=\; [\mathbb{C}^3/\mathbb{Z}_3] \;\sqcup\; B\mathbb{Z}_3 \;\sqcup\; B\mathbb{Z}_3.
\]
The untwisted sector carries the ambient CY$_3$ structure; the two twisted sectors are point classes supported at the origin with age-shifts $1$ and $2$. McKay identifies them with the non-trivial cohomology of the exceptional divisor on the crepant resolution: $H^0(\bP^2_{\mathrm{exc}}) = \C$ absorbs into the untwisted origin, while $H^2(\bP^2_{\mathrm{exc}}) = \C$ and $H^4(\bP^2_{\mathrm{exc}}) = \C$ match the age-$1$ and age-$2$ twisted-sector point-contributions respectively. This is the Batyrev--Denef--Loeser stringy-cohomology shadow of BKR: $\chi_{\mathrm{stringy}}([\C^3/\Z_3]) = 3 = \chi(\bP^2_{\mathrm{exc}})$, the ranks $(1,1,1)$ of the three cohomology groups appearing on both sides.

\subsection{The Beilinson nine-arrow quiver}
\label{subsec:beilinson-quiver}

The skew-group algebra $\mathbb{C}[x,y,z] \# \mathbb{Z}_3$ has three indecomposable projectives corresponding to the three irreducible representations $\rho_0, \rho_1, \rho_2$ of $\mathbb{Z}_3$ (characters $1, \zeta, \zeta^2$). Its basic algebra is the path algebra of the quiver $Q$ with three vertices $\{0, 1, 2\}$ (one for each character) and nine arrows
\[
X_i, Y_i, Z_i : i \longrightarrow i + 1 \pmod 3, \qquad i \in \Z/3,
\]
one triple $(X_i, Y_i, Z_i)$ for each vertex, corresponding to the three coordinates $(x, y, z)$ of $\mathbb{C}^3$. Because $\sigma$ acts on each coordinate by the same weight $\zeta$, each coordinate raises the $\rho$-degree by $+1 \bmod 3$, producing the cyclic $0 \to 1 \to 2 \to 0$ arrow pattern with three arrows per step.

The Jacobi potential is the cubic
\[
W \;=\; \sum_{i \in \Z/3} \mathrm{tr}\bigl(X_i \, Y_{i+1} \, Z_{i+2} \;-\; X_i \, Z_{i+1} \, Y_{i+2}\bigr),
\]
the cyclic unfolding of $\mathrm{tr}(x [y, z])$ across the three vertex-types, canonically fixed by the $\Z_3$-equivariant Jordan-triple structure on $\mathbb{C}^3$ (see toric-CoHA Theorem on local $K_{\bP^2}$). The pair $(Q, W)$ is the \emph{nine-arrow Beilinson quiver with cubic potential}; the Jacobi algebra $\mathrm{Jac}(Q, W)$ is the Morita-reduction of $R \# \mathbb{Z}_3$. The potential is cubic, not quartic: the rank-one Jordan-triple structure of $\mathbb{C}^3$ gives degree-three cycles through the three vertices (an $X$, a $Y$, a $Z$), in contrast to the conifold's rank-two Kronecker setup where cycles pass through two nodes and close at length four.

The equivalence with Beilinson's derived description of $\mathbb{P}^2$:
\[
D^b(\mathbb{P}^2) \;\simeq\; D^b\bigl(\mathrm{End}(\mathcal{O} \oplus \mathcal{O}(1) \oplus \mathcal{O}(2))\text{-mod}\bigr),
\]
is the classical Be\u{\i}linson resolution. The three line bundles $\mathcal{O}, \mathcal{O}(1), \mathcal{O}(2)$ are the three indecomposable summands of the tilting bundle; the nine arrows are global sections $H^0(\mathbb{P}^2, \mathcal{O}(1)) = \langle x, y, z\rangle$ between consecutive summands. Under BKR, the Be\u{\i}linson tilting bundle on $\mathbb{P}^2$ pulls up along the projection $\mathrm{Tot}(K_{\mathbb{P}^2}) \to \mathbb{P}^2$ to a tilting bundle on local $\mathbb{P}^2$ whose endomorphism ring is precisely the Jacobi algebra $\mathrm{Jac}(Q, W)$.

\subsection{The CoHA via Schiffmann--Vasserot equivariantised}
\label{subsec:SV-equivariantised}

The Schiffmann--Vasserot theorem constructs the CoHA of $\mathbb{C}^3$ (with the Jordan-triple potential $W_{\mathbb{C}^3} = \mathrm{tr}(x[y,z])$) as the positive half $Y^+(\widehat{\mathfrak{gl}}_1)$ of the affine Yangian of $\mathfrak{gl}_1$:
\[
\mathcal{H}_{\mathbb{C}^3} \;=\; \bigoplus_n H^{T^3}_\bullet\bigl(\mathrm{crit}(W)|_{\mathcal{M}_n^{\mathbb{C}^3}}\bigr) \;\cong\; Y^+\bigl(\widehat{\mathfrak{gl}}_1\bigr).
\]
Equivariantising by the $\mathbb{Z}_3$-action gives the CoHA of $[\mathbb{C}^3/\mathbb{Z}_3]$:
\[
\mathcal{H}_{[\mathbb{C}^3/\mathbb{Z}_3]} \;=\; \bigl(\mathcal{H}_{\mathbb{C}^3}\bigr)^{\mathbb{Z}_3} \;\oplus\; \mathcal{H}_{\mathbb{C}^3}^{\mathrm{tw}, \sigma} \;\oplus\; \mathcal{H}_{\mathbb{C}^3}^{\mathrm{tw}, \sigma^2},
\]
where the second and third summands are the twisted-sector contributions indexed by the non-trivial conjugacy classes. The twisted sectors are modules over the untwisted invariant ring $Y^+(\widehat{\mathfrak{gl}}_1)^{\mathbb{Z}_3}$, in a manner that matches the age-graded decomposition. Explicitly: the equivariant parameters $(\epsilon_1, \epsilon_2, \epsilon_3)$ of $T^3$ pick up a $\mathbb{Z}_3$-equivariant refinement; the $\mathbb{Z}_3$-invariant combination $\epsilon_1 + \epsilon_2 + \epsilon_3 = 0$ (the CY relation) survives, and the twisted-sector generators carry additional $\zeta$-charge compatible with age.

\subsection{The chart-wise cover of the crepant resolution}
\label{subsec:chartwise-crepant}

Local $\mathbb{P}^2$ as a toric variety has fan $\Sigma$ with three top-dimensional cones, corresponding to the three torus-fixed points of $\mathbb{P}^2$, each carrying a local $\mathbb{C}^3$-chart (the affine toric patch trivialising $K_{\mathbb{P}^2}$ over that chart). Label the charts $U_0, U_1, U_2$; each is $\cong \mathbb{C}^3$ with local coordinates $(u, v, w)$ such that the Jordan-triple potential $\mathrm{tr}(u[v,w])$ holds on the chart.

The atlas reconciles with the orbifold presentation as follows: on the singular base $\mathbb{C}^3/\mathbb{Z}_3$ the resolution $\pi : \mathrm{Tot}(K_{\mathbb{P}^2}) \to \mathbb{C}^3/\mathbb{Z}_3$ is an isomorphism away from the unique singular point $0 \in \mathbb{C}^3/\mathbb{Z}_3$, and the preimage $\pi^{-1}(0) = \mathbb{P}^2$ is the exceptional divisor. The three toric charts $U_0, U_1, U_2$ cover $\mathrm{Tot}(K_{\mathbb{P}^2})$ with pairwise overlaps $\cong \mathbb{G}_m \times \mathbb{C}^2$ and triple overlap the complement of a point. Each chart $U_i$ intersects the exceptional $\mathbb{P}^2$ in an $\mathbb{A}^2$-patch centred at a torus-fixed point of $\mathbb{P}^2$, and the three patches assemble to the standard toric $\mathbb{A}^2$-atlas of $\mathbb{P}^2$. The chart decomposition on the resolution side therefore reads: a single orbifold chart at the origin (the $[\mathbb{C}^3/\mathbb{Z}_3]$ stacky chart covering a neighbourhood of $0$) plus three smooth toric $\mathbb{C}^3$-patches resolving the $\mathbb{Z}_3$-singularity along the exceptional $\mathbb{P}^2$.

The gluing cocycle on the overlaps $U_i \cap U_j$ is the toric transition; the Schiffmann--Vasserot CoHA lives on each chart as $Y^+(\widehat{\mathfrak{gl}}_1)$, and the multi-chart CoHA is the homotopy limit over the \v{C}ech nerve. Because of the compact $\mathbb{P}^2$-divisor, this limit picks up \emph{internal-lattice contributions}: the Negu\c{t} wheel-condition corrections. Concretely, the D0-brane moduli on local $\mathbb{P}^2$ contain not just the torus-fixed points $U_0, U_1, U_2$ but also $1$-parameter families of brane configurations \emph{inside} the compact $\mathbb{P}^2$-divisor, whose cohomological contribution is non-trivial and must be integrated separately.

\subsection{Reconciling the two routes}
\label{subsec:reconcile-p2}

The orbifold CoHA via Schiffmann--Vasserot equivariantised on $[\mathbb{C}^3/\mathbb{Z}_3]$, and the chart-wise CoHA on the three-chart atlas of local $\mathbb{P}^2$, agree under the BKR-transport. The reconciliation proceeds through:
\begin{enumerate}
  \item BKR gives $D^b(\text{local } \mathbb{P}^2) \simeq D^b_{\mathbb{Z}_3}(\mathbb{C}^3)$ as triangulated categories.
  \item The CoHA Hall product is preserved, because Hall multiplication is defined from the extension stack $\mathrm{Ext}(\mathcal{E}, \mathcal{F})$, and extensions are preserved by BKR as an equivalence of derived categories.
  \item The orbifold side's inertia-graded decomposition $\mathcal{H}_{\mathrm{orb}} = \mathcal{H}_{\mathrm{untw}} \oplus \mathcal{H}_{\sigma} \oplus \mathcal{H}_{\sigma^2}$ matches the chart-wise side's cohomology-grading via the toric fan's three top-cones: each cone contributes one summand, with the Negu\c{t} wheel-condition corrections encoding exactly the age-shifted contributions of the twisted sectors.
  \item Explicitly: the age-$1$ twisted sector on the orbifold side corresponds to the contribution from the exceptional $\mathbb{P}^2$ of the resolution in the chart-wise route; age-$2$ corresponds to the Poincar\'e-dual class of $\mathbb{P}^2$.
\end{enumerate}

This is the first worked example of stratum overlap (i)$\cap$(iii): local $\mathbb{P}^2$ is simultaneously toric (stratum i) and orbifold (stratum iii), and the two gluing cocycles agree through BKR. The skew-group presentation via the nine-arrow quiver is the most economical form; the chart-wise \v{C}ech presentation with Negu\c{t} corrections is the most geometric form. Both compute the same CoHA, equal to the module-theoretic Jacobi-algebra CoHA of $(Q, W)$.

\section{Mathieu \texorpdfstring{$M_{24}$}{M24} on K3}
\label{sec:mathieu-m24}

The Mathieu moonshine of Eguchi--Ooguri--Tachikawa \texttt{arXiv:1004.0956} is an orbifold gluing not against a single finite group action on a local model, but against a sporadic simple group $M_{24}$ acting by symmetries of the elliptic genus of a K3 surface. This sits at the intersection of Stratum (iii) orbifold and Stratum (iv) lattice-polarised: $M_{24}$ emerges from the Mukai lattice $\Lambda_{K3} = \mathrm{II}_{4, 20}$ via Mukai's classification (Mukai 1988, \emph{Invent.\ Math.}\ 94) of symplectic automorphism groups, and the assembly of twined elliptic genera is indexed by the $M_{24}$-conjugacy classes with each $\phi_g$ a weak Jacobi form on $\mathbb{H} \times \mathbb{C}$ whose mock modular $\mathcal{N} = 4$-massive component is a weight-$1/2$ mock form on $\mathbb{H}$.

\subsection{The Eguchi--Ooguri--Tachikawa observation}
\label{subsec:eot-observation}

The K3 elliptic genus is the weight-$0$ index-$1$ weak Jacobi form
\[
\phi_{K3}(\tau, z) \;=\; 2 \, \phi_{0, 1}(\tau, z) \;=\; 8 \, \sum_{\alpha = 2, 3, 4} \Bigl(\frac{\theta_\alpha(\tau, z)}{\theta_\alpha(\tau, 0)}\Bigr)^2,
\]
with $\phi_{0, 1}$ the generator of $J_{0, 1}$ normalised by Eichler--Zagier, and the factor of two encoding the Euler characteristic $\chi(K3) = 24 = \phi_{K3}(\tau, 0)$ (Eguchi--Ooguri--Taormina--Yang 1989, \emph{Nucl.\ Phys.\ B} 315).
The $\mathcal{N}=4$ superconformal character-decomposition of $\phi_{K3}$ reads
\[
\phi_{K3}(\tau, z) \;=\; 24 \cdot \chi_{1/4, 0}(\tau, z) \;+\; \Sigma(\tau) \cdot \chi_{1/4, 1/2}(\tau, z),
\]
where $\chi_{1/4, h}$ are the $\mathcal{N}=4$ Ramond characters at central charge $c = 6$ and conformal weight $h$, and
\[
\Sigma(\tau) \;=\; 2 q^{-1/8}\bigl(-1 + 45 q + 231 q^2 + 770 q^3 + 2277 q^4 + \cdots\bigr).
\]
The numerical observation of Eguchi--Ooguri--Tachikawa: the Fourier coefficients $\{45, 231, 770, 2277, \dots\}$ are dimensions of irreducible representations of $M_{24}$.

\subsection{Twined elliptic genera and mock modularity}
\label{subsec:twined-mock}

For each $g \in M_{24}$, define the \emph{twined elliptic genus}
\[
\phi_g(\tau, z) \;=\; \mathrm{tr}_{\mathcal{H}_{K3}}\bigl(g \cdot (-1)^{F_L + F_R} q^{L_0 - c/24} \overline{q}^{\bar{L}_0 - \bar{c}/24} y^{J_0}\bigr),
\]
with $\mathcal{H}_{K3}$ the Hilbert space of the K3 sigma model, $F_{L/R}$ left/right fermion numbers, $J_0$ the $\mathrm{U}(1)$ current of the $\mathcal{N}=4$ superconformal algebra. The $g = e$ case recovers $\phi_{K3}$. For general $g$, the $\mathcal{N}=4$ decomposition gives
\[
\phi_g(\tau, z) \;=\; \chi_g \cdot \chi_{1/4, 0}(\tau, z) \;+\; \Sigma_g(\tau) \cdot \chi_{1/4, 1/2}(\tau, z),
\]
with $\chi_g$ the character of $g$ acting on $H^\bullet(K3)$ (a finite integer depending only on the $M_{24}$-conjugacy class) and $\Sigma_g$ a weight-$1/2$ \emph{mock modular form} of a specified shadow and multiplier system.

Cheng \texttt{arXiv:1005.5415} exhibited the twined $\Sigma_g$ and their $\Gamma_0(|g|)$ transformation laws; Cheng--Duncan--Harvey (\emph{Res.\ Math.\ Sci.}\ 1, \texttt{arXiv:1204.2779}) packaged the result into the umbral-moonshine programme, assigning each conjugacy class of $M_{24}$ a weight-$1/2$ mock modular form with specified shadow and multiplier. Gannon (\emph{Commun.\ Number Theory Phys.}\ 10, \texttt{arXiv:1211.5531}) proved positivity and uniqueness of the $M_{24}$-representation decomposition of the Fourier coefficients. Gaberdiel--Hohenegger--Volpato \texttt{arXiv:1008.0725} established that $M_{24}$ acts on the BPS states of the K3 sigma model by exhausting the possible symmetry groups of consistent $\mathcal{N} = (4, 4)$ K3 sigma models and matching them against the Mukai classification.

\subsection{CoHA twisted by an \texorpdfstring{$M_{24}$}{M24}-conjugacy class}
\label{subsec:coha-m24-twist}

For each conjugacy class $[g] \in M_{24}$, there is a twisted CoHA $\mathcal{H}_{K3}^g$ whose graded character is the twined elliptic genus $\phi_g$ after $\mathcal{N}=4$-decomposition. The twist-structure:
\begin{itemize}
  \item The untwisted CoHA $\mathcal{H}_{K3}^e$ has graded character $\phi_{K3}$, with BPS multiplicity $\phi_{K3}(\tau, 0) = \chi(K3) = 24$ in the massless $\mathcal{N} = 4$ sector and non-BPS multiplicities $2 \cdot \{45, 231, 770, 2277, \ldots\}$ in the massive sector (from $\Sigma(\tau)$), both encoding the vertex-algebra structure of the K3 Yangian.
  \item A twisted CoHA $\mathcal{H}_{K3}^g$ has graded character $\phi_g$ and interprets the $g$-fixed BPS states: for $g \in M_{24}$ of cycle shape $\prod_i k_i^{n_i}$, the character $\chi_g(K3) = \phi_g(\tau, 0) \in \Z$ is the Frame-shape-graded Euler characteristic, and the mock modular discrepancy $\Sigma_g$ encodes the twisted-sector shadow.
  \item The full orbifold CoHA on $[K3 / M_{24}]$ assembles as
  \[
  \mathcal{H}_{K3}^{M_{24}\text{-orb}} \;=\; \bigoplus_{[g] \in \mathrm{Conj}(M_{24})} \bigl(\mathcal{H}_{K3}^g\bigr)^{C_{M_{24}}(g)},
  \]
  with $C_{M_{24}}(g)$ the centraliser of $g$ in $M_{24}$, acting on the twisted sector to extract the invariant piece. The cardinality of $\mathrm{Conj}(M_{24})$ is $26$, matching the table of $M_{24}$-conjugacy classes.
\end{itemize}

\subsection{Stratum (iii)$\cap$(iv) and automorphic assembly}
\label{subsec:m24-strata-overlap}

Mathieu $M_{24}$ on K3 sits in stratum (iii) because K3 has no global $M_{24}$-action on the variety itself but carries a derived $M_{24}$-symmetry of its $\mathcal{N}=4$ chiral algebra; the symmetry acts on the BPS Hilbert space (equivalently, on the CoHA). It sits in stratum (iv) because the twined elliptic genera $\phi_g$ are weak Jacobi forms on the K3 lattice $\mathrm{II}_{4,20}$, and their assembly into a consistent system is an automorphic datum.

The two strata mesh through the Gritsenko lift. The K3 elliptic genus $\phi_{K3} = 2 \phi_{0,1}$ lifts via Borcherds--Gritsenko to the paramodular form $\Delta_5 \in M_5(\mathrm{Sp}_4(\mathbb{Z}))$ on the Siegel upper half-space; this is the stratum-(iv) face. Each twined elliptic genus $\phi_g$ lifts analogously to a Borcherds--Gritsenko product $\Delta_g \in M_{k_g}(\Gamma_g)$ on a congruence Siegel modular group, and these $\{\Delta_g\}_{[g] \in \mathrm{Conj}(M_{24})}$ span the gluing cocycle for the $M_{24}$-orbifold CoHA. On the CHL sub-sequence $N \in \{1, 2, 3, 4, 6\}$ the weights are $k_N = c_N(0)/2 \in \{5, 4, 3, 2, 1\}$ (Gritsenko 1999, \emph{Abh.\ Math.\ Sem.\ Hamburg} 69, Thm.~6.1; Borcherds 1998, \emph{Invent.\ Math.}\ 132 Thm.~10.1), so $\kappa_{\mathrm{BKM}}(\Phi_N) = c_N(0)/2$ across the ladder, with $N = 1$ recovering $\Delta_e = \Delta_5$ at cycle shape $1^{24}$. The full Gritsenko--Cl\'ery 8-form atlas carries weights $(5, 2, 1, 1, 1/2, 1, 1/4, 0)$ (Gritsenko--Cl\'ery 2013, \emph{J.\ Reine Angew.\ Math.}\ 678, \S 3), half-integral and fractional entries lifted under double-cover paramodular groups.

The stratum overlap is therefore concrete: the orbifold side is the assembly
\[
\mathcal{H}_{K3}^{M_{24}\text{-orb}} \;=\; \bigoplus_{[g]} \mathcal{H}_{K3}^g,
\]
and the lattice-polarised side is the assembly
\[
\{\Delta_g\}_{[g]} \;\subset\; \bigoplus_{[g]} M_{k_g}(\Gamma_g),
\]
and the \emph{twining consistency} between the two sides is the assertion that the Fourier-coefficient matrix of $\{\Delta_g\}$ matches the $M_{24}$-character table evaluated on the BPS-spectrum graded dimensions. This is the content of Mathieu moonshine at the level of Vol~III CoHA theory.

\subsection{Weight-\texorpdfstring{$1/2$}{1/2} mock modular shadow and the $M_{24}$-twining cocycle}
\label{subsec:mock-shadow-twining}

The shadow $S_g$ of $\Sigma_g$ is a weight-$3/2$ holomorphic cusp form on $\Gamma_0(|g|)$ fixed by Zwegers' completion: $\widehat\Sigma_g(\tau, \bar\tau) = \Sigma_g(\tau) + R_g(\tau, \bar\tau)$ with $R_g$ the non-holomorphic Eichler-integral completion whose $\bar\partial$-derivative recovers $S_g(\tau)$ through the standard shadow-pairing $\bar\partial_{\bar\tau} R_g = (2i)^{1/2} \overline{S_g(\tau)} (\tau - \bar\tau)^{-3/2}$. For $g = e$, $S_e(\tau) \propto \eta(\tau)^3$ (the weight-$3/2$ generator of the shadow space; Eguchi--Hikami 2009, \emph{J.\ Phys.\ A} 42). For general $g$, $S_g$ is a specific linear combination of $\eta$-quotients and Jacobi theta series indexed by the centraliser structure of $g$ (Cheng--Duncan--Harvey 2014).

The shadow encodes the failure of $\Sigma_g$ to be modular: $\Sigma_g$ transforms under $\Gamma_0(|g|)$ modulo the lift of $S_g$ across $\mathbb{H}$. In CoHA terms, the shadow records the deviation from a holomorphic global assembly: the $M_{24}$-twined CoHA is only mock-modularly consistent, with a precise shadow controlling the obstruction.

The twining cocycle is the assignment $g \mapsto (\phi_g, \Sigma_g, S_g)$ compatible with $M_{24}$-conjugation: $\phi_{hgh^{-1}} = \phi_g$, $\Sigma_{hgh^{-1}} = \Sigma_g$, etc. This cocycle is the orbifold-gluing datum of Stratum (iii) refined by the mock-modular Stratum (iv) data.

\subsection{K3 CoHA as Yangian of \texorpdfstring{$\mathfrak{osp}(4|20)$}{osp(4|20)}}
\label{subsec:k3-super-yangian}

The K3 CoHA is conjecturally the super-Yangian $Y_{\mathfrak{osp}}(4|20) = Y(\mathfrak{gl}(4|20))^\theta$, with $(4, 20)$ the signature of the Mukai lattice's Hodge involution-eigenspace decomposition: the $+1$-eigenspace is $\widetilde{H}^{1,1}(K3, \mathbb{C})$ of dimension $20$, and the $-1$-eigenspace is $(H^{2,0} \oplus H^{0,2} \oplus H^{0,0} \oplus H^{2,2})_{\mathbb{R}} = \mathbb{R}^4$. The super-Yangian $Y_{\mathfrak{osp}}(4|20)$ carries an $R$-matrix on $V = \mathbb{C}^{4|20}$ satisfying the graded Yang--Baxter equation; its positive half is the K3-CoHA.

The $M_{24}$-twining is compatible with this super-Yangian structure: conjugation by $g \in M_{24}$ preserves the $(4|20)$-supersplit (because $M_{24}$ acts by symplectic automorphisms of K3, preserving the Hodge involution), and the twisted sector $\mathcal{H}_{K3}^g$ is a $Y_{\mathfrak{osp}}(4|20)^{\langle g \rangle}$-module, with the centraliser $C_{M_{24}}(g)$-invariant part giving the global twisted CoHA.

\section{Summary: orbifold gluing in the four-stratum taxonomy}
\label{sec:orbifold-summary}

The orbifold stratum (iii) provides three distinct gluing modes, each illustrated in this part:
\begin{enumerate}
  \item \emph{Global abelian orbifold action}, e.g., $[\mathbb{C}^3/\mathbb{Z}_3]$ for local $\mathbb{P}^2$. Gluing cocycle: BKR equivalence with kernel the universal $G$-cluster on $G\text{-}\mathrm{Hilb}(\mathbb{C}^3)$. CoHA construction: either skew-group Schiffmann--Vasserot equivariantised, or chart-wise on the crepant resolution with Negu\c{t} wheel-corrections. The two agree as derived Hall structures.
  \item \emph{Finite non-abelian orbifold}, e.g., $[\mathbb{C}^3/G]$ for $G = \mathbb{Z}_n \times \mathbb{Z}_n$ or $G$ the binary tetrahedral / octahedral / icosahedral group. Gluing cocycle: BKR extended to non-abelian case; the inertia decomposition becomes non-trivial with twisted sectors indexed by conjugacy classes rather than individual elements.
  \item \emph{Sporadic / chiral orbifold symmetry}, e.g., Mathieu $M_{24}$ on K3. Gluing cocycle: weight-$1/2$ mock modular forms $\Sigma_g$ with shadow $S_g$, assembled into a twining system compatible with $M_{24}$-conjugation. This sits in Stratum (iii)$\cap$(iv): the orbifold group $M_{24}$ acts on the chiral-algebra side of $\Phi$, and the Borcherds--Gritsenko lifts $\{\Delta_g\}$ provide the automorphic gluing.
\end{enumerate}

Each example passes through the same categorical pattern: a single $G$-equivariant chart $X$ (or chiral-algebra action of $G$) carrying a CoHA structure, with the inertia-stack decomposition $I([X/G]) = \sqcup_{[g]} [X^g / C_G(g)]$ stratifying the global assembly into untwisted and twisted sectors. The global CoHA is the direct sum over conjugacy classes, each summand weighted by the age-shift and twisted-sector character. On the derived-category side, BKR (or its chiral generalisation for Mathieu) is the universal cocycle intertwining the orbifold presentation with the resolution presentation.

The K3 Yangian / $Y_{\mathfrak{osp}}(4|20)$ conjecture closes the frame: the super-Yangian's positive half is the K3-CoHA, the Mathieu $M_{24}$-action is the sporadic symmetry of this positive half, and the Borcherds-lift $\Delta_5$ is the resulting automorphic form on $\mathrm{Sp}_4(\mathbb{Z})$, with $\kappa_{\mathrm{BKM}}(\mathfrak{g}_{\Delta_5}) = 5$ the Gritsenko weight. These three pieces --- the super-Yangian, the $M_{24}$-orbifold twining, the $\Delta_5$ Borcherds lift --- are the three faces of a single gluing cocycle straddling Strata (ii), (iii), (iv).

The compact CY$_3$ manifold $K3 \times E$, treated in Part~VIII of this chapter, inherits all three faces and adds a fourth via the $E$-factor's translation equivariance. Part~III establishes the orbifold-gluing tools needed; Part~VIII assembles them against the other strata to produce the global CoHA of $K3 \times E$ with its seven $r_{\mathrm{CY}}$-faces.

\paragraph*{Transition to Part IV.}
The orbifold stratum admits a refinement: the gluing cocycle can also be presented as a \emph{derived Morita equivalence} by a tilting object whose endomorphism algebra is the skew-group algebra. Part~IV develops this Van den Bergh tilting-bundle viewpoint, with local $\mathbb{P}^2$ reappearing there as the three-line-bundle tilting example, and the Serre-equivariant quasi-NCCR of $K3 \times E$ emerging as the Morita-gluing obstruction that Part~III's orbifold methods alone cannot resolve.


% Vol III --- Gluing Chapter, Part IV: Van den Bergh derived Morita gluing.
% Architecture: notes/gluing_chapter_architecture_2026_04_23.md, Part IV (§§4.1--4.4).
% Primary sources:
%   Van den Bergh 2004 (arXiv:math/0211064), Non-commutative crepant resolutions
%   Bondal--Kapranov 1990 (Math. USSR Izv.), Enhanced triangulated categories
%   Bondal--Orlov 2002 (arXiv:math/0206295), Derived categories of coherent sheaves
%   King 1997, Moduli of representations of finite-dimensional algebras
%   Davison--Hennecart--Schlegel Mejia 2023 (arXiv:2303.12592), BPS Lie algebras
%   Davison 2017 (arXiv:1601.02479), Critical CoHA and PBW
%   Derksen--Weyman--Zelevinsky 2008 (arXiv:0704.0649), Quivers with potentials
%   Keller--Yang 2011 (arXiv:0906.0761), Derived equivalences from mutations
%   Szendroi 2008 (arXiv:0705.3419), Non-commutative DT on resolved conifold
%   Craw--Ishii 2004, and Craw 2008 on (-3)-curve tilting
%   Kontsevich--Soibelman 2008 (arXiv:0811.2435), Motivic DT and wall-crossing
%   Bridgeland 2006 (arXiv:math/0502050), Flops and derived equivalence
%   Ginzburg 2006 (arXiv:math/0612139), Calabi--Yau algebras

\providecommand{\C}{\mathbb{C}}
\providecommand{\Z}{\mathbb{Z}}
\providecommand{\N}{\mathbb{N}}
\providecommand{\Q}{\mathbb{Q}}
\providecommand{\R}{\mathbb{R}}
\providecommand{\bP}{\mathbb{P}}
\providecommand{\cO}{\mathcal{O}}
\providecommand{\cA}{\mathcal{A}}
\providecommand{\cC}{\mathcal{C}}
\providecommand{\cD}{\mathcal{D}}
\providecommand{\cE}{\mathcal{E}}
\providecommand{\cF}{\mathcal{F}}
\providecommand{\cG}{\mathcal{G}}
\providecommand{\cH}{\mathcal{H}}
\providecommand{\cM}{\mathcal{M}}
\providecommand{\cT}{\mathcal{T}}
\providecommand{\cU}{\mathcal{U}}
\providecommand{\cY}{\mathcal{Y}}
\providecommand{\Tot}{\mathrm{Tot}}
\providecommand{\Coh}{\mathrm{Coh}}
\providecommand{\Perf}{\mathrm{Perf}}
\providecommand{\Hom}{\mathrm{Hom}}
\providecommand{\End}{\mathrm{End}}
\providecommand{\Ext}{\mathrm{Ext}}
\providecommand{\RHom}{R\mathrm{Hom}}
\providecommand{\mod}{\mathrm{mod}}
\providecommand{\Mod}{\mathrm{Mod}}
\providecommand{\Spec}{\mathrm{Spec}}
\providecommand{\Proj}{\mathrm{Proj}}
\providecommand{\Jac}{\mathrm{Jac}}
\providecommand{\Crit}{\mathrm{Crit}}
\providecommand{\tr}{\mathrm{tr}}
\providecommand{\id}{\mathrm{id}}
\providecommand{\vphi}{\varphi}
\providecommand{\gl}{\mathfrak{gl}}
\providecommand{\fg}{\mathfrak{g}}
\providecommand{\ghat}{\widehat{\mathfrak{gl}}_{1}}
\providecommand{\Yplus}{Y^{+}}
\providecommand{\CoHA}{\mathrm{CoHA}}
\providecommand{\Hilb}{\mathrm{Hilb}}
\providecommand{\NCCR}{\mathrm{NCCR}}
\providecommand{\Mu}{\mathrm{M}}
\providecommand{\kch}{\kappa_{\mathrm{ch}}}
\providecommand{\kcat}{\kappa_{\mathrm{cat}}}
\providecommand{\kBKM}{\kappa_{\mathrm{BKM}}}
\providecommand{\kfib}{\kappa_{\mathrm{fibre}}}
\providecommand{\ClaimStatusTheorem}{\textsc{[theorem]}}
\providecommand{\ClaimStatusDefinition}{\textsc{[definition]}}
\providecommand{\ClaimStatusDerivation}{\textsc{[first-principles]}}
\providecommand{\ClaimStatusConjectured}{\textsc{[conjectural]}}
\providecommand{\ClaimStatusProposition}{\textsc{[proposition]}}

% =======================================================================
\part{Van den Bergh derived Morita gluing}
\label{part:gluing-vdb}
% =======================================================================

Parts I--III assembled global CoHAs from \v{C}ech, toric, and McKay atlases: each
overlap carried a geometric transition --- a Fourier--Mukai kernel, a toric
cocycle, a BKR equivalence. Part IV replaces the atlas by a single object. A
\emph{tilting bundle} $T$ on a projective CY$_3$ (or a local CY$_3$)
collapses the derived category of coherent sheaves onto the derived category of
modules over a finite-dimensional algebra $A = \End(T)$: all gluing
information on $X$ is encoded in the multiplicative structure of $A$. The Hall
product on critical cohomology transports along this collapse. Mutations of $T$
act as the chamber-changing automorphisms of the associated CoHA.

% =======================================================================
\section{Van den Bergh tilting bundles}
\label{sec:vdb-tilting}
% =======================================================================

\subsection*{The obstruction: $D^b(\Coh X)$ is not compactly generated from below}

For $X$ a smooth projective variety, $D^b(\Coh X)$ admits no canonical
generator. Locally one trivialises $\Coh X$ by $\Coh U \simeq \mod R_U$ with
$R_U = \Gamma(U,\cO_X)$ a commutative affine ring, but the local presentations
are incompatible: $R_U$ is Morita-trivial (its module category is generated by
the regular module $R_U$ itself), whereas $D^b(\Coh X)$ globally admits no such
single-generator description. A tilting bundle is the object that repairs this
failure \emph{globally}.

\begin{definition}[Tilting bundle; Bondal--Kapranov~1990]
\label{def:tilting-bundle}
\ClaimStatusDefinition\ Let $X$ be a smooth variety. A coherent sheaf
$T \in \Coh(X)$ is a \emph{tilting bundle} if
\begin{enumerate}
\item[(T1)] $T$ is locally free of finite rank;
\item[(T2)] $\Ext^{i}_{\cO_X}(T, T) = 0$ for all $i \geq 1$
           (``vanishing higher self-Ext'');
\item[(T3)] $T$ generates $D^{b}(\Coh X)$ as a triangulated category:
           the smallest full triangulated subcategory of $D^{b}(\Coh X)$ closed
           under shifts, cones, and direct summands and containing $T$ is
           $D^{b}(\Coh X)$ itself.
\end{enumerate}
The algebra $A = \End_{\cO_X}(T)^{\mathrm{op}}$ is called the \emph{tilting
algebra}.
\end{definition}

Conditions (T1)--(T2) force $\RHom(T, T) = \End(T) = A$ concentrated in degree
zero; condition (T3) is the generation property. The reason for passing to the
opposite algebra is that $\RHom(T, -)$ produces a \emph{right} $\End(T)$-module
and, by convention, one wishes to write these as left $A$-modules.

\begin{theorem}[Derived Morita via tilting; Bondal 1989; Rickard 1989 for algebras;
Bondal--Van den Bergh 2003, arXiv:\texttt{math/0204218} Thm.~3.1.1]
\label{thm:derived-morita-tilting}
\ClaimStatusTheorem\ Let $T$ be a tilting bundle on $X$ and $A = \End(T)^{\mathrm{op}}$.
The pair
\[
  L \;=\; T \otimes^{L}_{A} (-) \;:\; D(A\text{-}\mod) \;\longrightarrow\; D_{\mathrm{qc}}(X),
  \qquad
  R \;=\; \RHom_{\cO_X}(T, -) \;:\; D_{\mathrm{qc}}(X) \;\longrightarrow\; D(A\text{-}\mod)
\]
is an adjoint pair $L \dashv R$. Two hypotheses make it an equivalence after
restriction to the compact part:
\begin{enumerate}
\item[(H1)] (\emph{Compact generation}) $T$ is a compact generator of
$D_{\mathrm{qc}}(X)$: the smallest localising subcategory of $D_{\mathrm{qc}}(X)$
closed under arbitrary coproducts and containing $T$ is all of $D_{\mathrm{qc}}(X)$,
and $T$ is compact (Bondal--Van den Bergh \emph{loc.\ cit.}, Thm.~3.1.1);
\item[(H2)] (\emph{Higher Ext vanishing}) $\Ext^{>0}_{\cO_X}(T, T) = 0$,
so that the unit $A \to R(L(A)) = \RHom(T, T \otimes^L_A A) = \RHom(T, T)$ is an
isomorphism in $D(A\text{-}\mod)$ concentrated in degree zero.
\end{enumerate}
Under (H1)--(H2), $L \dashv R$ is a pair of mutually quasi-inverse equivalences
\[
  D_{\mathrm{qc}}(X) \;\simeq\; D(A\text{-}\mod), \qquad
  \Perf(X) \;\simeq\; \Perf(A) = D^{\mathrm{perf}}(A\text{-}\mod).
\]
Restricting to bounded coherent objects on the left and finitely presented
modules on the right yields the classical equivalence
$D^{b}(\Coh X) \simeq D^{b}(A\text{-}\mod)$.
\end{theorem}

\begin{proof}[Proof.]
The adjunction $L \dashv R$ is the standard tensor--Hom adjunction for the
$A$-$\cO_X$ bimodule $T$: for $M \in D(A\text{-}\mod)$ and $\cF \in D_{\mathrm{qc}}(X)$,
$\Hom_{D_{\mathrm{qc}}(X)}(T \otimes^L_A M, \cF) = \Hom_{D(A)}(M, \RHom_{\cO_X}(T, \cF))$
functorially. The unit $\eta_A \colon A \to R(L(A)) = \RHom(T, T)$ is an
isomorphism by (H2), which forces $\Ext^{>0}(T, T) = 0$ and identifies $\RHom(T,T)
\simeq \End(T) = A^{\mathrm{op}}$ concentrated in degree zero; since $A$ generates
$D(A\text{-}\mod)$, this upgrades to $\eta_M \colon M \xrightarrow{\sim} R(L(M))$
for all $M \in D(A\text{-}\mod)$. The counit $\varepsilon_\cF \colon L(R(\cF)) \to
\cF$ is an isomorphism on the compact generator $T$ by the same computation, and
$T$ compactly generates $D_{\mathrm{qc}}(X)$ by (H1), so $\varepsilon$ is an
isomorphism on all of $D_{\mathrm{qc}}(X)$ (Bondal--Van den Bergh
arXiv:\texttt{math/0204218}, Thm.~3.1.1: a compact generator with bounded Ext
produces an equivalence with modules over its endomorphism dg-algebra, here
concentrated in degree zero by (H2)).
\end{proof}

\subsection*{Non-commutative crepant resolutions}

The obstruction on singular $X$ is sharper: $X$ admits no smooth model within
its birational class without blowing up, but one may ask for a
\emph{non-commutative} smooth substitute.

\begin{definition}[NCCR; Van den Bergh~2004, arXiv:\texttt{math/0211064}]
\label{def:nccr}
\ClaimStatusDefinition\ Let $R$ be a commutative Noetherian normal Gorenstein
domain. A \emph{non-commutative crepant resolution} of $R$ is an $R$-algebra
$A = \End_R(M)$, for some reflexive $R$-module $M$, such that
\begin{enumerate}
\item[(N1)] $A$ is \emph{homologically homogeneous}: all simple $A$-modules have
the same projective dimension, equal to $\dim R$;
\item[(N2)] $A$ is a \emph{maximal Cohen--Macaulay} $R$-module;
\item[(N3)] $A$ is generically Azumaya and $M$ is a reflexive generator.
\end{enumerate}
\end{definition}

``Crepant'' means $\omega_A \cong A$ as $A$-bimodules: the canonical class is
trivial, matching the CY condition on $\Spec R$.

\begin{theorem}[Van den Bergh equivalence, arXiv:\texttt{math/0211064} Thm.~6.6.1
and \S8]
\label{thm:vdb-equivalence}
\ClaimStatusTheorem\ Let $Y \to X = \Spec R$ be a crepant resolution (in the
commutative sense) with $R$ as in Definition~\ref{def:nccr}, and suppose
$\dim R \leq 3$. Then there exists a tilting bundle $T_Y$ on $Y$ with
$\End(T_Y)^{\mathrm{op}} = A$ an NCCR of $R$ in the sense of
Definition~\ref{def:nccr}, and the adjoint pair $L \dashv R$ of
Theorem~\ref{thm:derived-morita-tilting} is an equivalence
\[
  D^{b}(\Coh Y) \;\simeq\; D^{b}(A\text{-}\mod).
\]
Both sides realise the same abstract triangulated category; when $R$ is
Gorenstein of dimension three with trivial dualising module, the
\emph{Ginzburg dg-lift} $\Gamma = \Gamma(Q, W)$ of $A$ (Ginzburg 2006,
arXiv:\texttt{math/0612139} Thm.~3.2.8) carries a CY$_{3}$ structure:
$\Gamma$ is exact $3$-Calabi--Yau as a dg-algebra, meaning the inverse
dualising complex satisfies
$\RHom_{\Gamma^e}(\Gamma, \Gamma \otimes \Gamma)[3] \simeq \Gamma$ in the
derived category of $\Gamma$-bimodules.

\emph{Scope of existence.} The hypothesis $\dim R \leq 3$ is sharp. Beyond
dimension three, the existence of an NCCR is not guaranteed and \emph{fails}
in the cases most of interest: projective K3 surfaces with generic polarisation
(where $H^{0,2}(X) \neq 0$ obstructs higher Ext-vanishing), the quintic and
other generic projective CY$_3$ without toric or orbifold presentation, and
compact hyperk\"ahler manifolds of dimension $\geq 4$. The Van den Bergh
existence theorem covers precisely the local three-dimensional crepant setting:
toric CY$_3$, orbifolds $\C^n/G$ with $G \subset SU(n)$ of dimension $n \leq 3$,
$(-3)$-curve contractions, and their relatives.
\end{theorem}

On the conifold $R = \C[x,y,z,w]/(xw - yz)$ the NCCR is the \emph{Klebanov--Witten
quiver algebra} with tilting bundle $T = \cO \oplus \cO(1)$ on the resolved
conifold $\tilde Y = \cO_{\bP^1}(-1) \oplus \cO_{\bP^1}(-1)$. On local $\bP^2 =
\Tot(\cO_{\bP^2}(-3))$ the NCCR is the nine-arrow Beilinson algebra, with
$T = \cO \oplus \cO(1) \oplus \cO(2)$. On $(-3)$-curve contractions
$\tilde Y \to X$ with $\tilde Y$ containing a $\bP^1$ of self-intersection
$(-1,-1,-1,-1)$ on its normal bundle, Craw's tilting bundles (Craw
2008, building on Craw--Ishii 2004) play the same role. In each case, an
explicit tilting object $T$ reduces all geometric CoHA questions on $X$ to
finite-dimensional Hall-algebra questions on $A = \End(T)^{\mathrm{op}}$; the
superpotential $W$ reconstructing $A$ as a Jacobi algebra is uniquely
determined (up to right-equivalence) from the cyclic $A_\infty$-structure on
the tilting algebra (Bocklandt 2008, arXiv:\texttt{math/0603558} Thm.~3.1).

\subsection*{Existence and failure}

The Van den Bergh theorem (Theorem~\ref{thm:vdb-equivalence}) guarantees the
existence of a tilting bundle in precisely the \emph{local three-dimensional}
crepant setting: $X$ obtained as a crepant resolution of an affine Gorenstein
singularity $\Spec R$ with $\dim R \leq 3$. This covers toric CY$_3$,
orbifolds $\C^n/G$ with $G \subset SU(n)$ of dimension $n \leq 3$,
$(-3)$-curve contractions, and their relatives. Beyond this scope, global
tilting \emph{fails}. Three cases are central:
\begin{itemize}
\item[(a)] \emph{Projective K3 surface $S$ with generic polarisation.} The
obstruction is $\Ext^{2}_{\cO_S}(\cF, \cF) \neq 0$ for any locally-free $\cF$:
Serre duality on a CY$_2$ identifies
$\Ext^2(\cF, \cF) \cong \Hom(\cF, \cF)^{\vee} \neq 0$. The higher-Ext vanishing
hypothesis (H2) of Theorem~\ref{thm:derived-morita-tilting} cannot hold. The
K3 category admits compact generation (by Bondal--Van den Bergh, Thm.~3.1.1)
but no tilting generator.
\item[(b)] \emph{Quintic threefold $X \subset \bP^4$.} The obstruction is the
non-triviality of $H^{0,3}(X) = \C$ (the holomorphic volume form), which
combined with Serre duality forces $\Ext^3_{\cO_X}(\cF, \cF) \neq 0$ for
non-zero $\cF$. No $\Ext^{>0}$-vanishing tilting generator exists globally.
\item[(c)] \emph{Compact hyperk\"ahler manifold of dimension $\geq 4$.}
Hyperk\"ahler symmetry forces $h^{0,p}(X) = 1$ for all $p$ even, $0 \leq p \leq
\dim X$; the higher $\Ext$-vanishing hypothesis is obstructed at every even
degree $\leq \dim X$.
\end{itemize}
The general pattern: at $d \geq 3$, if $H^{0,q}(X) \neq 0$ for some $1 \leq q
\leq d$, then $\Ext^q(T, T) \supseteq H^{0,q}(X) \neq 0$ for any
locally-free $T$ with $\cO_X$ as a summand; the hypothesis (H2) of
Theorem~\ref{thm:derived-morita-tilting} fails. In such cases, no global
tilting exists; only local tilting on Zariski patches, and the gluing of
Parts I--III (or Part V via factorisation algebras) applies.

% =======================================================================
\section{Tilting as gluing presentation}
\label{sec:tilting-as-gluing}
% =======================================================================

The tilting bundle is a \emph{presentation} of $X$ by a single
quiver-with-potential: the quiver vertices are the summands of
$T = \bigoplus_i T_i$, the arrows are extensions among them, and the potential
encodes the exact $3$-Calabi--Yau constraint on the Ginzburg dg-lift
$\Gamma(Q, W)$.

\subsection*{The decomposition $T = \bigoplus_i T_i$}

Write $T = \bigoplus_{i \in I} T_i$ as a finite direct sum of indecomposable
summands. This decomposition is canonical up to isomorphism of summands
(Krull--Schmidt on the additive category generated by $T$). Then
\[
  A \;=\; \End(T)^{\mathrm{op}} \;=\; \bigoplus_{i, j \in I}
  \Hom_{\cO_X}(T_i, T_j)^{\mathrm{op}}
\]
is a matrix algebra over the ``diagonal'' algebras $\End(T_i)$, indexed by
pairs $(i,j)$. The off-diagonal block $\Hom_{\cO_X}(T_i, T_j)$ is the
\emph{arrow space} from $i$ to $j$: its basis is the space of morphisms
$T_i \to T_j$ in $\Coh(X)$ modulo scaling.

In the conifold example $T = \cO \oplus \cO(1)$:
\[
  \Hom(\cO, \cO(1)) \;=\; H^{0}(\bP^1, \cO(1)) \oplus H^{0}(\bP^1, \cO(-1))
  \;=\; \C^{2} \oplus 0 \;=\; \C a_1 \oplus \C a_2,
\]
and $\Hom(\cO(1), \cO) = H^0(\cO(-1)) \oplus H^0(\cO(1)) = 0 \oplus \C^{2}
= \C b_1 \oplus \C b_2$ (the second term arises via the $\cO(-1) \oplus
\cO(-1)$ fibre direction). Similarly $\End(\cO) = \End(\cO(1)) = \C$. The
Klebanov--Witten quiver with vertices $\{\circ, \bullet\}$, arrows $a_1, a_2
\colon \circ \to \bullet$ and $b_1, b_2 \colon \bullet \to \circ$, and
superpotential $W = a_1 b_1 a_2 b_2 - a_1 b_2 a_2 b_1$ emerges \emph{directly}
from the cohomology count.

\subsection*{The Jacobi algebra, the potential, and the Ginzburg dg-lift}

The endomorphism algebra of a tilting bundle on a smooth local CY$_3$ carries
a cyclic derivation structure: there exists a superpotential
$W \in \C Q_{\mathrm{cyc}} = \C Q / [\C Q, \C Q]$ in the cyclic quotient of
the path algebra such that the relations among the arrows are precisely the
cyclic derivatives $\partial_a W = 0$. The pair $(Q, W)$ is the quiver with
potential; the \emph{Jacobi algebra} is
\[
  J = \Jac(Q, W) \;=\; \C Q / \langle \partial_a W \colon a \in Q_1 \rangle.
\]
The CY$_3$ property is carried not by $J$ but by its \emph{Ginzburg dg-lift}.

\begin{definition}[Ginzburg dg-algebra; Ginzburg 2006,
arXiv:\texttt{math/0612139} \S5]
\label{def:ginzburg-dg}
\ClaimStatusDefinition\ Given a quiver with potential $(Q, W)$, the
\emph{Ginzburg dg-algebra} $\Gamma(Q, W)$ is the tensor algebra on the graded
quiver
\begin{itemize}
\item arrows of $Q$ in degree $0$;
\item one reverse arrow $a^{*} \colon j \to i$ in degree $-1$ for each
$a \colon i \to j$ in $Q_1$;
\item one loop $t_i \colon i \to i$ in degree $-2$ at each vertex $i \in I$,
\end{itemize}
equipped with the differential
$d a = 0$, $d a^{*} = \partial_a W$, and $d t_i = \sum_{a \in Q_1}
[a, a^{*}] \cdot e_i$ (the total commutator at vertex $i$).
\end{definition}

\begin{proposition}[CY-3 property lives on $\Gamma$, not on $J$;
Ginzburg 2006, Thm.~3.2.8 and \S5.3]
\label{prop:cy3-on-gamma-not-j}
\ClaimStatusProposition\ For $(Q, W)$ a quiver-with-potential arising from a
local CY$_3$:
\begin{enumerate}
\item[(a)] The dg-algebra $\Gamma(Q, W)$ is \emph{exact $3$-Calabi--Yau}: its
inverse dualising complex satisfies
$\RHom_{\Gamma^e}(\Gamma, \Gamma \otimes \Gamma)[3] \simeq \Gamma$ as
$\Gamma$-bimodules, with the CY$_3$ shift degree $+3$ matching the PTVV shift
$(d, \mathrm{shift}, E_n^{\mathrm{cl}}) = (3, -1, E_1)$ under dg-Hochschild
Serre duality.
\item[(b)] The Jacobi algebra $J = H^0(\Gamma(Q, W))$ is the degree-zero
cohomology of $\Gamma$; it is \emph{not} by itself a CY$_3$ algebra. Calling
$J$ ``CY$_3$'' is a category error: $J$ has no intrinsic shift structure. The
three-dimensional duality is an assertion about $\Gamma$-bimodules.
\item[(c)] For $Q$ connected and $W$ suitably non-degenerate,
$H^{<0}(\Gamma(Q, W)) = 0$, so $\Gamma \simeq J$ in the derived category of
$\Gamma$ (hence derived Morita between $\Gamma$ and $J$), and only via this
quasi-isomorphism does the CY$_3$ structure on $\Gamma$ descend to the
derived categorical structure $D^b(J\text{-}\mod)$.
\end{enumerate}
\end{proposition}

\begin{theorem}[Tilting-via-quiver presentation; Ginzburg 2006, Keller 2011
arXiv:\texttt{1102.4148}]
\label{thm:tilting-jacobi}
\ClaimStatusTheorem\ Let $T = \bigoplus_{i \in I} T_i$ be a tilting bundle on a
smooth local CY$_3$ $X$. Let $Q$ be the quiver with vertex set $I$ and arrow
spaces $Q_1(i,j) = \Ext^{1}_{\cO_X}(T_i, T_j)^{*}$. There exists a
quasi-homogeneous superpotential $W \in \C Q_{\mathrm{cyc}}$, unique up to
gauge, such that
\[
  \Jac(Q, W) \;=\; \End_{\cO_X}(T)^{\mathrm{op}},
  \qquad
  \Gamma(Q, W) \;\simeq\; \RHom_{\cO_X}(T, T) \text{ as dg-algebras},
\]
and the CY$_3$ structure on $\Perf(X)$ transports (via
Theorem~\ref{thm:derived-morita-tilting} applied to $\Gamma$) to the exact
CY$_3$ structure on $\Gamma$ of Proposition~\ref{prop:cy3-on-gamma-not-j}(a).
\end{theorem}

\begin{proof}[Proof sketch.]
The arrow spaces are Ext-spaces: under (H2) the underived and derived
$\Hom(T_i, T_j)$ agree in degree zero, so the first non-trivial term for
$i \neq j$ is $\Ext^{1}(T_i, T_j)$, which Serre-dualises to
$\Ext^{2}(T_j, T_i)$ on the CY$_3$; the potential encodes the $m_3$ Massey
product. The construction of $W$ is the Kontsevich--Costello
reconstruction of a cyclic $A_\infty$-algebra from its degree-zero piece plus
Serre duality: the $m_3$ on $\Ext^1$ evaluated against Serre duality
reproduces the degree-3 piece of $W$, higher $m_k$ give higher-degree
corrections, and quasi-homogeneity follows from the internal CY$_3$ grading
on $\Gamma$. The quasi-isomorphism $\RHom(T, T) \simeq \Gamma(Q, W)$ is
Keller's derived-equivalence theorem (Keller 2011 \S A.15). The CY$_3$
property, as Proposition~\ref{prop:cy3-on-gamma-not-j} records, is a property
of $\Gamma$ as a dg-algebra; the statement ``$J$ is CY$_3$'' is shorthand for
``$\Gamma$ is CY$_3$ and $J = H^0(\Gamma)$''.
\end{proof}

\subsection*{Canonical examples}

\paragraph{Conifold $\tilde Y = \cO_{\bP^1}(-1)^{\oplus 2}$.} Tilting
$T = \cO \oplus \cO(1)$. Quiver: two vertices, two arrows each way. Potential
$W = a_1 b_1 a_2 b_2 - a_1 b_2 a_2 b_1$ (Klebanov--Witten, Szendr\H{o}i 2008).

\paragraph{Local $\bP^2 = \Tot(\cO_{\bP^2}(-3))$.} Tilting
$T = \cO \oplus \cO(1) \oplus \cO(2)$ (Beilinson 1978). Quiver: three
vertices, $3 \times 3 = 9$ arrows grouped as three Koszul strands (Beilinson
triple $\cO, \cO(1), \cO(2)$). Potential
$W_{P^2} = \sum_{\sigma \in S_3} \mathrm{sgn}(\sigma)\,
x_{\sigma(1)} y_{\sigma(2)} z_{\sigma(3)}$: the antisymmetric McKay cubic
(Beasley--Plesser 2001).

\paragraph{$(-3)$-curves.} Let $\tilde Y$ contain a rational curve $C \cong
\bP^1$ with normal bundle $N_{C/\tilde Y} = \cO(-1) \oplus \cO(-1)$ that can
be contracted to a singular point of $X$. Craw 2008 (building on the
Craw--Ishii 2004 analysis of $\bC^3/G$-Hilb schemes) constructs tilting
$T_C = \cO_{\tilde Y} \oplus \cO_{\tilde Y}(C)$ (or an explicit line-bundle
refinement with intersection numbers read off from $N_{C/\tilde Y}$). The
associated tilting algebra is derived-equivalent to the quiver from the
contracted singularity; tilting automorphisms realise the flop $\tilde Y
\dashrightarrow \tilde Y^+$ as derived Morita on the same CY$_3$ category.

\paragraph{$\C^3/\Z_n$ orbifolds.} Tilting
$T = \bigoplus_{k=0}^{n-1} \cO \otimes \chi_k$ on the resolved orbifold, where
$\chi_k$ are the characters of $\Z_n$. Quiver: $n$ vertices arranged on a
cycle; arrows per McKay graph; potential cubic (type A McKay).

\subsection*{Quiver-with-potential as universal gluing data}

The force of Theorem~\ref{thm:tilting-jacobi} is that for every tilting
presentation, a single pair $(Q, W)$ plays the role of the entire \v{C}ech
atlas plus all its transition functors. One algebra $A = \Jac(Q, W)$ encodes
everything. In particular:
\begin{itemize}
\item The \v{C}ech nerve $\{U_\alpha\} \cup \{U_\alpha \cap U_\beta\} \cup \ldots$
  on $X$ collapses to the vertex set $I$ and the arrow set $Q_1$;
\item The transition functors between charts collapse to the relations
  $\partial_a W = 0$;
\item The higher cocycle data on triple overlaps collapses to the higher
  $A_\infty$ structure on $\Gamma(Q, W)$.
\end{itemize}
The tilting presentation is therefore the \emph{economy} of the gluing: one
single finite-dimensional algebra replaces the infinite-dimensional \v{C}ech
diagram.

% =======================================================================
\section{Transport of Hall product along derived Morita}
\label{sec:hall-transport}
% =======================================================================

The CoHA on $X$ is defined via the moduli stack of coherent sheaves
$\cM(X)$ with its critical cohomology ring structure (Kontsevich--Soibelman
2008, Davison--Meinhardt 2020 \texttt{arXiv:1601.02479}). On the NCCR side, the
same construction yields the CoHA on the moduli stack $\cM(A)$ of
finite-dimensional $A$-modules. The derived Morita equivalence lifts to these
moduli stacks, and Davison--Hennecart--Schlegel Mejia (DHSM) 2023
\texttt{arXiv:2303.12592} prove that the Hall product transports. The
assertion of this section is the precise form of that transport and its
first-principles derivation.

\subsection*{Moduli-stack lift of derived Morita}

\begin{theorem}[Moduli-stack Morita; Davison--Hennecart--Schlegel Mejia 2023,
arXiv:\texttt{2303.12592} Thm.~A]
\label{thm:moduli-morita}
\ClaimStatusTheorem\ Let $X$ be a smooth local CY$_3$ admitting a tilting
bundle $T$ satisfying hypotheses (H1)--(H2) of
Theorem~\ref{thm:derived-morita-tilting}, and let $A = \End(T)^{\mathrm{op}} = \Jac(Q, W)$
be the associated Jacobi algebra. Then the adjoint pair $L \dashv R$ descends
to an equivalence of (derived) moduli stacks of semistable objects
\[
  \cM^{\mathrm{ss}}(X) \;\simeq\; \cM^{\mathrm{ss}}(A)
\]
under the class-level identification
\[
  \mathrm{cl}_T \;:\; K_0(\Coh X) \;\xrightarrow{\;\sim\;}\; \Z^{I}, \qquad
  [\cF] \;\longmapsto\; (\dim \RHom(T_i, \cF))_{i \in I}.
\]
The lift intertwines the universal family $\cE_X$ on $\cM(X)$ with the
universal representation $V_A$ on $\cM(A)$: $\RHom(T, \cE_X) = V_A$ as
coherent sheaves on the common moduli stack.
\end{theorem}

\begin{proof}[Sketch.]
Derived Morita transports the abelian heart (a $t$-structure on $D^b(\Coh X)$
whose heart is the image of $A$-$\mod$) to the standard heart on $A$-$\mod$.
Bridgeland's identification of stability conditions preserves semistability
under this $t$-structure alignment. The resulting map on moduli is an
equivalence of stacks because the universal family $\cE_X$ on $\cM(X)$ maps to
the universal representation $V_A$ on $\cM(A)$, and conversely, via the
inverse functor $T \otimes^L_A -$.
\end{proof}

\subsection*{Critical cohomology transports}

The CoHA on either side is
\[
  \cH(X) \;=\; H^{\bullet,\mathrm{BM}}_{T_{\C^\times}}(\cM(X), \vphi_{\mathrm{tr}\,W_X}),
\]
where $\vphi$ is the vanishing-cycle functor and $W_X$ is the canonical CY$_3$
superpotential (pulled back from the framing). On the NCCR side, the
superpotential is precisely $W = W_{Q,W}$ --- the potential of the
quiver-with-potential --- and $\vphi_{\tr W}$ is the vanishing-cycle sheaf on
the representation variety $\mathrm{Rep}(Q, \underline{d})$.

\begin{theorem}[Vanishing-cycle transport; DHSM 2023, extending Davison 2017]
\label{thm:vanishing-cycle-transport}
\ClaimStatusTheorem\ Under the moduli equivalence of
Theorem~\ref{thm:moduli-morita}, the vanishing-cycle sheaf
$\vphi_{\tr W_X}$ on $\cM(X)$ is canonically isomorphic to
$\vphi_{\tr W_Q}$ on $\cM(A)$, compatibly with the $T$-equivariant structure.
Consequently:
\[
  \cH(X) \;\simeq\; \cH(A) \quad\text{as equivariant Borel--Moore
  cohomology modules over the respective moduli stacks.}
\]
\end{theorem}

\begin{proof}[First-principles derivation.]
A coherent sheaf $\cF$ on $X$ corresponds, under $\Phi_T$, to an $A$-module
$M = \RHom(T, \cF)$. Both $\cF$ and $M$ admit \emph{the same} deformation
theory: the obstruction space is $\Ext^{2}(\cF, \cF) \cong \Ext^{2}_A(M, M)$
(Serre duality on $X$ matches the cyclic duality on $A$). The superpotential
$W_X$ on the deformation space matches $\tr W_Q$ pulled back along the tilting
isomorphism, because
\[
  W_X|_{T^{*}_{[\cF]} \cM(X)} \;=\;
  \text{degree-3 part of the $A_\infty$ structure on $\Ext^\bullet(\cF, \cF)$},
\]
and Theorem~\ref{thm:tilting-jacobi} identifies precisely this degree-3 piece
with the Jacobi potential $W$. The vanishing-cycle sheaf is therefore an
invariant of the pair (deformation space, superpotential), transporting
identically along the moduli equivalence.
\end{proof}

\subsection*{Extension stacks and Hall product transport}

The Hall product on $\cH(X)$ is built from the \emph{extension stack}
$\mathrm{Ext}(\cM \times \cM, \cM)$: for charge classes
$\gamma_1, \gamma_2 \in K_0(\Coh X)$, the extension stack parametrises short
exact sequences $0 \to \cF_1 \to \cF \to \cF_2 \to 0$ with
$[\cF_i] = \gamma_i$ and $[\cF] = \gamma_1 + \gamma_2$. The Hall product is
convolution along the push/pull diagram
\[
  \cM_{\gamma_1} \times \cM_{\gamma_2}
  \;\xleftarrow{\;\mathrm{ends}\;}\;
  \mathrm{Ext}_{\gamma_1, \gamma_2}
  \;\xrightarrow{\;\mathrm{middle}\;}\;
  \cM_{\gamma_1 + \gamma_2}.
\]

\begin{proposition}[Extension-stack transport]
\label{prop:extension-stack-transport}
\ClaimStatusProposition\ Under the moduli equivalence
$\cM(X) \simeq \cM(A)$ the extension stacks satisfy
\[
  \mathrm{Ext}_{\gamma_1, \gamma_2}^{X} \;\simeq\;
  \mathrm{Ext}_{\mathrm{cl}_T(\gamma_1), \mathrm{cl}_T(\gamma_2)}^{A},
\]
with the ``ends'' and ``middle'' morphisms intertwined by $\Phi_T$. Combined
with Theorem~\ref{thm:vanishing-cycle-transport}, this yields:
\end{proposition}

\begin{theorem}[Hall-product transport along derived Morita; DHSM 2023
arXiv:\texttt{2303.12592} Thm.~A, with super-grading extension Thm.~B]
\label{thm:hall-product-transport}
\ClaimStatusTheorem\ Under hypotheses (H1)--(H2) on the tilting bundle $T$,
there is an isomorphism of (super-graded) bialgebras
\[
  \star_T \;:\; \cH(X) \;\xrightarrow{\;\sim\;}\; \cH(A)
\]
intertwining the Hall product, the natural coproduct (where defined), and the
$T_{\C^\times}$-equivariant structures. On charges, $\star_T$ is the
$\Z$-linear extension of $\mathrm{cl}_T$. The \emph{$\Z/2$-super grading}
(cohomological parity on $\vphi_{\tr W}$) transports along the same
equivalence, by DHSM Thm.~B.
\end{theorem}

\begin{proof}[Proof.]
The Hall product is a composition of pullback-along-ends, intersection with
the vanishing-cycle sheaf, and pushforward-along-middle. Each factor transports
under derived Morita: the ends and middle morphisms transport by
Proposition~\ref{prop:extension-stack-transport}; the vanishing-cycle sheaf
transports by Theorem~\ref{thm:vanishing-cycle-transport}; the equivariant
structure transports because the tilting bundle is $T_{\C^\times}$-equivariant
in the toric setting (and equivariantly localisable at the fixed locus in the
general local CY$_3$ setting). The $\Z/2$-super grading transports because
$\vphi_{\tr W_X}$ and $\vphi_{\tr W_Q}$ carry identically-graded sheaves of
nearby cycles under the cohomological parity; DHSM Thm.~B upgrades Thm.~A from
an isomorphism of bialgebras to an isomorphism of super-bialgebras.
\end{proof}

Theorem~\ref{thm:hall-product-transport} is established unconditionally by
DHSM (2023, Thm.~A + Thm.~B) for $A = \Jac(Q, W)$ the Jacobi algebra of any
quiver-with-potential; the CY$_3$ hypothesis is absorbed into the assertion
that the Ginzburg dg-lift $\Gamma(Q, W)$ is exact $3$-Calabi--Yau
(Proposition~\ref{prop:cy3-on-gamma-not-j}). For a smooth local CY$_3$ $X$
with tilting presentation $(Q, W)$, this yields the full
$\cH(X) \simeq \cH(\Jac(Q, W))$ identification on $\C^3$, the conifold
(Klebanov--Witten $(Q, W)$), local $\bP^2$ (Beilinson cubic), and
$\C^3/\Z_n$ for all $n$.

\subsection*{Why Hall product preservation is forced}

An alternative way to see the transport is purely representation-theoretic:
both $\cH(X)$ and $\cH(A)$ are instances of the \emph{universal positive-
geometry grammar} $Y^+(X) = H^\bullet_{\mathrm{eq}}(\cM^+_{\mathrm{eff}}(X),
\vphi_W)$ (CLAUDE.md, ``Universal positive-geometry grammar''). The
grammar is invariant under any derived equivalence that preserves the
effective semistable moduli stack and the superpotential. Derived Morita
preserves both, by construction. Hence the CoHAs agree, and the Hall product
structure is preserved.

% =======================================================================
\section{Gluing-preservation theorem for mutations}
\label{sec:mutation-preservation}
% =======================================================================

If $T$ and $T'$ are two tilting bundles on the same CY$_3$ $X$, they give two
tilting algebras $A = \End(T)$ and $A' = \End(T')$, and by
Theorem~\ref{thm:hall-product-transport} the same CoHA $\cH(X)$: both
transports land on the same Hall-product object. This induces an isomorphism
$\cH(A) \simeq \cH(A')$. When $T$ and $T'$ are related by a \emph{mutation}
(Derksen--Weyman--Zelevinsky 2008), the induced isomorphism has an explicit
combinatorial presentation.

\subsection*{Mutation at a vertex}

\begin{definition}[Mutation of quiver-with-potential; DWZ 2008,
arXiv:\texttt{0704.0649} Defs.~5.1, 5.5, 7.1]
\label{def:mutation-qp}
\ClaimStatusDefinition\ Let $(Q, W)$ be a reduced quiver-with-potential
(no $2$-cycles in $Q$, no $2$-cycle terms in $W$) with vertex set $I$, and
fix a vertex $k \in I$ supporting no loops. The \emph{mutation at $k$},
$\mu_k(Q, W) = (Q', W')$, factors as a composite of two moves:
the \emph{premutation} $\tilde\mu_k$ followed by the \emph{reduction}
$\mathrm{red}$ of DWZ Theorem~4.6,
\[
  \mu_k(Q, W) \;=\; \mathrm{red}\bigl(\tilde\mu_k(Q, W)\bigr).
\]
The premutation $\tilde\mu_k(Q, W) = (\tilde Q, \tilde W)$ is the three-step
arrow surgery:
\begin{enumerate}
\item[(M1)] \emph{Compose through $k$.} For each pair of arrows
$a \colon i \to k$ and $b \colon k \to j$ with $i, j \neq k$, adjoin a
formal composite $[ba] \colon i \to j$ to the arrow set.
\item[(M2)] \emph{Reverse $k$-incident arrows.} Replace each
$a \colon i \to k$ by $a^{*} \colon k \to i$, and each $b \colon k \to j$
by $b^{*} \colon j \to k$; arrows not incident to $k$ are unchanged.
\item[(M3)] \emph{Update the potential.} Set
\[
  \tilde W \;=\; [W]_{\{[ba]\}} \;+\; \sum_{\substack{a \colon i \to k \\ b \colon k \to j}}
  [ba]\, b^{*}\, a^{*},
\]
where $[W]_{\{[ba]\}}$ is the cyclic representative of $W$ obtained by
replacing every length-two path $ba$ through $k$ by the formal symbol $[ba]$.
\end{enumerate}
The reduction splits the premutation: DWZ Theorem~4.6 provides a canonical
right-equivalence decomposition
$(\tilde Q, \tilde W) \simeq_{\mathrm{r.e.}} (Q_{\mathrm{red}},
W_{\mathrm{red}}) \oplus (Q_{\mathrm{triv}}, W_{\mathrm{triv}})$
into a \emph{reduced} summand (no $2$-cycle terms in $W_{\mathrm{red}}$) and
a \emph{trivial} summand (Jacobi algebra zero), unique up to right-equivalence.
Set $(Q', W') = (Q_{\mathrm{red}}, W_{\mathrm{red}})$.
\end{definition}

The reduction is not the physicist's ``Euler--Lagrange elimination'' but the
Derksen--Weyman--Zelevinsky right-equivalence: a $2$-cycle term
$c\, x\, y$ in $\tilde W$ is removed by the change of variable
$x \mapsto x - c^{-1}\,\partial_{y}(\text{rest of }\tilde W)$, which DWZ
prove is an automorphism of the completed path algebra preserving the cyclic
class of the potential (DWZ Thms.~4.2, 4.6).

The \emph{non-degeneracy hypothesis} on $(Q, W)$ is that every iterated
premutation admits a reduction. DWZ \S7 prove that non-degenerate QPs are
generic and that all QPs arising from local CY$_3$ tilting --- the conifold,
local $\bP^2$, $\C^3/\Z_n$ --- are non-degenerate (DWZ Cor.~7.4).

\begin{proposition}[Mutation is an involution in reduced QPs; DWZ 2008
Thm.~5.7]
\label{prop:mutation-involution}
\ClaimStatusProposition\ In the category whose objects are reduced QPs and
whose morphisms are right-equivalences, $\mu_k \circ \mu_k$ is naturally
right-equivalent to the identity: $\mu_k(\mu_k(Q, W)) \simeq_{\mathrm{r.e.}}
(Q, W)$. The involution holds \emph{up to right-equivalence}, not strictly:
the double premutation $\tilde\mu_k^2$ differs from the identity by a
trivial $2$-cycle summand that cancels only after reduction.
\end{proposition}

Geometrically, $T' = \mu_k(T)$ is \emph{not} a plain ``dual-and-shift''
$T_k^{\vee}[1]$ but a cone on the $\Ext^1$-arrow data at $k$. Replace the
summand $T_k$ by one of two mutation cones,
\[
  T_k^{+} \;=\; \mathrm{cone}\!\left(T_k \;\xrightarrow{\;\mathrm{ev}\;}\;
  \bigoplus_{b \colon k \to j} T_j^{\oplus \dim b}\right),
  \quad
  T_k^{-} \;=\; \mathrm{cone}\!\left(\bigoplus_{a \colon i \to k}
  T_i^{\oplus \dim a} \;\xrightarrow{\;\mathrm{coev}\;}\; T_k\right)[-1],
\]
the first built from arrows \emph{out of} $k$ and the second from arrows
\emph{into} $k$; the two differ by a global shift (Iyama--Reiten 2008,
arXiv:\texttt{0707.1669} \S5; Keller--Yang 2011, arXiv:\texttt{0906.0761}
\S3). The involution Proposition~\ref{prop:mutation-involution} corresponds
geometrically to $\mu_k^{-}(\mu_k^{+}(T)) \simeq T$, which is a
Serre-duality statement on the two cones, \emph{not} to the spurious
$T_k^{\vee}[2] = T_k$.

\begin{theorem}[Mutation gives derived equivalence; Keller--Yang 2011,
arXiv:\texttt{0906.0761} Thms.~3.2 and 6.3]
\label{thm:keller-yang-mutation}
\ClaimStatusTheorem\ Assume $(Q, W)$ is non-degenerate. The mutation
$\mu_k \colon (Q, W) \mapsto (Q', W') = \mu_k(Q, W)$ lifts to a pair of
triangulated equivalences
\[
  \Mu_k^{+} \;:\; \cD(\Gamma(Q, W)) \;\xrightarrow{\;\sim\;}\;
  \cD(\Gamma(Q', W')),
  \qquad
  \Mu_k^{-} \;:\; \cD(\Gamma(Q, W)) \;\xrightarrow{\;\sim\;}\;
  \cD(\Gamma(Q', W')),
\]
where $\cD(\Gamma) = D(\Gamma\text{-}\mathrm{dg\text{-}mod})$. The two equivalences
are constructed from two distinct tilting objects $T_k^{+}, T_k^{-}$ in
$\cD(\Gamma(Q, W))$, built from the decomposition $\Gamma(Q, W) =
\bigoplus_{i \in I} P_i$ into indecomposable projective dg-modules:
\[
  T_k^{+} \;=\; \mathrm{cone}\Bigl(P_k \;\xrightarrow{\;(b)_{b \colon k \to j}\;}\;
  \bigoplus_{b \colon k \to j} P_{t(b)}\Bigr) \;\oplus\; \bigoplus_{i \neq k} P_i,
\]
with the differential given by the arrows $b$ \emph{out of} $k$ in $Q$ (the
sum ranges over $b \in Q_1$ with source $k$ and target $t(b) = j$), and
similarly $T_k^{-}$ built from arrows \emph{into} $k$. Keller--Yang Thm.~3.2
establishes that $T_k^{\pm}$ is a compact generator of $\cD(\Gamma(Q, W))$
with $\Ext^{>0}(T_k^{\pm}, T_k^{\pm}) = 0$; Keller--Yang Thm.~6.3 identifies
$\End_{\cD(\Gamma(Q, W))}(T_k^{\pm}) \simeq \Gamma(Q', W')$ as dg-algebras,
and then Theorem~\ref{thm:derived-morita-tilting} yields the equivalences
$\Mu_k^{\pm}$.

Along $\Mu_k^{\pm}$, the simple $S_k$ (the simple at vertex $k$) over
$\Gamma(Q, W)$ maps to $S_k^{\pm}[\mp 1]$ over $\Gamma(Q', W')$: the
\emph{shift direction is opposite} for $\Mu_k^{+}$ and $\Mu_k^{-}$, and the
two equivalences differ by the spherical twist at $S_k$ (Keller--Yang Prop.~6.4).
\end{theorem}

\subsection*{Induced CoHA automorphism}

\begin{theorem}[Mutation-induced CoHA automorphism]
\label{thm:mutation-coha-automorphism}
\ClaimStatusTheorem\ Let $(Q, W)$ be a non-degenerate quiver-with-potential
whose Ginzburg dg-lift $\Gamma(Q, W)$ is exact $3$-Calabi--Yau, with Jacobi
algebra $A = \Jac(Q, W) = H^0(\Gamma(Q, W))$, and let $\mu_k$ denote mutation
at vertex $k$. Let $\Mu_k^{+} \colon \cD(\Gamma(Q, W)) \to
\cD(\Gamma(\mu_k(Q, W)))$ be the positive Keller--Yang equivalence of
Theorem~\ref{thm:keller-yang-mutation}. Transporting Hall products via
Theorem~\ref{thm:hall-product-transport}:
\[
  \cH(A) \;\xrightarrow{\;\Mu_{k,\star}^{+}\;}\; \cH(A')
\]
is an isomorphism of super-bialgebras. If $(Q, W)$ and $(Q', W') =
\mu_k(Q, W)$ arise as tilting presentations of the \emph{same} geometric
CY$_3$ $X$ --- equivalently, $\mu_k$ relates two tilting bundles $T, T'$ on
$X$ --- then $\Mu_{k,\star}^{+}$ is an automorphism of the single CoHA
$\cH(X)$. The opposite mutation $\Mu_k^{-}$ gives the inverse automorphism
$(\Mu_{k,\star}^{+})^{-1}$.
\end{theorem}

\begin{proof}[Proof.]
By Theorem~\ref{thm:keller-yang-mutation}, $\Mu_k^{+}$ is a triangulated
equivalence between the derived categories of the two CY$_3$ Ginzburg
algebras, constructed as a derived Morita equivalence along the tilting
object $T_k^{+}$. By Theorem~\ref{thm:hall-product-transport}, any derived
Morita equivalence between CY$_3$ Ginzburg algebras induces a super-bialgebra
isomorphism on the critical CoHAs, with equivariant and $\Z/2$-supergrading
structure (DHSM 2023 Thm.~B). When $(Q, W)$ and $(Q', W')$ arise as two
tilting presentations of the same $X$, the corresponding tilting objects
$T, T'$ give two factorisations of the identity
$\mathrm{id}_{\cH(X)} = (\star_{T'})^{-1} \circ \Mu_{k,\star}^{+} \circ \star_T$,
so $\Mu_{k,\star}^{+}$ is an endomorphism of $\cH(X)$. Non-degeneracy of
$(Q, W)$ (DWZ Cor.~7.4) ensures $\Mu_k^{-}$ exists and gives the inverse
at the derived level, hence $\Mu_{k,\star}^{+}$ is an automorphism.
\end{proof}

\subsection*{Geometric interpretation: Seiberg duality}

In the physics literature, mutation at a vertex is \emph{Seiberg duality}: a
quiver gauge theory undergoes an exact IR equivalence under a node mutation
with appropriate matter content adjustment (Seiberg 1994; Berenstein--Douglas
2002). The present theorem is the derived-category / CoHA shadow of this
physical statement: the non-perturbative BPS algebra is invariant under
Seiberg duality.

\begin{example}[Conifold flop via mutation]
\label{ex:conifold-seiberg-duality}
\ClaimStatusDerivation\ The Klebanov--Witten quiver
$(Q_{\mathrm{con}}, W_{\mathrm{con}})$ has two vertices
$I = \{0, 1\}$, four arrows $a_1, a_2 \colon 0 \to 1$ and $b_1, b_2 \colon 1
\to 0$, and the cyclic quartic potential
\[
  W_{\mathrm{con}} \;=\; a_1\, b_1\, a_2\, b_2 \;-\; a_1\, b_2\, a_2\, b_1 \,.
\]
Premutation $\tilde\mu_0$ (DWZ recipe applied at vertex $k = 0$). Step
(M1) requires pairs of arrows $a \colon i \to k$ and $b \colon k \to j$
--- i.e., arrows into $0$ composed with arrows out of $0$. The arrows into
$0$ are $b_1, b_2 \colon 1 \to 0$, and the arrows out of $0$ are $a_1,
a_2 \colon 0 \to 1$; the composites $a \cdot b$ through $0$ run from
$1$ to $1$, adjoined as formal arrows $[a_i b_j] \colon 1 \to 1$ for
$i, j \in \{1, 2\}$ --- four new arrows.
Step (M2) reverses the $0$-incident arrows: $a_i \colon 0 \to 1$ becomes
$a_i^* \colon 1 \to 0$, and $b_j \colon 1 \to 0$ becomes $b_j^* \colon 0
\to 1$. Step (M3) gives
\[
  \tilde W \;=\; [a_1 b_1]\,[a_2 b_2] \;-\; [a_1 b_2]\,[a_2 b_1]
  \;+\; \sum_{i, j}\, [a_i b_j]\, a_i^{*}\, b_j^{*},
\]
where the first two terms are the original $W_{\mathrm{con}}$ rewritten
in the composite variables, and the last sum is the universal $[ba]\, b^*\,
a^*$ piece. The premutation is a cyclic quartic-plus-cubic.

Reduction: the premutation has no $2$-cycles at vertex $0$ or $1$ (the
paths $a_i^* b_j^*$ run $1 \to 0 \to 1$ so sit in the $(1, 1)$-block, but
in pairs $a_i^*, b_j^*$ which are single arrows, not length-two cycles
within a single vertex). However, the cubic terms $[a_i b_j] a_i^* b_j^*$
together with the quartic pair up under the right-equivalence of DWZ \S4.
Explicitly, setting $c_{ij} = [a_i b_j]$, the cyclic derivatives
$\partial_{c_{ij}} \tilde W = \partial_{c_{ij}}\!\bigl(c_{11} c_{22} -
c_{12} c_{21} + \sum_{k\ell} c_{k\ell}\, a_k^* b_\ell^*\bigr)$ solve for
the $c_{ij}$ in terms of $a_i^* b_j^*$, and the reduced potential is
\[
  W' \;=\; a_1^*\, b_1^*\, a_2^*\, b_2^* \;-\; a_1^*\, b_2^*\, a_2^*\, b_1^*,
\]
the same Klebanov--Witten quartic with roles $a \leftrightarrow a^*$ and
$b \leftrightarrow b^*$ swapped. The result $\mu_0(Q_{\mathrm{con}},
W_{\mathrm{con}})$ is right-equivalent to $(Q_{\mathrm{con}}, W_{\mathrm{con}})$
with the vertex labels $\{0, 1\}$ swapped --- the $\Z/2$ symmetry
interchanging the two chambers of the Klebanov--Witten quiver.

The derived equivalence $\Mu_0^{+} \colon \cD(\Gamma(Q_{\mathrm{con}},
W_{\mathrm{con}})) \xrightarrow{\sim} \cD(\Gamma(\mu_0(Q_{\mathrm{con}},
W_{\mathrm{con}})))$ is exactly the Bridgeland--Chen derived equivalence
$D^b(\tilde Y) \simeq D^b(\tilde Y^{+})$ between the two small resolutions
of the conifold singularity $xy = zw$ (Bridgeland 2002, arXiv:\texttt{math/0009053};
Chen 2002). On the Hall-algebra side, $\Mu_{0,\star}^{+} \colon
\cH(\tilde Y) \xrightarrow{\sim} \cH(\tilde Y^{+})$ transports the CoHA
structure; under the identification $\tilde Y \simeq \tilde Y^{+}$ as
abstract varieties (the two small resolutions of the same conifold are
abstractly isomorphic, related by the Atiyah flop), this becomes a
non-trivial automorphism
\[
  \Mu_{0,\star}^{+} \;\in\; \mathrm{Aut}_{\mathrm{bialg}}(\cH(\tilde Y))
\]
equal to the Kontsevich--Soibelman wall-crossing cocycle at the flop wall
(KS 2008 \S6; Nagao 2010, arXiv:\texttt{1002.4884}).

\emph{Braid structure at two vertices.} For $|I| = 2$, the braid group
$B_2 \simeq \Z$ is freely generated by a single element $\mu_0 \mu_1$
(the product of the two mutations, acting as translation on the cluster
tube). There is no ``$\mu_0 \mu_1 \mu_0 = \mu_1 \mu_0 \mu_1$'' relation
because $B_2$ has no braid relation; the correct statement is that at the
conifold, the braid-type subgroup of $\mathrm{Aut}_{\mathrm{bialg}}(\cH(\tilde Y))$
generated by the mutation cocycle is cyclic, and its infinite order
corresponds to the infinite stack of conifold walls in the KS wall-crossing
diagram. This matches Nagao's tilting-theoretic derivation of the conifold
partition function (Nagao--Nakajima, arXiv:\texttt{0809.2992}).
\end{example}

\begin{example}[Local $\bP^2$ nine-arrow Beilinson quiver and $B_3$-braid]
\label{ex:local-p2-mutations}
\ClaimStatusDerivation\ The Beilinson tilting
$T = \cO_{\bP^2} \oplus \cO_{\bP^2}(1) \oplus \cO_{\bP^2}(2)$ on
$K_{\bP^2} = \Tot(\cO_{\bP^2}(-3))$ gives a quiver $Q_{\bP^2}$ on three
vertices $I = \{0, 1, 2\}$, with arrow spaces
\[
  Q_{\bP^2, 1}(i, j) \;=\;
  \begin{cases}
    \Hom(\cO(i), \cO(j)) \;=\; H^{0}(\bP^2, \cO(j - i)) \;=\; \C^{3} &
    \text{if } j = i + 1, \\
    0 & \text{otherwise,}
  \end{cases}
\]
yielding three strands of three parallel arrows $x_\alpha \colon 0 \to 1$,
$y_\alpha \colon 1 \to 2$, and (by descent through the $(-3)$ twist)
$z_\alpha \colon 2 \to 0$, $\alpha = 1, 2, 3$, indexed by a basis of
$H^0(\bP^2, \cO(1)) = \C^3$. The superpotential is the antisymmetric cubic
\[
  W_{\bP^2} \;=\; \sum_{\sigma \in S_3} \mathrm{sgn}(\sigma)\,
  x_{\sigma(1)}\, y_{\sigma(2)}\, z_{\sigma(3)}
\]
(Beilinson 1978; Beasley--Plesser 2001). The Jacobi algebra
$\Jac(Q_{\bP^2}, W_{\bP^2})$ is the NCCR of $\C^3/\Z_3$ in the McKay
convention via Craw--Ishii 2004.

\emph{Three generators.} Mutation at each vertex gives an operator
$\mu_k \in \mathrm{Aut}_{\mathrm{r.e.}}(\text{QP})$ for $k \in \{0, 1, 2\}$.
Because $W_{\bP^2}$ is cyclically symmetric and $Q_{\bP^2}$ is
vertex-transitive under $\Z/3 \subset S_3$, the three $\mu_k$ are pairwise
related by the cyclic automorphism of $Q_{\bP^2}$.

\emph{$B_3$-braid as Bondal--Polishchuk autoequivalences.} The braid group
$B_3 = \langle \mu_0, \mu_1 \mid \mu_0 \mu_1 \mu_0 = \mu_1 \mu_0 \mu_1 \rangle$
(two-generator presentation) acts on $D^b(\mathrm{Coh}\, \bP^2)$ via the
Bondal--Polishchuk exceptional-collection twists $T_{\cO(i)}$
(Bondal--Polishchuk 1993, ``Homological properties of associative algebras:
the method of helices''), extended to $D^b(\mathrm{Coh}\, K_{\bP^2})$ via the
total-space pullback. On the Ginzburg dg-algebra side, each Bondal--Polishchuk
twist $T_i$ corresponds to a spherical twist at the simple $S_i$, and
Keller--Yang 2011 Prop.~6.4 identifies the mutation $\mu_i$ with the inverse
of this spherical twist in $\cD(\Gamma(Q_{\bP^2}, W_{\bP^2}))$.

The braid relation $\mu_0 \mu_1 \mu_0 \simeq_{\mathrm{r.e.}} \mu_1 \mu_0
\mu_1$ is the cluster-combinatorial shadow of the $\mathrm{SL}_2(\Z)$-covariance
of $\bP^2$: the $B_3$-braid is precisely $B_3 = \pi_1(\mathrm{Conf}_3(\C))$,
acting on the $3$-helix $(\cO, \cO(1), \cO(2))$ via the monodromy of the
Kuznetsov exceptional-collection moduli (Kuznetsov 2007, arXiv:\texttt{math/0610957};
Bridgeland 2009, arXiv:\texttt{0909.4299}).

\emph{Transport to CoHA.} By
Theorem~\ref{thm:mutation-coha-automorphism}, each $\mu_k$ transports to an
automorphism $\mu_{k, \star} \in \mathrm{Aut}_{\mathrm{bialg}}(\cH(K_{\bP^2}))$,
and the braid relation $\mu_{0,\star} \mu_{1,\star} \mu_{0,\star} =
\mu_{1,\star} \mu_{0,\star} \mu_{1,\star}$ holds in the automorphism group.
Under the identification $\cH(K_{\bP^2}) \simeq Y_{1, 0, 0}[\psi]$ with the
Gaiotto--Rap\v{c}\'ak corner VOA (Gaiotto--Rap\v{c}\'ak 2018,
arXiv:\texttt{1703.00982}), this $B_3$-braid realises the triality
$\epsilon_1 \leftrightarrow \epsilon_2 \leftrightarrow \epsilon_3$ of the
three toric $\C^\times$-weights at the McKay vertex of $\C^3/\Z_3$.
\end{example}

\subsection*{The mutation cocycle as a braid-group representation}

The mutations $\{\mu_k\}_{k \in I}$ satisfy the braid relations of $B_{|I|}$
at two distinct layers: at the level of derived equivalences, via Keller--Yang
Theorem~\ref{thm:keller-yang-mutation}; and at the level of CoHA bialgebra
automorphisms, via transport through
Theorem~\ref{thm:mutation-coha-automorphism}. The combined statement is a
bialgebra representation of the braid group.

\begin{theorem}[Mutation cocycle; braid-group representation on the CoHA]
\label{thm:mutation-cocycle}
\ClaimStatusTheorem\ Let $X$ be a smooth local CY$_3$ with tilting bundle
$T = \bigoplus_{i \in I} T_i$ and non-degenerate QP presentation $(Q, W)$.
Denote by $B_{|I|}$ the Artin braid group on $|I|$ strands, with standard
generators $\sigma_0, \ldots, \sigma_{|I| - 2}$ subject to
\[
  \sigma_i \sigma_{i+1} \sigma_i \;=\; \sigma_{i+1} \sigma_i \sigma_{i+1},
  \qquad
  \sigma_i \sigma_j \;=\; \sigma_j \sigma_i \quad (|i - j| \geq 2).
\]
The assignment $\sigma_k \mapsto \Mu_k^{+}$ (Keller--Yang equivalence at
vertex $k$) extends to a group homomorphism
\[
  \rho_X^{\mathrm{der}} \;:\; B_{|I|} \;\longrightarrow\;
  \mathrm{Auteq}\bigl(\cD(\Gamma(Q, W))\bigr)
\]
into the group of triangulated autoequivalences, and, composing with the
Hall-product transport $\Mu_{k,\star}$ of
Theorem~\ref{thm:mutation-coha-automorphism}, to
\[
  \rho_X \;=\; (-)_{\star} \circ \rho_X^{\mathrm{der}} \;:\; B_{|I|}
  \;\longrightarrow\; \mathrm{Aut}_{\mathrm{bialg}}(\cH(X)),
  \qquad \sigma_k \;\longmapsto\; \Mu_{k,\star}^{+}.
\]
The image $\rho_X(B_{|I|})$ is the \emph{mutation cocycle subgroup} of
bialgebra automorphisms of $\cH(X)$.
\end{theorem}

\begin{proof}[Proof.]
Two parts: the braid relations at the derived level, and the functoriality
of Hall-product transport.

(i) \emph{Braid relations at the derived level.} By
Theorem~\ref{thm:keller-yang-mutation}, each $\Mu_k^{+}$ is the derived
equivalence induced by the tilting object $T_k^{+} = \mathrm{cone}(P_k \to
\bigoplus_{b \colon k \to j} P_{t(b)}) \oplus \bigoplus_{i \neq k} P_i$.
Keller--Yang 2011 Prop.~6.4 identifies $\Mu_k^{+}$ with the inverse of the
spherical twist $\mathrm{Tw}_{S_k}^{-1}$ at the simple module $S_k$ over
$\Gamma(Q, W)$. The spherical-twist functors $\mathrm{Tw}_{S_k}$ satisfy
the Seidel--Thomas braid relations
\[
  \mathrm{Tw}_{S_k} \mathrm{Tw}_{S_\ell} \mathrm{Tw}_{S_k}
  \;=\;
  \mathrm{Tw}_{S_\ell} \mathrm{Tw}_{S_k} \mathrm{Tw}_{S_\ell}
  \quad\text{if}\quad \mathrm{Ext}^{*}(S_k, S_\ell) = \C[-1] \oplus \C[-2],
\]
i.e., when the two simples are ``adjacent'' in the sense of Seidel--Thomas
2001 arXiv:\texttt{math/0001043} Thm.~1.2, and commute when
$\mathrm{Ext}^{*}(S_k, S_\ell) = 0$. For $\Gamma(Q, W)$ a CY$_3$ Ginzburg
algebra, Keller--Yang Thm.~6.4 verifies that $\Ext^{*}(S_k, S_\ell)$ has
this form precisely when $k$ and $\ell$ are connected by an arrow in $Q$
(i.e.\ adjacent in the $2$-coloured exchange graph); the commutation holds
when $k, \ell$ are disconnected.

Therefore the assignment $\sigma_k \mapsto \Mu_k^{+}$ respects both the
braid relation (for adjacent $k, \ell$) and the commutation relation (for
non-adjacent), and extends to the group homomorphism $\rho_X^{\mathrm{der}}$
on $B_{|I|}$.

(ii) \emph{Functoriality of Hall-product transport.} By
Theorem~\ref{thm:hall-product-transport}, every derived Morita equivalence
between two CY$_3$ Ginzburg algebras induces a bialgebra isomorphism on the
critical CoHAs. This assignment is functorial: the composition
$\Mu_k^{+} \circ \Mu_\ell^{+}$ transports to $\Mu_{k,\star}^{+} \circ
\Mu_{\ell,\star}^{+}$ by the universal property of critical cohomology with
respect to derived Morita, which descends from the functoriality of
DHSM 2023 arXiv:\texttt{2303.12592} Thm.~A.

Combining (i) and (ii): $\rho_X = (-)_{\star} \circ \rho_X^{\mathrm{der}}$
is a group homomorphism from $B_{|I|}$ to
$\mathrm{Aut}_{\mathrm{bialg}}(\cH(X))$.
\end{proof}

\begin{corollary}[Chamber-independence of the CoHA bialgebra]
\label{cor:chamber-independence}
\ClaimStatusTheorem\ Let $X$ be as above. The CoHA bialgebra $\cH(X)$,
viewed as an abstract bialgebra, is \emph{invariant} under the mutation
cocycle: for any $\gamma \in B_{|I|}$ and any homogeneous elements $x, y
\in \cH(X)$, writing $\star$ for the Hall product,
\[
  \rho_X(\gamma)(x \star y) \;=\; \rho_X(\gamma)(x) \star \rho_X(\gamma)(y),
  \qquad
  \rho_X(\gamma) \Delta \;=\; \Delta\, \rho_X(\gamma).
\]
The chamber-dependent structure carried by $\cH(X)$ is not the bialgebra
itself but the \emph{PBW/BPS decomposition} of $\cH(X)$: in each chamber
$\mathfrak{C} \subset \mathrm{Stab}(X)$, the BPS sheaf
$\mathcal{BPS}_{\mathfrak{C}}$ on $\cM(X)$ picks out a different
$t$-structure heart, and the PBW basis of $\cH(X)$ depends on this choice.
The mutation cocycle translates between PBW bases but leaves the bialgebra
structure constants invariant.
\end{corollary}

\begin{proof}[Proof.]
Invariance under each generator $\sigma_k$: $\Mu_{k,\star}^{+}$ is by
Theorem~\ref{thm:mutation-coha-automorphism} a bialgebra isomorphism. The
braid-group extension $\rho_X$ is a composition of generators, hence
bialgebra-preserving.

Chamber dependence of the PBW basis: KS 2008 \S6 proves that the BPS sheaf
$\mathcal{BPS}_{\mathfrak{C}}$ depends on the stability condition
$\mathfrak{C}$, and that the PBW-type basis of $\cH(X)$ indexed by
BPS states is chamber-indexed. Davison--Meinhardt 2020
arXiv:\texttt{1601.02479} Thm.~A establishes that the \emph{bialgebra}
$\cH(X)$ is chamber-independent, with the PBW decomposition providing the
chamber-dependent structure.
\end{proof}

The mutation cocycle is the algebraic shadow of the Donaldson--Thomas
wall-crossing formula (Kontsevich--Soibelman 2008 \S6): each wall-crossing
transformation is a chamber-to-chamber BPS reindexing, and the mutation
cocycle realises this at the level of bialgebra automorphisms, preserving
the underlying $\cH(X)$.

% =======================================================================
\section*{Summary and placement}
% =======================================================================

Part IV has collapsed the \v{C}ech and McKay atlases of Parts I--III onto a
single finite-dimensional algebra $A = \End(T)^{\mathrm{op}}$ with Ginzburg
dg-lift $\Gamma(Q, W)$ via the Van den Bergh derived Morita equivalence
(Theorem~\ref{thm:vdb-equivalence}). The CY$_3$ property is carried by the
dg-lift $\Gamma$, not by its degree-zero cohomology $\Jac(Q, W) = H^0(\Gamma)$
(Proposition~\ref{prop:cy3-on-gamma-not-j}). The tilting bundle
$T = \bigoplus_i T_i$ serves as universal gluing data
(Theorem~\ref{thm:tilting-jacobi}): all chart-wise transitions fold into the
Jacobi relations $\partial_a W = 0$. The derived Morita equivalence is the
adjoint pair $L \dashv R$ satisfying compact generation and higher-Ext
vanishing (hypotheses (H1)--(H2) of Theorem~\ref{thm:derived-morita-tilting}).
The Hall product transports along any such equivalence between tilting algebras
(Theorems~\ref{thm:moduli-morita}--\ref{thm:hall-product-transport}, DHSM 2023
arXiv:\texttt{2303.12592} Thms.~A--B), and the
Derksen--Weyman--Zelevinsky--Keller--Yang mutation machinery
(Theorems~\ref{thm:keller-yang-mutation}--\ref{thm:mutation-cocycle},
arXiv:\texttt{0704.0649} and arXiv:\texttt{0906.0761})
realises the Artin braid group $B_{|I|}$ as explicit CoHA bialgebra
automorphisms $\rho_X \colon B_{|I|} \to \mathrm{Aut}_{\mathrm{bialg}}(\cH(X))$,
with chamber-independence
(Corollary~\ref{cor:chamber-independence}) of the bialgebra structure under
this action. The
four subsections above furnish the tilting toolkit that downstream parts of
the gluing chapter deploy: Part V (factorisation-algebra gluing) extends the
tilting picture to non-tilting-admitting varieties via local factorisation
presentations; Part VI (wall-crossing) frames the mutation cocycle as the
chamber-to-chamber KS quantum-dilogarithm factorisation; Part VIII (the
K3 $\times$ E master example) invokes the \emph{Serre-equivariant quasi-NCCR}
as the obstruction-aware analogue of Van den Bergh, where a global tilting
fails but a locally-defined Serre-equivariant tilting patches via the
factorisation pushforward.


\input{gluing_ch_part_5_factorization}

\input{gluing_ch_part_6_wall_crossing}

% !TEX root = ../main.tex
%
% Gluing Chapter, Part VII
% Lattice-polarised / automorphic gluing (Stratum iv)
%
% First-principles elementary derivation. CG voice.

\part{Lattice-polarised / automorphic gluing (Stratum iv)}
\label{part:gluing-lattice-automorphic}

The preceding parts assembled CoHAs by Čech, orbifold-McKay, tilting, factorisation-algebra, and wall-crossing cocycles. Each of those cocycles is \emph{geometric}: it lives on the variety $X$ or on a stack quotient of $X$. Part VII opens a different register. The gluing datum we now unpack does not live on $X$ at all; it lives on the period domain of $X$. Its Fourier coefficients, read off a single meromorphic Siegel modular form, determine the automorphic/BKM target arithmetic and the Euler or superdimension shadow visible from that target. CoHA multiplication constants require a finite Hall source fixture and a recognition comparison with this shadow. This is the automorphic lane of the gluing atlas, and it is the register on which the Borcherds-Monster and K3~$\times$~E data were always going to be read.

We give the derivation in four steps. First, the Mukai lattice and the Kuga--Satake period map (\S\ref{gluech:7.1}). Second, weak Jacobi forms and the Borcherds lift (\S\ref{gluech:7.2}). Third, the Gritsenko--Cl\'ery eight-form family and the universal identity $\kappa_{\mathrm{BKM}}(\Phi_N) = c_N(0)/2$ (\S\ref{gluech:7.3}). Fourth, the $\Delta_5$ / K3~$\times$~E gluing, which assembles the local Hall--Drinfeld pieces into the global BKM algebra $\mathfrak{g}_{\Delta_5}$ (\S\ref{gluech:7.4}).

\section{The Mukai lattice and the period map}
\label{gluech:7.1}

\subsection{Even integral cohomology of K3 as a lattice}

Fix a K3 surface $X$ over $\C$. Its even integral cohomology
\[
  \widetilde H(X, \Z) \;:=\; H^0(X, \Z) \oplus H^2(X, \Z) \oplus H^4(X, \Z)
\]
has rank $1 + 22 + 1 = 24$. The cup product makes $H^2(X, \Z)$ into the $\mathrm{K3}$ lattice $E_8(-1)^{\oplus 2} \oplus U^{\oplus 3}$ of signature $(3, 19)$, where $U$ denotes the hyperbolic plane $\begin{psmallmatrix} 0 & 1 \\ 1 & 0 \end{psmallmatrix}$ and $E_8(-1)$ the negative-definite $E_8$ root lattice. The unit classes $1 \in H^0$ and $[\mathrm{pt}] \in H^4$ contribute a further hyperbolic plane when one pairs them by Poincar\'e duality. Altogether
\[
  \widetilde H(X, \Z) \;\cong\; U^{\oplus 4} \oplus E_8(-1)^{\oplus 2} \;=\; \mathrm{II}_{4, 20}.
\]
The subscripts $(4, 20)$ record the signature.

\subsection{The Mukai pairing}

Cup product alone does not see the grading $0 + 2 + 4$; it symmetrises between $H^0$ and $H^4$ but forgets them. Mukai's observation is that a \emph{sign-twist} on the middle degree yields the correct pairing for Hodge-isometric period geometry. For
\[
  v \;=\; (v_0, v_1, v_2), \qquad w \;=\; (w_0, w_1, w_2),
\]
write $v_0, w_0 \in H^0$, $v_1, w_1 \in H^2$, $v_2, w_2 \in H^4$. The \emph{Mukai pairing} is
\begin{equation}
  \langle v, w\rangle_M \;=\; v_0 \smile w_2 \;+\; v_2 \smile w_0 \;-\; v_1 \smile w_1
  \label{eq:mukai-pairing}
\end{equation}
with the minus sign on the middle contraction. Under the identifications $H^0 \cong \Z$, $H^4 \cong \Z$ (via the fundamental class) this is
\[
  \langle (r, L, s), (r', L', s')\rangle_M \;=\; r s' + r' s - L \cdot L'.
\]
The minus sign is forced by the requirement that the Chern character
\[
  \mathrm{ch}\colon K_0(X) \longrightarrow \widetilde H(X, \Q),
  \qquad
  \mathrm{ch}([\mathcal F]) = \bigl(\mathrm{rk}\, \mathcal F, \; c_1(\mathcal F), \; \mathrm{ch}_2(\mathcal F)\bigr)
\]
turn Euler characteristic $\chi(\mathcal F, \mathcal G) := \sum_i (-1)^i \dim \mathrm{Ext}^i(\mathcal F, \mathcal G)$ into Mukai pairing on the Mukai vector $v(\mathcal F) := \mathrm{ch}(\mathcal F) \cdot \sqrt{\mathrm{td}(X)}$: for K3 with trivial canonical bundle,
\[
  \chi(\mathcal F, \mathcal G) \;=\; -\langle v(\mathcal F), v(\mathcal G)\rangle_M.
\]
The Riemann--Roch computation is elementary: $\chi = \int_X \mathrm{ch}(\mathcal F^\vee \otimes \mathcal G) \cdot \mathrm{td}(X)$, expand, and the cross term between $\mathrm{ch}_1(\mathcal F^\vee)$ and $\mathrm{ch}_1(\mathcal G)$ produces the $-L \cdot L'$ summand. That single minus sign is what distinguishes $\mathrm{II}_{4, 20}$ from the totally split $(3, 21)$ form one might na\"ively write down; it encodes Poincar\'e duality plus Hirzebruch--Riemann--Roch simultaneously.

\subsection{Three-graded Mukai vectors and rank/$c_1$/$c_2$}

A coherent sheaf $\mathcal F$ on $X$ carries three primary discrete invariants: its rank $r = \mathrm{rk}\, \mathcal F$, first Chern class $c_1(\mathcal F) \in H^2(X, \Z)$, and Chern character $\mathrm{ch}_2(\mathcal F) = \tfrac12 c_1^2 - c_2 \in H^4(X, \Z)$. The Mukai vector repackages these into a single element of $\widetilde H$:
\[
  v(\mathcal F) \;=\; \mathrm{ch}(\mathcal F) \cdot \sqrt{\mathrm{td}(X)} \;=\; \bigl(r, \; c_1, \; \tfrac12 c_1^2 - c_2 + r\bigr).
\]
The degree-$4$ component acquires $+r$ from the Todd square root: for K3, $c_1(X) = 0$ and $c_2(X) = 24$ give $\mathrm{td}_2(X) = \tfrac{1}{12}(c_1(X)^2 + c_2(X)) = 2$, whence $\mathrm{td}(X) = 1 + 2[\mathrm{pt}]$ and $\sqrt{\mathrm{td}(X)} = 1 + [\mathrm{pt}]$ (the $[\mathrm{pt}]$-coefficient of the square root being $\mathrm{td}_2(X)/2 = 1$ since higher-order corrections vanish on a surface). Multiplication by $1 + [\mathrm{pt}]$ shifts the degree-$4$ Chern slot by the rank $r$. The \emph{three-graded structure} of $\widetilde H = H^0 \oplus H^2 \oplus H^4$ is visible in this display: rank sits in the first slot, $c_1$ in the second, the Todd-shifted $(\tfrac12 c_1^2 - c_2 + r)$ in the third. The Mukai pairing \eqref{eq:mukai-pairing} then computes Euler pairings directly.

\subsection{Kuga--Satake correspondence to a Shimura variety}

The period domain of polarised K3 is the non-compact Hermitian symmetric space
\[
  \mathcal D_{\mathrm{K3}} \;=\; \bigl\{\,\Omega \in \mathbb P\bigl(H^2(X, \C)\bigr) : \Omega \cdot \Omega = 0,\; \Omega \cdot \overline\Omega > 0\,\bigr\}
\]
with action of the orthogonal group $\mathrm{O}(\Lambda_{K3})$ of the K3 lattice. Its dimension is $20$, matching $h^{1,1}(X) = 20$.

The \emph{Kuga--Satake construction} sends $X$ to a complex abelian variety $\mathrm{KS}(X)$ of dimension $2^{19}$ with $\End\bigl(\mathrm{KS}(X)\bigr) \supset \mathrm{Cl}^+(H^2(X, \R))$, the even Clifford algebra of the K3 lattice. The periods of $\mathrm{KS}(X)$ determine the periods of $X$: there is a canonical morphism of period domains
\[
  \mathrm{KS}\colon \mathcal D_{\mathrm{K3}} \hookrightarrow \mathcal H_{2^{19}}
\]
identifying K3 periods with a sub-locus of the genus-$2^{19}$ Siegel upper half-space. This embeds Hodge-theoretic K3 data into Shimura-varietal data, opening the door to automorphic forms.

For the Mukai-enhanced lattice $\mathrm{II}_{4, 20}$ the relevant period domain enlarges to
\[
  \mathcal D_{4, 20} \;=\; \bigl\{\, \Omega \in \mathbb P(\widetilde H_\C) : \langle \Omega, \Omega\rangle_M = 0, \; \langle \Omega, \overline\Omega\rangle_M > 0 \,\bigr\}
\]
which is of complex dimension $22$. This is the natural home for Bridgeland stability conditions: the positive cone $\{\,\Omega : \langle \Omega, \overline\Omega\rangle_M > 0\,\}$ is precisely where $Z(v) = \langle \Omega, v\rangle_M$ defines a central charge satisfying the support property. The period map to $\mathcal D_{4, 20}$, modulo the arithmetic group $\mathrm{O}(\mathrm{II}_{4, 20})$, is the target of the Borcherds lift.

\subsection{Humbert divisors as polarisation strata}

Inside $\mathcal A_2 = \mathcal H_2 / \mathrm{Sp}_4(\Z)$ the \emph{Humbert divisors} $\mathcal H_D \subset \mathcal A_2$ parametrise abelian surfaces whose endomorphism ring contains an order of discriminant $D$. The divisor $\mathcal H_1$ corresponds to products $E_1 \times E_2$ of elliptic curves; $\mathcal H_4$ to $E \times E$ with $E$ having an order-$4$ endomorphism. Under Kuga--Satake these pull back to sub-loci of the K3 period domain where the N\'eron--Severi lattice acquires extra classes. The \emph{lattice-polarised} K3s of rank-$n$ N\'eron--Severi live on such Humbert-style loci, and the moduli space $\mathcal F_\Lambda$ of $\Lambda$-polarised K3s is an $(20 - \mathrm{rk}\, \Lambda)$-dimensional arithmetic quotient of $\mathcal D_{\mathrm{K3}}$. It is on these arithmetic quotients that Borcherds lifts live.

\section{Borcherds lifts}
\label{gluech:7.2}

\subsection{Weak Jacobi forms: the elementary definition}

A \emph{weak Jacobi form} of weight $k \in \Z$ and index $m \in \tfrac12 \Z$ is a holomorphic function
\[
  \phi\colon \mathcal H \times \C \longrightarrow \C,
  \qquad (\tau, z) \longmapsto \phi(\tau, z),
\]
satisfying two transformation laws and a Fourier expansion.

\textbf{Modular law.} For $\gamma = \begin{psmallmatrix} a & b \\ c & d\end{psmallmatrix} \in \mathrm{SL}_2(\Z)$,
\[
  \phi\!\left(\frac{a\tau + b}{c\tau + d}, \; \frac{z}{c\tau + d}\right) \;=\; (c\tau + d)^k \, e^{2\pi i m c z^2 / (c\tau+d)} \, \phi(\tau, z).
\]

\textbf{Heisenberg (elliptic) law.} For $(\lambda, \mu) \in \Z^2$,
\[
  \phi(\tau, z + \lambda \tau + \mu) \;=\; e^{-2\pi i m (\lambda^2 \tau + 2\lambda z)} \, \phi(\tau, z).
\]

\textbf{Fourier expansion.} The double periodicity plus the modular law force a Fourier series
\[
  \phi(\tau, z) \;=\; \sum_{n, \ell} c(n, \ell) \, q^n y^\ell,
  \qquad q := e^{2\pi i \tau}, \quad y := e^{2\pi i z},
\]
where $(n, \ell)$ runs over $\Z \times \Z$ (resp.\ $\Z \times (\Z + \tfrac12)$ if $m$ is half-integral). A \emph{weak} Jacobi form allows negative $n$ with $4mn - \ell^2 \geq -N$ for some $N \geq 0$; equivalently, the polar part at the cusp is controlled.

The space $J_{k, m}^{\mathrm{weak}}$ is a finite-dimensional $\C$-vector space, and its structure for small $(k, m)$ is computable by a Riemann--Roch argument on the modular curve; classical references are Eichler--Zagier. For K3 the critical instance is
\[
  \phi_{0, 1}^{\mathrm{K3}} \;\in\; J_{0, 1}^{\mathrm{weak}},
\]
a weight-$0$, index-$1$ form whose Fourier coefficients compute the elliptic genus of K3:
\[
  \phi_{0, 1}^{\mathrm{K3}}(\tau, z) \;=\; 4 \left[ \frac{\theta_2(\tau, z)^2}{\theta_2(\tau, 0)^2} + \frac{\theta_3(\tau, z)^2}{\theta_3(\tau, 0)^2} + \frac{\theta_4(\tau, z)^2}{\theta_4(\tau, 0)^2} \right].
\]
The Laurent expansion at the cusp reads
\[
  \phi_{0, 1}^{\mathrm{K3}}(\tau, z) \;=\; q^{-1}(y^{-1} + 10 + y) \;+\; O(q),
\]
so the polar-part discriminant is $c(-1, \pm 1) = 1$ and the singular-theta constant term is $c_1(0, 0) = 10$; the Witten index at $y = 1$ is the Euler number $\chi(K3) = 24$. The constant $c_1(0, 0) = 10$ is the input to the Borcherds weight formula $k_{\Psi} = c_\phi(0, 0)/2$, yielding $k_{\Delta_5} = 5$.

\subsection{The Borcherds product formula}

The passage from Jacobi form to automorphic form on a bigger orthogonal group is Borcherds's \emph{singular theta lift}. For a weak Jacobi form $\phi = \sum c(n, \ell) q^n y^\ell$ of index $m$, the Borcherds lift is the function
\begin{equation}
  \Psi_\phi(v) \;=\; e^{2\pi i \langle \rho, v\rangle} \prod_{\alpha > 0} \bigl(1 - e^{2\pi i \langle \alpha, v\rangle}\bigr)^{c(|\alpha|^2/2)}
  \label{eq:borcherds-product}
\end{equation}
on a Grassmannian $\mathcal G$ attached to a lattice $L$ of signature $(n, 2)$. Unpacking:

\begin{itemize}
  \item The product runs over positive roots $\alpha$ of an infinite-dimensional root system inside $L$; \emph{positive} means lying on a fixed side of a Weyl wall.
  \item The Weyl vector $\rho$ is determined by $\phi$; its role is to absorb the leading term of the product.
  \item The exponent $c(|\alpha|^2 / 2)$ is a Fourier coefficient of $\phi$; only the value of $n - \ell^2/(4m)$, i.e.\ $|\alpha|^2/2$, enters.
  \item The product converges in a neighbourhood of the cusp; the singular theta lift extends it meromorphically to the full Hermitian symmetric domain.
\end{itemize}

The key fact is Borcherds's \emph{automorphy theorem} (Borcherds, \emph{Invent.\ Math.}~132 (1998), Theorem~13.3): $\Psi_\phi$ is a meromorphic automorphic form of weight $k_\Psi$ on $\mathrm{O}^+(L)$, with divisor supported on Heegner divisors (rational quadratic divisors of $\mathcal G$) and order of zero/pole along each Heegner divisor determined by the principal-part Fourier coefficients of $\phi$. The weight is
\begin{equation}
  k_\Psi \;=\; \tfrac12 \, c_\phi(0, 0),
  \label{eq:borcherds-weight}
\end{equation}
where $c_\phi(0, 0)$ is the cusp constant term of the vector-valued input at the principal $0$-dimensional cusp. This is Theorem 13.3 of Borcherds 1998 (``The singular theta correspondence and automorphic forms''), not Theorem 10.1, which records only the singular-theta automorphicity without the weight formula. Identity \eqref{eq:borcherds-weight} is the direct source of the $\kappa_{\mathrm{BKM}}(\Phi_N) = c_N(0, 0)/2$ identity we pursue below.

\subsection{Modular covariance under $\mathrm{O}(\mathrm{II}_{n, 2})$}

The lift $\Psi_\phi$ transforms under the arithmetic subgroup
\[
  \mathrm{O}^+(\mathrm{II}_{n, 2}) \;\subset\; \mathrm{O}(\mathrm{II}_{n, 2})
\]
of the even unimodular lattice of signature $(n, 2)$, with a specified character (typically trivial modulo a cover, detailed for each $\phi$). For our purposes $n = 2$ (Igusa $\Phi_{10}$ lives on $\mathrm{II}_{2, 2}$) and $n = 3$ (Gritsenko $\Delta_5$ on $\mathrm{II}_{3, 2}$) are decisive. Modular covariance means: the Fourier coefficients of $\Psi_\phi$, read in a neighbourhood of any cusp of $\mathcal D / \mathrm{O}^+$, determine the \emph{same} form globally; equivalently, the BKM root-multiplicity and Euler/superdimension data agree across cusps of the period domain. A Hall-algebra multiplication table is compared with this automorphic table only after choosing a finite source fixture.

\subsection{Automorphic coefficients as BKM target data}

The gluing statement we extract from all this is concrete. Consider a BKM-type algebra $\mathfrak g$ whose denominator identity takes the form
\begin{equation}
  \Psi_\phi(v) \;=\; \sum_{w \in W} \mathrm{sgn}(w) \, w\!\left(e^{2\pi i \langle \rho, v\rangle} \prod_{\alpha > 0}^{\mathrm{real}} \bigl(1 - e^{2\pi i \langle \alpha, v\rangle}\bigr)\right)
  \label{eq:bkm-denominator}
\end{equation}
with Weyl group $W$ acting on positive roots. Expanding both sides of \eqref{eq:borcherds-product} = \eqref{eq:bkm-denominator} and comparing Fourier coefficients determines the denominator-side arithmetic of $\mathfrak g$: imaginary-root multiplicities and the superdimension/Euler coefficients of root spaces appear as Fourier coefficients $c_\phi(n, \ell)$ with specific $(n, \ell)$. The Cartan matrix, parity conventions, and Serre relations are part of the BKM presentation; they are not recovered as Hall multiplication constants by coefficient projection alone.

For a CoHA-side identification, suppose $\mathfrak g = \mathfrak g(X)$ is the BKM associated to a CY$_3$ compactification $X$ with lattice-polarised period (Stratum iv). The CoHA product
\[
  \mathcal H(X) \otimes \mathcal H(X) \longrightarrow \mathcal H(X)
\]
is source-side data. One first fixes a finite Hall source fixture, computes its multiplication constants in that fixture, and then projects to the Euler/superdimension shadow indexed by a rank-$2$ Cartan-torus direction $v \in \widetilde H(X)$. The theorem to prove is equality of that projected table with the automorphic table $c_\phi(n, \ell)$, not an a priori identification of CoHA products with Fourier coefficients. This is the \emph{automorphic cocycle} on the period side: CoHA locally looks like a shuffle algebra on each lattice cusp; globally it is tested against the Borcherds-lifted target, and the target data are the Fourier coefficients of $\Psi_\phi$. The explicit chain-level statement, for the K3~$\times$~E case, is the content of \S\ref{gluech:7.4}.

\section{The Gritsenko--Cl\'ery eight-form family}
\label{gluech:7.3}

\subsection{Two scopes of the Borcherds identity}

Write $J_{0, m}^{\mathrm{weak}}$ for the space of weight-$0$, index-$m$ weak Jacobi forms. Two indexings of the Gritsenko--Cl\'ery programme recur in the literature, enumerating distinct scopes of the universal identity $\kappa_{\mathrm{BKM}}(\Phi_N) = c_N(0, 0)/2$.

\textbf{Scope (A): CHL ladder.} At the integer levels
\[
  N \in \{1, 2, 3, 4, 6\},
\]
cyclic $\Z/N$ symplectic orbifolds of $K3 \times E$ (Chaudhuri--Hockney--Lykken 1995) admit Borcherds products $\Phi_N$ whose Jacobi inputs are $\phi_{0, 1}^{\mathrm{K3}} \otimes T_N$ for explicit Hecke-type twists $T_N$. The cusp constant terms tabulate as
\begin{equation}
  \bigl(c_N(0, 0)\bigr)_{N \in \{1, 2, 3, 4, 6\}} \;=\; (10, \; 8, \; 6, \; 4, \; 2),
  \label{eq:cn-zero-chl}
\end{equation}
matching Gritsenko--Nikulin \emph{Automorphic forms and Lorentzian Kac--Moody algebras} (Part II, Theorem~2.1, 1995), Eguchi--Hikami 2011 for the $N = 2$ Mathieu-twined entry, and Govindarajan--Krishna 2010 for $N \in \{3, 4, 6\}$. The level $N = 5$ is absent: the elliptic curve $E$ carries no automorphism of order $5$, so the CHL orbifold construction stops at $N = 6$.

\textbf{Scope (B): Full Gritsenko--Cl\'ery eight-form atlas.} The paramodular-form classification programme (Gritsenko--Cl\'ery, \emph{The eight-form Gritsenko--Cl\'ery catalogue}, 2013 and 2018) enumerates eight canonical reflective automorphic products by paramodular type:
\[
  \phi_{0, m} \in J_{0, m}^{\mathrm{weak}}, \qquad m \in \{1, 2, 3, 4, 5, 6\}, \qquad \phi_{0, 1/2}, \; \phi_{0, 3/2}.
\]
Indexing the eight forms by their position in the catalogue and denoting the Borcherds weights by $w_N$, the per-form table reads
\begin{equation}
  \bigl(w_N\bigr)_{N = 1, \ldots, 8} \;=\; (5, \; 2, \; 1, \; 1, \; 1/2, \; 1, \; 1/4, \; 0), \qquad \bigl(c_N(0, 0)\bigr) \;=\; (10, \; 4, \; 2, \; 2, \; 1, \; 2, \; 1/2, \; 0),
  \label{eq:cn-zero-gc-atlas}
\end{equation}
with the per-position identity $w_N = c_N(0, 0)/2$ holding verbatim row by row (Gritsenko--Cl\'ery 2013 \emph{J.\ Reine Angew.\ Math.} 678 \S 3). The terminal weight-$0$ entry is the degenerate fibre without BKM content.

\textbf{Two scopes, one identity.} The two indexings coincide at cardinality $8$ by accident, not by functoriality. They agree at $N = 1$ (where both give $c_1(0,0) = 10$ and $\kappa_{\mathrm{BKM}} = 5$) and diverge at $N \in \{2, 3, 4, 6\}$: Scope (A) reads the Jacobi input $\phi_{0,1}^{\mathrm{K3}} \otimes T_N$ of weight $k(N) = N-1$ type (yielding $(8, 6, 4, 2)$), while Scope (B) reads an independent twined weight-$0$ weak Jacobi form $\phi_{0, N}^{(g_N)}$ (yielding $(4, 2, 2, 2)$ at $N = 2, 3, 4, 6$). Every invocation of the universal formula $\kappa_{\mathrm{BKM}} = c_N(0, 0)/2$ must name which scope.

\subsection{The universal identity $\kappa_{\mathrm{BKM}}(\Phi_N) = c_N(0)/2$}

Each $\phi_{0, N}$ has a Borcherds-style lift $\Phi_N$ to a Siegel modular form on $\mathrm{Sp}_4(\Z)$ (or its $\mathrm{Mp}_4$ / $\widetilde{\mathrm{Mp}}_4$ cover). The lift $\Phi_N$ has weight
\[
  k_{\Phi_N} \;=\; \tfrac12 \, c_N(0, 0)
\]
by the Borcherds automorphy theorem from \S\ref{gluech:7.2}. Identifying this weight with the Cartan--Killing invariant of the associated BKM algebra gives
\begin{equation}
  \kappa_{\mathrm{BKM}}(\Phi_N) \;=\; \tfrac12 \, c_N(0, 0)
  \label{eq:kappa-bkm-universal}
\end{equation}
in both scopes. Numerical witnesses:
\[
  \text{Scope (A, CHL):}\quad \bigl(\kappa_{\mathrm{BKM}}(\Phi_N)\bigr)_{N \in \{1, 2, 3, 4, 6\}} \;=\; (5, \; 4, \; 3, \; 2, \; 1),
\]
via the Gritsenko--Nikulin 1995 Pt II Thm 2.1 CHL ladder;
\[
  \text{Scope (B, GC atlas):}\quad \bigl(w_N\bigr)_{N = 1, \ldots, 8} \;=\; (5, \; 2, \; 1, \; 1, \; 1/2, \; 1, \; 1/4, \; 0),
\]
via Gritsenko--Cl\'ery 2013 \S 3.

\noindent The universality of \eqref{eq:kappa-bkm-universal} is that \emph{one} extraction rule --- weight equals half the cusp constant term of the Jacobi input --- governs all eight forms across both scopes. The scope declares which Jacobi input $\phi_{0, N}$ is being lifted; the extraction rule is the same.

\subsection{Cover group stratification}

The eight-form family stratifies naturally by the cover group of $\mathrm{Sp}_4$ under which the lift transforms:
\[
  \begin{array}{c|l}
    \text{weight type} & \text{cover group} \\
    \hline
    \text{integral}, \; k \in \Z_{\geq 1} & \mathrm{Sp}_4(\Z) \\
    \text{half-integral}, \; k \in \Z + \tfrac12 & \mathrm{Mp}_4(\Z) \\
    \text{quarter-integral}, \; k \in \tfrac14 \Z \setminus \tfrac12 \Z & \widetilde{\mathrm{Mp}}_4(\Z) \\
    \text{weight zero} & \text{degenerate}
  \end{array}
\]
The degenerate weight-$0$ form is the terminal fibre of the GC atlas; its BKM algebra is trivial (one-dimensional abelian), and it marks the endpoint, not content. The integer-weight forms at positions $N \in \{1, 2, 3, 4, 6\}$ sit on classical $\mathrm{Sp}_4(\Z)$-Siegel modular forms; the half-weight form (at position $5$ of the GC atlas, weight $w_5 = 1/2$) sits on a genuine $\mathrm{Mp}_4$-object; the quarter-weight form (at position $7$, weight $w_7 = 1/4$) sits on a $\widetilde{\mathrm{Mp}}_4$ cover requiring the spin double of $\mathrm{Sp}_4$ to state coherently.

\subsection{Each form is a gluing datum}

Take $\Phi_N$ fixed. The Borcherds product formula \eqref{eq:borcherds-product} expresses $\Phi_N$ as a product over positive roots of a BKM algebra $\mathfrak g_N$. The Cartan of $\mathfrak g_N$ is the Mukai lattice of a specific CY$_3$ geometry $X_N$: for $N = 1$, $X_1 = K3 \times E$ as elaborated below; for $N = 2, 3, 4, 6$, the CY$_3$ is the Nikulin--Scheithauer CHL compactification on orbifolds $(K3 \times T^2) / \Z_N$ with Mukai-type lattice reduction. The \emph{gluing datum} is: $\Phi_N$ supplies the automorphic/BKM target table for the positive-root arithmetic of $\mathfrak g_N$. A CoHA product on $\mathcal H(X_N)$ matches this datum only after a finite Hall source fixture has been chosen, its multiplication constants computed, and their Euler/superdimension projection compared with the Fourier coefficients of $\Phi_N$.

This is the Stratum-(iv) gluing in its most compact form. The local data on each equivariant chart of $X_N$ produces local Hall-algebra pieces; $\Phi_N$ is the automorphic recognition target against which their finite-source shadows are tested. Different $N$ give different gluings because different $N$ produce different lattice-polarisations on the period domain.

\section{The $\Delta_5$ and K3~$\times$~E gluing}
\label{gluech:7.4}

\subsection{From $\phi_{0, 1}^{\mathrm{K3}}$ to $\Delta_5$}

The K3 weak Jacobi form $\phi_{0, 1}^{\mathrm{K3}}$ introduced in \S\ref{gluech:7.2} is the input to Gritsenko's lift in the $N = 1$ case. Gritsenko's construction, refined from Maass's additive lift, produces a holomorphic \emph{Siegel modular form} $\mathrm{Grit}(\phi)$ of weight $\tfrac12 c_\phi(0, 0)$ on the genus-$2$ Siegel upper half-space $\mathcal H_2 = \{Z \in M_{2\times 2}(\C) : Z = Z^T,\; \mathrm{Im}\, Z > 0\}$. For $\phi = \phi_{0, 1}^{\mathrm{K3}}$ we have $c_\phi(0, 0) = 10$, hence
\[
  \Delta_5 \;:=\; \mathrm{Grit}(\phi_{0, 1}^{\mathrm{K3}}) \;\in\; M_5\!\left(\mathrm{Sp}_4(\Z)\right),
\]
a Siegel cusp form of weight $5$.

Explicitly, Gritsenko's formula writes the additive lift as
\[
  \mathrm{Grit}(\phi)(Z) \;=\; \sum_{m \geq 1} \phi|V_m(\tau, z) \, e^{2\pi i m w},
  \qquad Z = \begin{pmatrix} \tau & z \\ z & w \end{pmatrix},
\]
with $V_m$ the $m$-th Hecke operator on weak Jacobi forms, and this additive lift agrees with the Borcherds multiplicative lift on the common domain by uniqueness of weight-$5$ Siegel cusp forms.

The form $\Delta_5$ is the unique (up to scalar) weight-$5$ Siegel cusp form, vanishing on the diagonal locus $z = 0$ (which is the Humbert divisor $\mathcal H_1$ of $E_1 \times E_2$-products), and with Fourier expansion starting as
\[
  \Delta_5(Z) \;=\; \bigl(q^{1/2} \eta^{-1}(\tau) \eta^{-1}(w)\bigr) \cdot \bigl(1 - \text{O}(q)\bigr) \cdot \ldots
\]
after the Gritsenko normalisation.

\subsection{The BKM algebra $\mathfrak g_{\Delta_5}$}

The denominator identity for $\Delta_5$, written in Borcherds-product form \eqref{eq:borcherds-product} and in Weyl-sum form \eqref{eq:bkm-denominator}, produces a BKM superalgebra $\mathfrak g_{\Delta_5}$ with:

\begin{itemize}
  \item Cartan subalgebra $\mathfrak h_{\Delta_5}$ of rank $3$, with real-simple-root Gram matrix $\mathrm{diag}(2) - 2 \cdot (\mathbb 1 - \mathrm{diag}(1))$ having eigenvalues $(4, 4, -2)$ and hence signature $(2, 1)$ (Lorgat 2020 Lemma 1 and the $\Lambda^{2,1}_{II}$ sublattice of $\Lambda^{3,2}$);
  \item Real simple roots $\{\delta_1, \delta_2, \delta_3\}$ forming an $\widetilde A_2$-like Dynkin diagram with off-diagonal entries $-2$;
  \item Imaginary simple roots with multiplicities read off the Fourier coefficients $C_1(D)$ of $\phi_{0, 1}^{\mathrm{K3}}$, enumerated by the discriminant $D = 4nm - \ell^2$ of a positive root $\alpha = (n, \ell, m)$;
  \item Weight (Borcherds singular-theta normalisation) $\kappa_{\mathrm{BKM}}(\mathfrak g_{\Delta_5}) = k_{\Delta_5} = 5$.
\end{itemize}

The value $5 = c_1(0, 0)/2 = 10/2$ is the exemplar of the universal identity \eqref{eq:kappa-bkm-universal}. First-principles derivation: by Borcherds's theorem the weight of $\Delta_5$ is half the constant term of $\phi_{0, 1}^{\mathrm{K3}}$, which in the normalisation $\phi_{0, 1}^{\mathrm{K3}}(0, 0) = c_1(0, 0) = 10$ gives $5$. On the BKM side the Killing form restricted to the Cartan, together with the twist by the Weyl vector, produces precisely $c_1(0, 0)/2$ as the normalisation integer. The matching is not a coincidence; it is a \emph{structural identity} forced by the Howe-duality pair $(\mathrm{SL}_2, \mathrm{O}(n, 2))$ underlying the theta lift.

\subsection{The K3~$\times$~E geometry}

The CY$_3$ carrying $\mathfrak g_{\Delta_5}$ is
\[
  X \;=\; K3 \times E,
\]
with $E$ an elliptic curve. Its Mukai-enhanced even cohomology is
\[
  \widetilde H(X, \Z) \;=\; \widetilde H(K3, \Z) \otimes H^\bullet(E, \Z) \;=\; \mathrm{II}_{4, 20} \otimes \bigl(\Z \oplus \Z[1] \oplus \Z\bigr),
\]
and the K\"unneth-multiplicative Calabi--Yau Euler invariant is
\[
  \kappa_{\mathrm{cat}}(X) \;=\; \chi(\mathcal O_X) \;=\; \chi(\mathcal O_{K3}) \cdot \chi(\mathcal O_E) \;=\; 2 \cdot 0 \;=\; 0.
\]
(The K3 fibre value $\chi(\mathcal O_{K3}) = 2$ is \emph{not} the total-space value; K\"unneth on products is multiplicative, not additive.) The K3-fibre Mukai rank $24$ remains visible as $\kappa_{\mathrm{fiber}}(X) = 24$ of the K3 factor.

On the chiral side, Costello--Gaiotto $(2, 2)$ supersymmetric $\sigma$-model data on $K3 \times E$ produces a Heisenberg chiral shadow of central charge $3$:
\[
  \kappa_{\mathrm{ch}}^{\mathrm{Heis}}(K3 \times E) \;=\; 3,
\]
computed via the Heisenberg--Mukai specialisation of $\Phi_3$ along the $E$-factor. On the BKM side, the Borcherds lift $\Delta_5$ yields
\[
  \kappa_{\mathrm{BKM}}(\mathfrak g_{\Delta_5}) \;=\; 5.
\]

Four separate constructions produce four separate $\kappa_\bullet$ invariants, altogether
\[
  \{0, \; 3, \; 5, \; 24\} \;=\; \bigl\{\kappa_{\mathrm{cat}}(K3 \times E), \; \kappa_{\mathrm{ch}}^{\mathrm{Heis}}, \; \kappa_{\mathrm{BKM}}(\mathfrak g_{\Delta_5}), \; \kappa_{\mathrm{fiber}}(K3)\bigr\}.
\]
No two of these are numerically equal; no two of these arise from a shared formula. They are four \emph{distinct constructions} producing four distinct invariants. The $\Delta_5$-gluing speaks only to the third, $\kappa_{\mathrm{BKM}} = 5$.

\subsection{Automorphic cocycle: assembling local Hall--Drinfeld pieces}

The local $\mathbb C^3$ charts of $K3 \times E$ (in the elliptic-K3 presentation, twenty-four $I_1$-fibre loci each fibered over $E$, producing curve-supported local data of Jordan-triple type) each carry a Schiffmann--Vasserot CoHA $Y^+(\widehat{\mathfrak{gl}}_1)$. Locally these are the chart-wise CoHAs of Part II of the gluing chapter. Globally they do not glue by a Čech cocycle alone: the curve-supported nature of the $I_1 \times E$ stalks, the Aut($K3 \times E$)-equivariance, and the Humbert-monodromy action combine to enforce an automorphic gluing constraint.

The \emph{automorphic cocycle} is $\Delta_5$. The statement, to be established chain-level in the K3~$\times$~E master-example chapter (Part VIII of the gluing atlas), is as follows.

\begin{conjecture}[Automorphic assembly of $K3 \times E$ CoHA]
\label{conj:automorphic-assembly-k3xe}
Let $\mathcal H_\alpha$ denote the local Schiffmann--Vasserot CoHA on the $\alpha$-th $\mathbb C^3$-chart in the elliptic-K3 atlas, $\alpha = 1, \ldots, 24$. Then the global CoHA $\mathcal H(K3 \times E)$ fits into an exact sequence
\[
  \bigoplus_\alpha \mathcal H_\alpha \xrightarrow{\;\partial\;} \mathcal H(K3 \times E) \xrightarrow{\;\pi\;} U(\mathfrak g_{\Delta_5})^+ \longrightarrow 0
\]
in which the cokernel is the positive half of the universal enveloping algebra of $\mathfrak g_{\Delta_5}$, and the gluing cocycle $\partial$ is recognised by comparing a finite Hall source fixture with the Fourier coefficients $c_1(n, \ell)$ of $\phi_{0, 1}^{\mathrm{K3}}$. The Fourier expansion of $\Delta_5$ determines the BKM target arithmetic and the Euler/superdimension shadow of the desired assembly; the multiplication constants of $\mathcal H(K3 \times E)$ are those of the source fixture whose projected table passes this recognition test.
\end{conjecture}

The first inscription of this statement is due to the author (Lorgat, arXiv:2007.14218), where the explicit Borcherds lift of $\phi_{0, 1}^{\mathrm{K3}}$ to $\Delta_5$ was paired with the GKM superalgebra $\mathfrak g_{\Delta_5}$ and a conjectural identification with eight Gritsenko--Cl\'ery forms matching DT zeta roots and twined elliptic genera. The content of that paper is the base case $N = 1$ of the universal eight-form ladder.

The conjecture above is the \emph{gluing} restatement of Lorgat 2020: $\Delta_5$ labels the BKM algebra and supplies the explicit automorphic target for local CoHA charts. The chain-level verification involves, for each pair of adjacent $I_1$-fibre loci $\alpha \neq \beta$, choosing a finite Hall source fixture on the overlap, computing its Čech gluing 2-cocycle, and comparing its Euler/superdimension projection with the Fourier coefficient $c_1\bigl((\alpha - \beta)^2/2\bigr)$ of $\phi_{0, 1}^{\mathrm{K3}}$; the assertion is that these agree up to a coboundary, so that the projected 2-cocycle class is canonically represented by the Jacobi-Fourier expansion.

\subsection{Two scopes of the universal Borcherds statement, revisited}

Across both scopes of \S\ref{gluech:7.3}, each form $\Phi_N$ gives the automorphic/BKM target against which local CoHA charts of its host CY$_3$ must be recognised. Scope (A), the CHL ladder $N \in \{1, 2, 3, 4, 6\}$, hosts the integer-weight $\mathrm{Sp}_4(\Z)$-Siegel cusp forms of weights $(5, 4, 3, 2, 1)$ with principal-part constants $(c_N(0, 0)) = (10, 8, 6, 4, 2)$; these are the Gritsenko--Nikulin (1995, Part~II Theorem~2.1) BKM siblings on $\Z/N$-orbifolded $K3 \times E$. Scope (B), the Gritsenko--Cl\'ery 8-form atlas $N \in \{1, \ldots, 8\}$, hosts paramodular reflective products of weights $(5, 2, 1, 1, 1/2, 1, 1/4, 0)$ with principal-part constants $(10, 4, 2, 2, 1, 2, 1/2, 0)$; these live on the $\mathrm{Sp}_4$ / $\mathrm{Mp}_4$ / $\widetilde{\mathrm{Mp}}_4$ tower and include metaplectic and spin-cover forms absent from Scope (A).

Both scopes are theorems. Both are load-bearing. Neither is the \emph{only} valid reading of \eqref{eq:kappa-bkm-universal}; the two scopes together are how the formula deserves to be quoted, and every citation of $\kappa_{\mathrm{BKM}} = c_N(0, 0)/2$ names which scope it invokes.

\subsection{Global denominator identity and Fourier seed values}

The denominator identity for $\mathfrak g_{\Delta_5}$ reads, in compressed notation,
\[
  \sum_{w \in W} \mathrm{sgn}(w) \, w\bigl(e^{2\pi i \langle \rho, Z\rangle} \prod_{\alpha > 0}^{\mathrm{real}}(1 - e^{2\pi i \langle \alpha, Z\rangle})\bigr)
  \;=\; \Delta_5(Z),
\]
with $W$ the Weyl group generated by the three real simple reflections. The left-hand Weyl-sum is the Jacobi denominator, which for the $\widetilde A_2$ base rewrites as a product over positive roots $\alpha$ of $\widetilde A_2$ in the Mukai lattice. The right-hand Borcherds product is the multiplicative Fourier expansion:
\[
  \Delta_5(Z) \;=\; q^{1/2} r^{1/2} s^{1/2} \prod_{(n, \ell, m) > 0} \bigl(1 - q^n r^\ell s^m\bigr)^{c_1(4nm - \ell^2)},
\]
with $(q, r, s) = \bigl(e^{2\pi i \tau}, e^{2\pi i z}, e^{2\pi i w}\bigr)$, the triple $(n, \ell, m) > 0$ running over a fundamental chamber, and $c_1(D)$ the Fourier coefficient of $\phi_{0, 1}^{\mathrm{K3}}$ at discriminant $D = 4nm - \ell^2$.

For a weight-$0$, index-$1$ weak Jacobi form the discriminant $D = 4n - \ell^2$ reduces modulo $4$ to $0$ or $3$; equivalently, $D \in \{-1, 0, 3, 4, 7, 8, \ldots\}$. Eichler--Zagier (\emph{The Theory of Jacobi Forms}, Theorem~2.2) shows that $c_1(n, \ell)$ depends only on the pair $(D, \ell \bmod 2)$; write $C_1(D)$ for the resulting discriminant-indexed Fourier coefficient. The direct Laurent expansion
\[
  \phi_{0, 1}^{\mathrm{K3}}(\tau, z) \;=\; q^{-1}(y^{-1} + 10 + y) \;+\; q\bigl(10 y^{-2} - 64 y^{-1} + 108 - 64 y + 10 y^2\bigr) \;+\; O(q^2)
\]
yields the first discriminant values
\[
  C_1(-1) = 1, \quad C_1(0) = 10, \quad C_1(3) = -64, \quad C_1(4) = 108, \quad \ldots,
\]
with $C_1(0) = 10$ realised both at $(n, \ell) = (0, 0)$ and at $(1, \pm 2)$ (Eichler--Zagier discriminant coherence). Discriminants $D \equiv 1, 2 \pmod 4$ are absent from the index-$1$ lattice. The convention $c_1(D = -1) = 2$ adopted in some denominator-identity presentations (Gritsenko--Nikulin 1998, \emph{Duke}~94) counts both $\ell = +1$ and $\ell = -1$ at $n = 0$ additively; we keep the single-$\ell$ normalisation $C_1(-1) = 1$ throughout the present section and flag the doubling at point of use. The integer $C_1(0) = 10$ appears both as the weight of the \emph{imaginary Weyl vector} $\rho$ of $\mathfrak g_{\Delta_5}$ and as the singular-theta constant term entering the Borcherds weight formula \eqref{eq:borcherds-weight}; $C_1(D)$ for $D \geq 3$ gives positive-root multiplicities
\[
  \mathrm{mult}(\alpha) \;=\; C_1\bigl(4nm - \ell^2\bigr), \qquad \alpha = (n, \ell, m) \in \Delta^+(\mathfrak g_{\Delta_5}).
\]

A central-$L$-value reading of the dyonic square $\Phi_{10} = \Delta_5^2$ produces a simple pole whose residue is the automorphic input to the BKM weight. The programme-scale parameter $s_{\mathrm{prog}} = 5/2$ rescales to Andrianov's convention $s_{\mathrm{Andr}}$ by a factor of $4$ (not $2$): one factor from $\Delta_5 \mapsto \Delta_{10} = \Delta_5^2$ Siegel-weight doubling, one from the unshifted-Andrianov $s_{\mathrm{Andr}} = k$ vs programme $s = k/2$ convention. Hence $s_{\mathrm{prog}} = 5/2 \mapsto s_{\mathrm{Andr}} = 10$. The Saito--Kurokawa lift factors the Andrianov spinor $L$-function as
\[
  L_{\mathrm{spin}}(s, \Phi_{10}) \;=\; \zeta(s - 8) \, \zeta(s - 9) \, L(s, g),
\]
where $g = \Delta \cdot E_6 \in S_{18}(\mathrm{SL}_2(\Z))$ is the weight-$18$ elliptic cusp form (Andrianov 1974). At $s = 10$, the factor $\zeta(s - 9) = \zeta(1) = \infty$ forces a simple pole; the residue is
\begin{equation}
  \mathrm{Res}_{s = 10} L_{\mathrm{spin}}(s, \Phi_{10}) \;=\; -15120 \cdot a_{10}(g) \cdot \Omega^-(g),
  \label{eq:saito-kurokawa-residue}
\end{equation}
with $a_{10}(g) \in \Q$ the $10$-th Manin-polynomial Hecke coefficient of $g$ and $\Omega^-(g)$ the Manin minus-period. Ichino--Ikeda's theta-lift period formula identifies the Deligne period $c^+(\Phi_{10}, 5/2)$ of the SK Siegel form with the Manin period of the elliptic source $g$, not with a Siegel Petersson norm; this is a structural consequence of the SK lift being the theta correspondence from $\mathrm{Mp}_2$ to $\mathrm{Sp}_4$. The constant $15120 = 2^4 \cdot 3^3 \cdot 5 \cdot 7$ and the Hecke coefficient $a_{10}(g)$ are read from the Hecke expansion of $g$; the detailed derivation is in the K3~$\times$~E BKM chapter (\texttt{chapters/examples/k3e\_bkm\_chapter.tex}, remark \texttt{rem:wave15-fourier-seed-central-value}).

\subsection{The closing synthesis}

The reader has now seen four modes of CoHA gluing (Parts II--V) produce local-to-global assemblies on geometries with enough local structure (toric, orbifold, tilting-presented, or factorisation-locally-smooth). The present Part VII is the register which takes over when \emph{none} of the above suffice on their own: the lattice-polarised / automorphic stratum. Here the cocycle does not live on the variety at all; it lives on the period domain, and it is a \emph{single Siegel modular form} whose Fourier coefficients determine the BKM target arithmetic and the Euler/superdimension shadow. CoHA structure constants enter only through a finite source fixture whose projected table is recognised by that automorphic target.

The one identity $\kappa_{\mathrm{BKM}}(\Phi_N) = c_N(0) / 2$ summarises what the automorphic cocycle contributes to the Vol III ledger. Across the Gritsenko--Cl\'ery family, eight distinct automorphic cocycles produce eight distinct $\kappa_{\mathrm{BKM}}$ values; in the master K3~$\times$~E instance, the cocycle is $\Delta_5$, the value is $5$, and the global algebra $\mathfrak g_{\Delta_5}$ is the BKM superalgebra whose positive half receives the collected local Hall--Drinfeld pieces from the $24$ $I_1 \times E$-stalks. The further assembly of these with the Mathieu orbifold stratum (iii) and the reduced Aut-equivariance stratum (ii) is the content of Part VIII; the present part establishes only the automorphic contribution, in the precise form that survives without the other strata's help.

The period domain $\mathcal D_{4, 20}$, the Mukai pairing \eqref{eq:mukai-pairing}, the Jacobi form $\phi_{0, 1}^{\mathrm{K3}}$, the Borcherds product \eqref{eq:borcherds-product}, the Siegel form $\Delta_5$, and the BKM algebra $\mathfrak g_{\Delta_5}$ -- all these fit into a single commuting diagram in which target-coefficient extraction is forced by Borcherds's theta-lift theorem combined with Gritsenko's additive-multiplicative lift matching. The Hall multiplication table is a separate finite-source input; the diagram recognises it after projection to the automorphic shadow. This diagram is the gluing cocycle of Stratum (iv). We have now derived it from first principles.


\input{gluing_ch_part_8_k3xe_master}

\input{gluing_ch_part_9_obstructions}

\input{gluing_ch_part_10_unifying}

% ============================================================================
% Chapter-closing synthesis (three sentences, CG convention)
% ============================================================================

\part*{Closing synthesis}

The ten Parts have reduced the assembly question to a single invariant of $X$:
the stratum-set $\Str(X)$, read as a subset of $\{\mathrm{i}, \mathrm{ii},
\mathrm{iii}, \mathrm{iv}\}$, with each occupied stratum supplying a cocycle
theory whose class determines the global Hall bialgebra up to an
$(2,1)$-coherent equivalence. The reduction is not uniformly theorem-grade:
Parts~II, III, IV establish full theorems on their respective strata;
Parts~V, VI, VII establish the $\Phi_3$, wall-crossing, and Borcherds
cocycles up to explicit primary-literature-cited conjectures (two-stage
factorisation for non-simply-laced super, the $c_{11} = -1$ Hopf-pairing
correction at all-mode propagation, the universal $\kBKM(\Phi_N) = c_N(0)/2$
identity across both Gritsenko--Nikulin CHL and Gritsenko--Cl\'ery 8-form
scopes); Part~VIII verifies the master example, Part~IX sharpens the failure
boundary, and Part~X names the sheaf. The next chapter applies the gluing
framework to the seven faces of $r_{\CY}$, extracting from each stratum-cell
the precise chiral-algebra shadow of the classical $r$-matrix.

\end{document}
